% Generated by matins/tools/build_latex.py from data/corpus.json — do not edit.
\documentclass[10pt,twoside,openright]{book}
\usepackage{matins}
\begin{document}
\BodySize
\frontmatter
\thispagestyle{empty}
\vspace*{0.25\textheight}
{\centering
{\Huge\rubr{Lectiones ad Matutinum}}\par\vspace{0.4em}
{\large\itshape ex Breviario Romano anni MCMLIV}\par\vspace{2em}
{\LARGE The Lessons of Matins}\par\vspace{0.4em}
{\large\itshape from the Roman Breviary of 1954, in Latin and English}\par\vspace{3em}
{\small Sample: Advent}\par}
\clearpage
\thispagestyle{empty}
{\footnotesize\setlength{\parindent}{0pt}\setlength{\parskip}{0.35em}\raggedright
\par\addvspace{0.6em}\textsc{Divinum Officium}\par
The texts, and the program that arranges them by the rubrics, are the work of the Divinum Officium Project (divinumofficium.com), created by the late Laszlo Kiss and continued since by the Project's team of volunteers. Its code and data are freely available at github.com/DivinumOfficium/divinum-officium under the MIT License.\par
\par\addvspace{0.6em}\textsc{Latin}\par
The Latin text derives from the Ratisbon 1888 edition of the Breviarium Romanum, scanned by the University of St Michael's College, Toronto, amended from the 1906 edition, and corrected by Laszlo Kiss and the Project team against later printings: Pustet (Regensburg, 1943), Benziger Brothers (New York, 1945), Burns Oates \& Washbourne (London, 1946), La Presse Catholique Panaméricaine (Montréal, 1943), and Desclée and Mame (1962).\par
The accented text of the Sixto-Clementine Vulgate was supplied by Frère Romain-Marie of the Abbaye Saint-Joseph de Clairval, Flavigny-sur-Ozerain.\par
\par\addvspace{0.6em}\textsc{English}\par
Holy Scripture is given in the Douay-Rheims translation.\par
The other lessons and the responsories are from the translation of the Roman Breviary by John, Marquess of Bute (edition of 1908), made available by the University of St Michael's College, Toronto.\par
\par\addvspace{0.6em}\textsc{This edition}\par
The lessons were gathered by running the Divinum Officium program, unchanged, for every day from 1950 to 2100 under the rubrics of 1954, and tracing each lesson it read to the office it belongs to. Errors of arrangement in this edition are not the Project's.\par
\par\addvspace{0.6em}\textsc{Typeface}\par
Set in EB Garamond, designed by Georg Duffner and Octavio Pardo, under the SIL Open Font License 1.1.\par
}
\mainmatter
\PartTitle{Proprium de Tempore}{Proper of Time}

\SeasonTitle{Tempus Adventus}{Advent}

\Office{Dominica I Adventus}{First Sunday of Advent}{Semiduplex I classis}
\begin{paracol}{2}
\LA
\LessonHead{Lectio I}
\switchcolumn
\EN
\LessonHead{Lesson I}
\switchcolumn*

\LA
\LessonTitle{Incipit liber Isaíæ Prophétæ}
\switchcolumn
\EN
\LessonTitle{Lesson from the book of Isaias}
\switchcolumn*

\LA
\Cite{Isa 1:1-3}
\switchcolumn
\EN
\Cite{Isa 1:1-3}
\switchcolumn*

\LA
\noindent Vísio Isaíæ fílii Amos, quam vidit super Judam et Jerúsalem, in diébus Ozíæ, Jóatham, Achaz, et Ezechíæ, regum Juda. Audíte, cæli, et áuribus pércipe, terra, quóniam Dóminus locútus est: Fílios enutrívi, et exaltávi: ipsi autem sprevérunt me. Cognóvit bos possessórem suum, et ásinus præsépe dómini sui: Israël autem me non cognóvit, et pópulus meus non intelléxit.\par
\switchcolumn
\EN
\noindent The vision of Isaias the son of Amos, which he saw concerning Juda and Jerusalem in the days of Ozias, Joathan, Achaz, and Ezechias, kings of Juda Hear, O ye heavens, and give ear, O earth, for the Lord hath spoken. I have brought up children, and exalted them: but they have despised me. The ox knoweth his owner, and the ass his master's crib: but Israel hath not known me, and my people hath not understood.\par
\switchcolumn*

\LA
\begin{Resp}\Rsign Aspíciens a longe, ecce vídeo Dei poténtiam veniéntem, et nébulam totam terram tegéntem. \Star{} Ite óbviam ei, et dícite: * Núntia nobis, si tu es ipse, * Qui regnatúrus es in pópulo Israël.\end{Resp}
\begin{Resp}\Vsign Quique terrígenæ, et fílii hóminum, simul in unum dives et pauper. Ite óbviam ei, et dícite.\end{Resp}
\begin{Resp}\Vsign Qui regis Israël, inténde, qui dedúcis velut ovem Joseph. Núntia nobis, si tu es ipse.\end{Resp}
\begin{Resp}\Vsign Tóllite portas, príncipes, vestras, et elevámini portæ æternáles, et introíbit Rex glóriæ. Qui regnatúrus es in pópulo Israël.\end{Resp}
\begin{Resp}\Vsign Glória Patri, et Fílio, \Star{} et Spirítui Sancto.\end{Resp}
\begin{Resp}\Rsign Aspíciens a longe, ecce vídeo Dei poténtiam veniéntem, et nébulam totam terram tegéntem. * Ite óbviam ei, et dícite: * Núntia nobis, si tu es ipse, * Qui regnatúrus es in pópulo Israël.\end{Resp}
\switchcolumn
\EN
\begin{Resp}\Rsign I look from afar, and behold I see the Power of God, coming like as a cloud to cover the land with the hosts of his People: \Star{} Go ye out to meet him and say: * Tell us if thou art he, * That shalt reign over God's people Israel.\end{Resp}
\begin{Resp}\Vsign All ye that dwell in the world, all ye children of men, high and low, rich and poor, one with another. Go ye out to meet him and say.\end{Resp}
\begin{Resp}\Vsign Hear, O thou Shepherd of Israel, thou that leadest Joseph like a sheep. Tell us if thou art he.\end{Resp}
\begin{Resp}\Vsign Lift up your heads, O ye gates, and be ye lift up, ye everlasting doors, and the King of Glory shall come in. That shalt reign over God's people Israel.\end{Resp}
\begin{Resp}\Vsign Glory be to the Father, and to the Son, \Star{} and to the Holy Ghost.\end{Resp}
\begin{Resp}\Rsign I look from afar, and behold I see the Power of God, coming like as a cloud to cover the land with the hosts of his People: * Go ye out to meet him and say: * Tell us if thou art he, * That shalt reign over God's people Israel.\end{Resp}
\switchcolumn*

\LA
\LessonHead{Lectio II}
\switchcolumn
\EN
\LessonHead{Lesson II}
\switchcolumn*

\LA
\Cite{Isa 1:4-6}
\switchcolumn
\EN
\Cite{Isa 1:4-6}
\switchcolumn*

\LA
\noindent Væ genti peccatríci, pópulo gravi iniquitáte, sémini nequam, fíliis scelerátis: dereliquérunt Dóminum, blasphemavérunt Sanctum Israël, abalienáti sunt retrórsum. Super quo percútiam vos ultra, addéntes prævaricatiónem? omne caput lánguidum, et omne cor mærens. A planta pedis usque ad vérticem non est in eo sánitas: vulnus, et livor, et plaga tumens non est circumligáta, nec curáta medicámine, neque fota óleo.\par
\switchcolumn
\EN
\noindent Woe to the sinful nation, a people laden with iniquity, a wicked seed, ungracious children: they have forsaken the Lord, they have blasphemed the Holy One of Israel, they are gone away backwards. For what shall I strike you any more, you that increase transgression? the whole head is sick, and the whole heart is sad. From the sole of the foot unto the top of the head, there is no soundness therein: wounds and bruises and swelling sores: they are not bound up, nor dressed, nor fomented with oil.\par
\switchcolumn*

\LA
\begin{Resp}\Rsign Aspiciébam in visu noctis, et ecce in núbibus cæli Fílius hóminis veniébat: et datum est ei regnum, et honor: \Star{} Et omnis pópulus, tribus, et linguæ sérvient ei.\end{Resp}
\begin{Resp}\Vsign Potéstas ejus, potéstas ætérna, quæ non auferétur: et regnum ejus, quod non corrumpétur.\end{Resp}
\begin{Resp}\Rsign Et omnis pópulus, tribus, et linguæ sérvient ei.\end{Resp}
\switchcolumn
\EN
\begin{Resp}\Rsign I saw in the night visions, and, behold, the Son of man came with the clouds of heaven, and there was given Him a Kingdom, and glory; \Star{} And all people, nations, and languages shall serve Him.\end{Resp}
\begin{Resp}\Vsign His dominion is an everlasting dominion which shall not pass away, and His Kingdom that which shall not be destroyed.\end{Resp}
\begin{Resp}\Rsign And all people, nations, and languages shall serve Him.\end{Resp}
\switchcolumn*

\LA
\LessonHead{Lectio III}
\switchcolumn
\EN
\LessonHead{Lesson III}
\switchcolumn*

\LA
\Cite{Isa 1:7-9}
\switchcolumn
\EN
\Cite{Isa 1:7-9}
\switchcolumn*

\LA
\noindent Terra vestra desérta, civitátes vestræ succénsæ igni: regiónem vestram coram vobis aliéni dévorant, et desolábitur sicut in vastitáte hostíli. Et derelinquétur fília Sion ut umbráculum in vínea, et sicut tugúrium in cucumerário, et sicut cívitas quæ vastátur. Nisi Dóminus exercítuum reliquísset nobis semen, quasi Sódoma fuissémus, et quasi Gomórrha símiles essémus.\par
\switchcolumn
\EN
\noindent Your land is desolate, your cities are burnt with fire: your country strangers devour before your face, and it shall be desolate as when wasted by enemies. And the daughter of Sion shall be left as a covert in a vineyard, and as a lodge in a garden of cucumbers, and as a city that is laid waste. Except the Lord of hosts had left us seed, we had been as Sodom, and we should have been like to Gomorrha.\par
\switchcolumn*

\LA
\begin{Resp}\Rsign Missus est Gábriel Angelus ad Maríam Vírginem desponsátam Joseph, núntians ei verbum; et expavéscit Virgo de lúmine: ne tímeas, María, invenísti grátiam apud Dóminum: \Star{} Ecce concípies, et páries, et vocábitur Altíssimi Fílius.\end{Resp}
\begin{Resp}\Vsign Dabit ei Dóminus Deus sedem David, patris ejus, et regnábit in domo Jacob in ætérnum.\end{Resp}
\begin{Resp}\Rsign Ecce concípies, et páries, et vocábitur Altíssimi Fílius.\end{Resp}
\begin{Resp}\Vsign Glória Patri, et Fílio, \Star{} et Spirítui Sancto.\end{Resp}
\begin{Resp}\Rsign Ecce concípies, et páries, et vocábitur Altíssimi Fílius.\end{Resp}
\switchcolumn
\EN
\begin{Resp}\Rsign The Angel Gabriel was sent to Mary, a Virgin espoused to Joseph, to bring unto her the word of the Lord, and when the Virgin saw the light, she was afraid. Fear not, Mary, for thou hast found grace from the Lord. \Star{} Behold, thou shalt conceive and bring forth a son, and He shall be called the Son of the Highest.\end{Resp}
\begin{Resp}\Vsign The Lord God shall give unto Him the throne of His father David, and He shall reign over the house of Jacob for ever.\end{Resp}
\begin{Resp}\Rsign Behold, thou shalt conceive, and bring forth a son, and He shall be called the Son of the Highest.\end{Resp}
\begin{Resp}\Vsign Glory be to the Father, and to the Son, \Star{} and to the Holy Ghost.\end{Resp}
\begin{Resp}\Rsign Behold, thou shalt conceive and bring forth a son, and He shall be called the Son of the Highest.\end{Resp}
\switchcolumn*

\LA
\LessonHead{Lectio IV}
\switchcolumn
\EN
\LessonHead{Lesson IV}
\switchcolumn*

\LA
\LessonTitle{Sermo sancti Leónis Papæ}
\switchcolumn
\EN
\LessonTitle{From the Sermons of Pope St. Leo (the Great,)}
\switchcolumn*

\LA
\Source{Sermo 8 de jejunio decimi mensis et eleemosynis}
\switchcolumn
\EN
\Source{8th on the December Fast, and almsgiving.}
\switchcolumn*

\LA
\noindent Cum de advéntu regni Dei, et de mundi fine ac témporum, discípulos suos Salvátor instrúeret, totámque Ecclésiam suam in Apóstolis erudíret: Cavéte, inquit, ne forte gravéntur corda vestra in crápula, et ebrietáte, et cogitatiónibus sæculáribus. Quod útique præcéptum, dilectíssimi, ad nos speciálius pertinére cognóscimus, quibus denuntiátus dies, etiámsi est occúltus, non dubitátur esse vicínus.\par
\switchcolumn
\EN
\noindent Our Saviour Himself instructed His disciples concerning the times and seasons of the coming of the Kingdom of God and the end of the world, and He hath given the same teaching to the Church by the mouth of His Apostles. In connection with this subject then, Our Lord biddeth us beware lest we let our hearts grow heavy through excess of meat and drink, and worldly thoughts. Dearly beloved brethren, we know how that this warning applieth particularly to us. We know that that day is coming, and though for a season we know not the very hour, yet this we know, that it is near.\par
\switchcolumn*

\LA
\begin{Resp}\Rsign Ave, María, grátia plena, Dóminus tecum: \Star{} Spíritus Sanctus supervéniet in te, et virtus Altíssimi obumbrábit tibi: quod enim ex te nascétur Sanctum, vocábitur Fílius Dei.\end{Resp}
\begin{Resp}\Vsign Quómodo fiet istud, quóniam virum non cognósco? Et respóndens Angelus, dixit ei.\end{Resp}
\begin{Resp}\Rsign Spíritus Sanctus supervéniet in te, et virtus Altíssimi obumbrábit tibi: quod enim ex te nascétur Sanctum, vocábitur Fílius Dei.\end{Resp}
\switchcolumn
\EN
\begin{Resp}\Rsign Hail, Mary, full of grace; the Lord is with thee; \Star{} The Holy Ghost shall come upon thee, and the power of the Highest shall overshadow thee; therefore, also that Holy Thing Which shall be born of thee shall be called the Son of God.\end{Resp}
\begin{Resp}\Vsign How shall this be, seeing I know not a man? And the Angel answered and said unto her,\end{Resp}
\begin{Resp}\Rsign The Holy Ghost shall come upon thee, and the power of the Highest shall overshadow thee; therefore, also that Holy Thing Which shall be born of thee shall be called the Son of God.\end{Resp}
\switchcolumn*

\LA
\LessonHead{Lectio V}
\switchcolumn
\EN
\LessonHead{Lesson V}
\switchcolumn*

\LA
\noindent Ad cujus advéntum omnem hóminem cónvenit præparári: ne quem aut ventri déditum, aut curis sæculáribus invéniat implicátum. Quotidiáno enim, dilectíssimi, experiménto probátur, potus satietáte áciem mentis obtúndi, et cibórum nimietáte vigórem cordis hebetári; ita ut delectátio edéndi étiam córporum contrária sit salúti, nisi rátio temperántiæ obsístat illécebræ, et quod futúrum est óneri, súbtrahat voluptáti.\par
\switchcolumn
\EN
\noindent Let every man then make himself ready against the coming of the Lord, so that He may not find him making his belly his god, or the world his chief care. Dearly beloved brethren, it is a matter of every day experience that fulness of drink dulleth the keenness of the mind, and that excess of eating unnerveth the strength of the will. The very stomach protesteth that gluttony doth harm to the bodily health, unless temperance get the better of desire, and the thought of the indigestion afterward check the indulgence of the moment.\par
\switchcolumn*

\LA
\begin{Resp}\Rsign Salvatórem exspectámus, Dóminum Jesum Christum, \Star{} Qui reformábit corpus humilitátis nostræ configurátum córpori claritátis suæ.\end{Resp}
\begin{Resp}\Vsign Sóbrie, et juste, et pie vivámus in hoc sǽculo, exspectántes beátam spem, et advéntum glóriæ magni Dei.\end{Resp}
\begin{Resp}\Rsign Qui reformábit corpus humilitátis nostræ configurátum córpori claritátis suæ.\end{Resp}
\switchcolumn
\EN
\begin{Resp}\Rsign We look for the Saviour, the Lord Jesus Christ; \Star{} Who shall change our vile body, that it may be fashioned like unto His glorious Body.\end{Resp}
\begin{Resp}\Vsign We should live soberly, and righteously, and godly in this present world, looking for that blessed hope, and the glorious appearing of the great God.\end{Resp}
\begin{Resp}\Rsign Who shall change our vile body, that it may be fashioned like unto His glorious Body.\end{Resp}
\switchcolumn*

\LA
\LessonHead{Lectio VI}
\switchcolumn
\EN
\LessonHead{Lesson VI}
\switchcolumn*

\LA
\noindent Quamvis enim sine ánima nihil caro desíderet, et inde accípiat sensus, unde sumit et motus: ejúsdem tamen est ánimæ, quædam sibi súbditæ negáre substántiæ, et interióri judício ab inconveniéntibus exterióra frenáre: ut a corpóreis cupiditátibus sǽpius líbera, in aula mentis possit divínæ vacáre sapiéntiæ: ubi omni strépitu terrenárum silénte curárum, in meditatiónibus sanctis, et in delíciis lætétur ætérnis.\par
\switchcolumn
\EN
\noindent The body without the soul hath no desires; its sensibility cometh from the same source as its movements. And it is the duty of a man with a reasonable soul to deny something to his lower nature and to keep back the outer man from things unseemly. Then will his soul, free from fleshly cravings, sit often at leisure in the palace of the mind, dwelling on the wisdom of God. There, when the roar and rattle of earthly cares are stilled, will she feed on holy thoughts and entertain herself with the expectation of the everlasting joy.\par
\switchcolumn*

\LA
\begin{Resp}\Rsign Obsecro, Dómine, mitte quem missúrus es: vide afflictiónem pópuli tui: \Star{} Sicut locútus es, veni, * Et líbera nos.\end{Resp}
\begin{Resp}\Vsign Qui regis Israël, inténde, qui dedúcis velut ovem Joseph, qui sedes super Chérubim.\end{Resp}
\begin{Resp}\Rsign Sicut locútus es, veni.\end{Resp}
\begin{Resp}\Vsign Glória Patri, et Fílio, \Star{} et Spirítui Sancto.\end{Resp}
\begin{Resp}\Rsign Et líbera nos.\end{Resp}
\switchcolumn
\EN
\begin{Resp}\Rsign O my Lord, send I pray thee, Him Whom Thou wilt send; see the affliction of thy people. As Thou hast promised, come; \Star{} And deliver us.\end{Resp}
\begin{Resp}\Vsign Give ear, O Shepherd of Israel, Thou That leadest Joseph like a flock, Thou That sittest upon the Cherubim!\end{Resp}
\begin{Resp}\Rsign As Thou hast promised, come.\end{Resp}
\begin{Resp}\Vsign Glory be to the Father, and to the Son, \Star{} and to the Holy Ghost.\end{Resp}
\begin{Resp}\Rsign And deliver us.\end{Resp}
\switchcolumn*

\LA
\LessonHead{Lectio VII}
\switchcolumn
\EN
\LessonHead{Lesson VII}
\switchcolumn*

\LA
\LessonTitle{Léctio sancti Evangélii secúndum Lucam}
\switchcolumn
\EN
\LessonTitle{From the Holy Gospel according to Luke}
\switchcolumn*

\LA
\Cite{Luc 21:25-33}
\switchcolumn
\EN
\Cite{Luke 21:25-33}
\switchcolumn*

\LA
\noindent In illo témpore: Dixit Jesus discípulis suis: Erunt signa in sole, et luna, et stellis, et in terris pressúra géntium. Et réliqua.\par
\switchcolumn
\EN
\noindent In that time, Jesus said to his disciples: And there shall be signs in the sun, and in the moon, and in the stars; and upon the earth distress of nations. And so on.\par
\switchcolumn*

\LA
\LessonTitle{Homilía sancti Gregórii Papæ}
\switchcolumn
\EN
\LessonTitle{Homily by Pope St. Gregory (the Great,)}
\switchcolumn*

\LA
\Source{Homilia 1 in Evangelia}
\switchcolumn
\EN
\Source{1st on the Gospels}
\switchcolumn*

\LA
\noindent Dóminus ac Redémptor noster parátos nos inveníre desíderans, senescéntem mundum quæ mala sequántur denúntiat, ut nos ab ejus amóre compéscat. Appropinquántem ejus términum quantæ percussiónes prævéniant, innotéscit: ut, si Deum metúere in tranquillitáte nólumus, saltem vicínum ejus judícium vel percussiónibus attríti timeámus.\par
\switchcolumn
\EN
\noindent Our Lord and Saviour wisheth to find us ready at His second coming. Therefore He telleth us what will be the evils of the world as it groweth old, that He may wean our hearts from worldly affections. Here we read what great convulsions will go before the end, that, if we will not fear God in our prosperity, we may at least be scourged into fearing His judgment when it is at hand.\par
\switchcolumn*

\LA
\begin{Resp}\Rsign Ecce virgo concípiet, et páriet fílium, dicit Dóminus: \Star{} Et vocábitur nomen ejus Admirábilis, Deus, Fortis.\end{Resp}
\begin{Resp}\Vsign Super sólium David, et super regnum ejus sedébit in ætérnum.\end{Resp}
\begin{Resp}\Rsign Et vocábitur nomen ejus Admirábilis, Deus, Fortis.\end{Resp}
\switchcolumn
\EN
\begin{Resp}\Rsign Behold, the Virgin shall conceive, and bear a son, saith the Lord; \Star{} And His name shall be called Wonderful, the Mighty God.\end{Resp}
\begin{Resp}\Vsign He shall sit upon the throne of David, and upon his kingdom for ever.\end{Resp}
\begin{Resp}\Rsign And His name shall be called Wonderful, the Mighty God.\end{Resp}
\switchcolumn*

\LA
\LessonHead{Lectio VIII}
\switchcolumn
\EN
\LessonHead{Lesson VIII}
\switchcolumn*

\LA
\noindent Huic étenim lectióni sancti Evangélii, quam modo vestra fratérnitas audívit, paulo supérius Dóminus præmísit, dicens: Exsúrget gens contra gentem, et regnum advérsus regnum: et erunt terræmótus magni per loca, et pestiléntiæ, et fames. Et quibúsdam interpósitis, hoc, quod modo audístis, adjúnxit: Erunt signa in sole, et luna, et stellis, et in terris pressúra géntium præ confusióne sónitus maris, et flúctuum. Ex quibus profécto ómnibus ália jam facta cérnimus, ália in próximo ventúra formidámus.\par
\switchcolumn
\EN
\noindent Immediately before the passage which hath just been read from the Gospel, are found the following words of our Lord, Nation shall rise against nation, and kingdom against kingdom, and great earthquakes shall be in diverse places, and pestilences and famines. Then, after a few more verses, cometh today's Gospel. There shall be signs in the sun, and in the moon, and in the stars; and upon the earth distress of nations with perplexity, the sea and the waves roaring. Now some of these things are come to pass already, and we fear the others are not far off.\par
\switchcolumn*

\LA
\begin{Resp}\Rsign Audíte verbum Dómini, gentes, et annuntiáte illud in fínibus terræ: \Star{} Et ínsulis, quæ procul sunt, dícite: Salvátor noster advéniet.\end{Resp}
\begin{Resp}\Vsign Annuntiáte, et audítum fácite: loquímini, et clamáte.\end{Resp}
\begin{Resp}\Rsign Et ínsulis, quæ procul sunt, dícite: Salvátor noster advéniet.\end{Resp}
\switchcolumn
\EN
\begin{Resp}\Rsign Hear the word of the Lord, O ye nations, and declare it in the ends of the earth; \Star{} And in the isles afar off, and say: Our Saviour shall come.\end{Resp}
\begin{Resp}\Vsign Declare it and make it known, lift up your voice and cry aloud.\end{Resp}
\begin{Resp}\Rsign And in the isles afar off, and say: Our Saviour shall come.\end{Resp}
\switchcolumn*

\LA
\LessonHead{Lectio IX}
\switchcolumn
\EN
\LessonHead{Lesson IX}
\switchcolumn*

\LA
\noindent Nam gentem contra gentem exsúrgere, earúmque pressúram terris insístere, plus jam in nostris tempóribus cérnimus, quam in codícibus légimus. Quod terræmótus urbes innúmeras óbruat, ex áliis mundi pártibus scitis quam frequénter audívimus. Pestiléntias sine cessatióne pátimur. Signa vero in sole, et luna, et stellis, adhuc apérte mínime vídimus: sed quia et hæc non longe sint, ex ipsa jam áëris immutatióne collígimus.\par
\switchcolumn
\EN
\noindent In these our days we see nation rise against nation, and their distress over all the earth, more than we read in books hath ever come to pass of old time. Ye know also how often we hear of earthquakes overwhelming countless cities in other parts of the world. As for pestilences, we suffer from them ourselves, with hardly any intermission. As yet we do not see signs in the sun, and in the moon, and in the stars; but the changes of seasons and climates warn us that we may look for these also before long.\par
\switchcolumn*

\LA
\begin{Resp}\Rsign Ecce dies véniunt, dicit Dóminus, et suscitábo David germen justum: et regnábit rex, et sápiens erit, et fáciet judícium et justítiam in terra: \Star{} Et hoc est nomen quod vocábunt eum: * Dóminus justus noster.\end{Resp}
\begin{Resp}\Vsign In diébus illis salvábitur Juda, et Israël habitábit confidénter.\end{Resp}
\begin{Resp}\Rsign Et hoc est nomen quod vocábunt eum:\end{Resp}
\begin{Resp}\Vsign Glória Patri, et Fílio, \Star{} et Spirítui Sancto.\end{Resp}
\begin{Resp}\Rsign Dóminus justus noster.\end{Resp}
\switchcolumn
\EN
\begin{Resp}\Rsign Behold, the days come, saith the Lord, that I will raise unto David a righteous Branch; and a King shall reign in wisdom and shall execute judgment and justice in the earth: \Star{} And this is His name whereby He shall be called: * The Lord our Righteous one.\end{Resp}
\begin{Resp}\Vsign In His days Judah shall be saved, and Israel shall dwell safely.\end{Resp}
\begin{Resp}\Rsign And this is His name whereby He shall be called.\end{Resp}
\begin{Resp}\Vsign Glory be to the Father, and to the Son, \Star{} and to the Holy Ghost.\end{Resp}
\begin{Resp}\Rsign The Lord our Righteous one.\end{Resp}
\switchcolumn*

\end{paracol}

\Office{Feria II infra Hebdomadam I Adventus}{Monday of the First Week of Advent}{Feria major}
\begin{paracol}{2}
\LA
\LessonHead{Lectio I}
\switchcolumn
\EN
\LessonHead{Lesson I}
\switchcolumn*

\LA
\LessonTitle{De Isaía Prophéta}
\switchcolumn
\EN
\LessonTitle{Lesson from the book of Isaias}
\switchcolumn*

\LA
\Cite{Isa 1:16-18}
\switchcolumn
\EN
\Cite{Isa 1:16-18}
\switchcolumn*

\LA
\noindent Lavámini, mundi estóte, auférte malum cogitatiónum vestrárum ab óculis meis: quiéscite ágere pervérse, díscite benefácere: quǽrite judícium, subveníte opprésso, judicáte pupíllo, deféndite víduam. Et veníte, et argúite me, dicit Dóminus: si fúerint peccáta vestra ut cóccinum, quasi nix dealbabúntur: et si fúerint rubra quasi vermículus, velut lana alba erunt.\par
\switchcolumn
\EN
\noindent Wash yourselves, be clean, take away the evil of your devices from my eyes: cease to do perversely, Learn to do well: seek judgment, relieve the oppressed, judge for the fatherless, defend the widow. And then come, and accuse me, saith the Lord: if your sins be as scarlet, they shall be made as white as snow: and if they be red as crimson, they shall be white as wool.\par
\switchcolumn*

\LA
\begin{Resp}\Rsign Súscipe verbum, Virgo María, quod tibi a Dómino per Angelum transmíssum est: concípies et páries Deum páriter et hóminem, \Star{} Ut benedícta dicáris inter omnes mulíeres.\end{Resp}
\begin{Resp}\Vsign Páries quidem fílium, et virginitátis non patiéris detriméntum: efficiéris grávida, et eris mater semper intácta.\end{Resp}
\begin{Resp}\Rsign Ut benedícta dicáris inter omnes mulíeres.\end{Resp}
\switchcolumn
\EN
\begin{Resp}\Rsign Receive, O Virgin Mary, receive the word of the Lord, which is sent thee by His Angel; thou shalt conceive, and shalt bring forth God and Man together. \Star{} And thou shalt be called blessed among all women.\end{Resp}
\begin{Resp}\Vsign Thou shalt bring forth a son, and remain a maiden undefiled thou shalt conceive and be a Mother, still Virgin unspotted.\end{Resp}
\begin{Resp}\Rsign And thou shalt be called blessed among all women.\end{Resp}
\switchcolumn*

\LA
\LessonHead{Lectio II}
\switchcolumn
\EN
\LessonHead{Lesson II}
\switchcolumn*

\LA
\Cite{Isa 1:19-23}
\switchcolumn
\EN
\Cite{Isa 1:19-23}
\switchcolumn*

\LA
\noindent Si voluéritis, et audiéritis me, bona terræ comedétis. Quod si noluéritis, et me ad iracúndiam provocavéritis, gládius devorábit vos, quia os Dómini locútum est. Quómodo facta est méretrix cívitas fidélis, plena judícii? Justítia habitávit in ea, nunc autem homicídæ. Argéntum tuum versum est in scóriam: vinum tuum mixtum est aqua. Príncipes tui infidéles, sócii furum: omnes díligunt múnera, sequúntur retributiónes. Pupíllo non júdicant: et causa víduæ non ingréditur ad illos.\par
\switchcolumn
\EN
\noindent if you be willing, and will hearken to me, you shall eat the good things of the land. But if you will not, and will provoke me to wrath: the sword shall devour you because the mouth of the Lord hath spoken it. How is the faithful city, that was full of judgment, become a harlot? justice dwelt in it, but now murderers. thy silver is turned into dross: thy wine is mingled with water. thy princes are faithless, companions of thieves: they all love bribes, they run after rewards. They judge not for the fatherless: and the widow's cometh not in to them.\par
\switchcolumn*

\LA
\begin{Resp}\Rsign Læténtur cæli, et exsúltet terra, jubiláte, montes, laudem: quia Dóminus noster véniet, \Star{} Et páuperum suórum miserébitur.\end{Resp}
\begin{Resp}\Vsign Oriétur in diébus ejus justítia, et abundántia pacis.\end{Resp}
\begin{Resp}\Rsign Et páuperum suórum miserébitur.\end{Resp}
\switchcolumn
\EN
\begin{Resp}\Rsign Sing, O heavens; and be joyful, O earth; and break forth into singing, O mountains, for our Lord will come; \Star{} And will have mercy on His afflicted.\end{Resp}
\begin{Resp}\Vsign In His days shall righteousness flourish and abundance of peace.\end{Resp}
\begin{Resp}\Rsign And will have mercy upon His afflicted.\end{Resp}
\switchcolumn*

\LA
\LessonHead{Lectio III}
\switchcolumn
\EN
\LessonHead{Lesson III}
\switchcolumn*

\LA
\Cite{Isa 1:24-28}
\switchcolumn
\EN
\Cite{Isa 1:24-28}
\switchcolumn*

\LA
\noindent Propter hoc ait Dóminus Deus exercítuum fortis Israël: Heu! consolábor super hóstibus meis, et vindicábor de inimícis meis. Et convértam manum meam ad te, et éxcoquam ad purum scóriam tuam, et áuferam omne stannum tuum. Et restítuam júdices tuos ut fuérunt prius, et consiliários tuos sicut antíquitus: post hæc vocáberis cívitas justi, urbs fidélis. Sion in judício redimétur, et redúcent eam in justítia: et cónteret sceléstos, et peccatóres simul: et qui dereliquérunt Dóminum, consuméntur.\par
\switchcolumn
\EN
\noindent Therefore saith the Lord the God of hosts, the mighty one of Israel: Ah! I will comfort myself over my adversaries: and I will be revenged of my enemies. And I will turn my hand to thee, and I will clean purge away thy dross, and I will take away all thy tin. And I will restore thy judges as they were before, and thy counsellors as of old. After this thou shalt be called the city of the just, a faithful city. Sion shall be redeemed in judgment, and they shall bring her back in justice. And he shall destroy the wicked, and the sinners together: and they that have forsaken the Lord, shall be consumed.\par
\switchcolumn*

\LA
\begin{Resp}\Rsign Aliéni non transíbunt per Jerúsalem ámplius: \Star{} Nam in illa die stillábunt montes dulcédinem, et colles fluent lac et mel, dicit Dóminus.\end{Resp}
\begin{Resp}\Vsign Deus a Líbano véniet, et Sanctus de monte umbróso et condénso.\end{Resp}
\begin{Resp}\Rsign Nam in illa die stillábunt montes dulcédinem, et colles fluent lac et mel, dicit Dóminus.\end{Resp}
\begin{Resp}\Vsign Glória Patri, et Fílio, \Star{} et Spirítui Sancto.\end{Resp}
\begin{Resp}\Rsign Nam in illa die stillábunt montes dulcédinem, et colles fluent lac et mel, dicit Dóminus.\end{Resp}
\switchcolumn
\EN
\begin{Resp}\Rsign There shall no strangers pass through Jerusalem any more \Star{} For in that day the mountains shall drop down sweet wine, and the hills shall flow with milk and honey, saith the Lord.\end{Resp}
\begin{Resp}\Vsign God shall come from Lebanon, and the Holy One from the thick and shady mountain.\end{Resp}
\begin{Resp}\Rsign For in that day the mountains shall drop down sweet wine, and the hills shall flow with milk and honey, saith the Lord.\end{Resp}
\begin{Resp}\Vsign Glory be to the Father, and to the Son, \Star{} and to the Holy Ghost.\end{Resp}
\begin{Resp}\Rsign For in that day the mountains shall drop down sweet wine, and the hills shall flow with milk and honey, saith the Lord.\end{Resp}
\switchcolumn*

\end{paracol}

\Office{Feria III infra Hebdomadam I Adventus}{Tuesday of the First Week of Advent}{Feria major}
\begin{paracol}{2}
\LA
\LessonHead{Lectio I}
\switchcolumn
\EN
\LessonHead{Lesson I}
\switchcolumn*

\LA
\LessonTitle{De Isaía Prophéta}
\switchcolumn
\EN
\LessonTitle{Lesson from the book of Isaias}
\switchcolumn*

\LA
\Cite{Isa 2:1-3}
\switchcolumn
\EN
\Cite{Isa 2:1-3}
\switchcolumn*

\LA
\noindent Verbum, quod vidit Isaías, fílius Amos, super Judam et Jerúsalem. Et erit in novíssimis diébus præparátus mons domus Dómini in vértice móntium, et elevábitur super colles; et fluent ad eum omnes gentes. Et ibunt pópuli multi, et dicent: Veníte, et ascendámus ad montem Dómini, et ad domum Dei Jacob, et docébit nos vias suas, et ambulábimus in sémitis ejus: quia de Sion exíbit lex, et verbum Dómini de Jerúsalem.\par
\switchcolumn
\EN
\noindent The word that Isaias the son of Amos saw, concerning Juda and Jerusalem. And in the last days the mountain of the house of the Lord shall be prepared on the top of mountains, and it shall be exalted above the hills, and all nations shall flow unto it. And many people shall go, and say: Come and let us go up to the mountain of the Lord, and to the house of the God of Jacob, and he will teach us his ways, and we will walk in his paths: for the law shall come forth from Sion, and the word of the Lord from Jerusalem.\par
\switchcolumn*

\LA
\begin{Resp}\Rsign Montes Israël, ramos vestros expándite, et floréte, et fructus fácite: \Star{} Prope est ut véniat dies Dómini.\end{Resp}
\begin{Resp}\Vsign Roráte, cæli, désuper, et nubes pluant justum: aperiátur terra, et gérminet Salvatórem.\end{Resp}
\begin{Resp}\Rsign Prope est ut véniat dies Dómini.\end{Resp}
\switchcolumn
\EN
\begin{Resp}\Rsign O ye mountains of Israel, shoot forth your branches and blossom and bring forth fruit. \Star{} The day of the Lord is at hand to come.\end{Resp}
\begin{Resp}\Vsign Drop down, ye heavens, from above, and let the skies pour down the Righteous One; let the earth open, and let her bring forth the Saviour.\end{Resp}
\begin{Resp}\Rsign The day of the Lord is at hand to come.\end{Resp}
\switchcolumn*

\LA
\LessonHead{Lectio II}
\switchcolumn
\EN
\LessonHead{Lesson II}
\switchcolumn*

\LA
\Cite{Isa 2:4-6}
\switchcolumn
\EN
\Cite{Isa 2:4-6}
\switchcolumn*

\LA
\noindent Et judicábit gentes, et árguet pópulos multos: et conflábunt gládios suos in vómeres, et lánceas suas in falces: non levábit gens contra gentem gládium, nec exercebúntur ultra ad prǽlium. Domus Jacob, veníte, et ambulémus in lúmine Dómini. Projecísti enim pópulum tuum, domum Jacob: quia repléti sunt ut olim, et áugures habuérunt ut Philísthiim, et púeris aliénis adhæsérunt.\par
\switchcolumn
\EN
\noindent And he shall judge the Gentiles, and rebuke many people: and they shall turn their swords into ploughshares, and their spears into sickles: nation shall not lift up sword against nation, neither shall they be exercised any more to war. O house of Jacob, come ye, and let us walk in the light of the Lord. For thou hast cast off thy people, the house of Jacob: because they are filled as in times past, and have had soothsayers as the Philistines, and have adhered to strange children.\par
\switchcolumn*

\LA
\begin{Resp}\Rsign Erúmpant montes jucunditátem, et colles justítiam: \Star{} Quia lux mundi Dóminus cum poténtia venit.\end{Resp}
\begin{Resp}\Vsign De Sion exíbit lex, et verbum Dómini de Jerúsalem.\end{Resp}
\begin{Resp}\Rsign Quia lux mundi Dóminus cum poténtia venit.\end{Resp}
\switchcolumn
\EN
\begin{Resp}\Rsign Let the mountains break forth into singing, and the hills bring forth righteousness \Star{} For the Lord, the Light of the world, cometh with power.\end{Resp}
\begin{Resp}\Vsign Out of Zion shall go forth the law, and the word of the Lord from Jerusalem.\end{Resp}
\begin{Resp}\Rsign For the Lord, the Light of the world, cometh with power.\end{Resp}
\switchcolumn*

\LA
\LessonHead{Lectio III}
\switchcolumn
\EN
\LessonHead{Lesson III}
\switchcolumn*

\LA
\Cite{Isa 2:7-9}
\switchcolumn
\EN
\Cite{Isa 2:7-9}
\switchcolumn*

\LA
\noindent Repléta est terra argénto et auro: et non est finis thesaurórum ejus: et repléta est terra ejus equis: et innumerábiles quadrígæ ejus. Et repléta est terra ejus idólis: opus mánuum suárum adoravérunt, quod fecérunt dígiti eórum. Et incurvávit se homo, et humiliátus est vir, ne ergo dimíttas eis.\par
\switchcolumn
\EN
\noindent Their land is filled with silver and gold: and there is no end of their treasures. And their land is filled with horses: and their chariots are innumerable. Their land also is full of idols: they have adored the work of their own hands, which their own fingers have made. And man hath bowed himself down, and man hath been debased: therefore forgive them not.\par
\switchcolumn*

\LA
\begin{Resp}\Rsign Ecce ab Austro vénio, ego Dóminus Deus vester, \Star{} Visitáre vos in pace.\end{Resp}
\begin{Resp}\Vsign Aspíciam vos, et créscere fáciam: multiplicabímini, et firmábo pactum meum vobíscum.\end{Resp}
\begin{Resp}\Rsign Visitáre vos in pace.\end{Resp}
\begin{Resp}\Vsign Glória Patri, et Fílio, \Star{} et Spirítui Sancto.\end{Resp}
\begin{Resp}\Rsign Visitáre vos in pace.\end{Resp}
\switchcolumn
\EN
\begin{Resp}\Rsign Behold, I, the Lord your God, come from the South \Star{} To visit you in peace.\end{Resp}
\begin{Resp}\Vsign I will look again upon you and make you to increase ye shall be multiplied, and I will establish My covenant with you.\end{Resp}
\begin{Resp}\Rsign To visit you in peace.\end{Resp}
\begin{Resp}\Vsign Glory be to the Father, and to the Son, \Star{} and to the Holy Ghost.\end{Resp}
\begin{Resp}\Rsign To visit you in peace.\end{Resp}
\switchcolumn*

\end{paracol}

\Office{Feria IV infra Hebdomadam I Adventus}{Wednesday of the First Week of Advent}{Feria major}
\begin{paracol}{2}
\LA
\LessonHead{Lectio I}
\switchcolumn
\EN
\LessonHead{Lesson I}
\switchcolumn*

\LA
\LessonTitle{De Isaía Prophéta}
\switchcolumn
\EN
\LessonTitle{Lesson from the book of Isaias}
\switchcolumn*

\LA
\Cite{Isa 3:1-4}
\switchcolumn
\EN
\Cite{Isa 3:1-4}
\switchcolumn*

\LA
\noindent Ecce enim Dominátor Dóminus exercítuum áuferet a Jerúsalem et a Juda válidum et fortem, omne robur panis, et omne robur aquæ: fortem, et virum bellatórem, Júdicem, et prophétam, et aríolum, et senem: príncipem super quinquagínta, et honorábilem vultu, et consiliárium, et sapiéntem de architéctis, et prudéntem elóquii mýstici. Et dabo púeros príncipes eórum, et effemináti dominabúntur eis.\par
\switchcolumn
\EN
\noindent For behold, the sovereign the Lord of hosts shall take away from Jerusalem, and from Juda the valiant and the strong, the whole strength of bread, and the whole strength of water. The strong man, and the man of war, the judge, and the prophet, and the cunning man, and the ancient. The captain over fifty, and the honourable in countenance, and the counsellor, and the architect, and the skilful in eloquent speech. And I will give children to be their princes, and the effeminate shall rule over them.\par
\switchcolumn*

\LA
\begin{Resp}\Rsign Rex noster advéniet Christus, \Star{} Quem Joánnes prædicávit Agnum esse ventúrum.\end{Resp}
\begin{Resp}\Vsign Super ipsum continébunt reges os suum, ipsum gentes deprecabúntur.\end{Resp}
\begin{Resp}\Rsign Quem Joánnes prædicávit Agnum esse ventúrum.\end{Resp}
\switchcolumn
\EN
\begin{Resp}\Rsign Christ our King cometh. \Star{} And John hath testified of Him, that He is the Lamb that should come!\end{Resp}
\begin{Resp}\Vsign The kings shall shut their mouths at Him, all nations shall serve Him.\end{Resp}
\begin{Resp}\Rsign And John hath testified of Him, that He is the Lamb that should come!\end{Resp}
\switchcolumn*

\LA
\LessonHead{Lectio II}
\switchcolumn
\EN
\LessonHead{Lesson II}
\switchcolumn*

\LA
\Cite{Isa 3:5-7}
\switchcolumn
\EN
\Cite{Isa 3:5-7}
\switchcolumn*

\LA
\noindent Et írruet pópulus vir ad virum, et unusquísque ad próximum suum: tumultuábitur puer contra senem, et ignóbilis contra nóbilem. Apprehéndet enim vir fratrem suum domésticum patris sui: Vestiméntum tibi est, princeps esto noster, ruína autem hæc sub manu tua. Respondébit in die illa, dicens: Non sum médicus, et in domo mea non est panis, neque vestiméntum: nolíte constitúere me príncipem pópuli.\par
\switchcolumn
\EN
\noindent And the people shall rush one upon another, and every man against his neighbour: the child shall make it tumult against the ancient, and the base against the honourable. For a man shall take hold of his brother, one of the house of his father, saying: Thou hast a garment, be thou our ruler, and let this ruin be under thy hand. In that day he shall answer, saying: I am no healer, and in my house there is no bread, nor clothing: make me not ruler of the people.\par
\switchcolumn*

\LA
\begin{Resp}\Rsign Ante multum tempus prophetávit Ezéchiel: Vidi portam clausam; ecce Deus ante sǽcula ex ea procedébat pro salúte mundi: \Star{} Et erat íterum clausa, demónstrans Vírginem, quia post partum permánsit virgo.\end{Resp}
\begin{Resp}\Vsign Porta quam vidísti, Dóminus solus transíbit per illam.\end{Resp}
\begin{Resp}\Rsign Et erat íterum clausa, demónstrans Vírginem, quia post partum permánsit virgo.\end{Resp}
\switchcolumn
\EN
\begin{Resp}\Rsign Of a long time, said Ezekiel the Prophet, I saw the gate shut behold, God went forth from it before the ages for the salvation of the world. \Star{} And it was shut again, for it is a figure of the Virgin, in that after child-birth she remained a Virgin still.\end{Resp}
\begin{Resp}\Vsign The Lord alone shall enter by the gate that thou savest.\end{Resp}
\begin{Resp}\Rsign And it was shut again, for it is a figure of the Virgin, in that after child-birth she remained a Virgin still.\end{Resp}
\switchcolumn*

\LA
\LessonHead{Lectio III}
\switchcolumn
\EN
\LessonHead{Lesson III}
\switchcolumn*

\LA
\Cite{Isa 3:8-11}
\switchcolumn
\EN
\Cite{Isa 3:8-11}
\switchcolumn*

\LA
\noindent Ruit enim Jerúsalem, et Judas cóncidit: quia lingua eórum et adinventiónes eórum contra Dóminum, ut provocárent óculos majestátis ejus. Agnítio vultus eórum respóndit eis: et peccátum suum quasi Sódoma prædicavérunt, nec abscondérunt: væ ánimæ eórum, quóniam réddita sunt eis mala. Dícite justo quóniam bene, quóniam fructum adinventiónum suárum cómedet. Væ ímpio in malum: retribútio enim mánuum ejus fiet ei.\par
\switchcolumn
\EN
\noindent For Jerusalem is ruined, and Juda is fallen: because their tongue, and their devices are against the Lord, to provoke the eyes of his majesty. The shew of their countenance hath answered them: and they have proclaimed abroad their sin as Sodom, and they have not hid it: woe to their souls, for evils are rendered to them. Say to the just man that it is well, for he shall eat the fruit of his doings. Woe to the wicked unto evil: for the reward of his hands shall be given him.\par
\switchcolumn*

\LA
\begin{Resp}\Rsign Ecce dies véniunt, dicit Dóminus, et suscitábo David germen justum: et regnábit rex, et sápiens erit, et fáciet judícium et justítiam in terra: \Star{} Et hoc est nomen quod vocábunt eum: * Dóminus justus noster.\end{Resp}
\begin{Resp}\Vsign In diébus illis salvábitur Juda, et Israël habitábit confidénter.\end{Resp}
\begin{Resp}\Rsign Et hoc est nomen quod vocábunt eum:\end{Resp}
\begin{Resp}\Vsign Glória Patri, et Fílio, \Star{} et Spirítui Sancto.\end{Resp}
\begin{Resp}\Rsign Dóminus justus noster.\end{Resp}
\switchcolumn
\EN
\begin{Resp}\Rsign Behold, the days come, saith the Lord, that I will raise unto David a righteous Branch; and a King shall reign in wisdom and shall execute judgment and justice in the earth: \Star{} And this is His name whereby He shall be called: * The Lord our Righteous one.\end{Resp}
\begin{Resp}\Vsign In His days Judah shall be saved, and Israel shall dwell safely.\end{Resp}
\begin{Resp}\Rsign And this is His name whereby He shall be called.\end{Resp}
\begin{Resp}\Vsign Glory be to the Father, and to the Son, \Star{} and to the Holy Ghost.\end{Resp}
\begin{Resp}\Rsign The Lord our Righteous one.\end{Resp}
\switchcolumn*

\end{paracol}

\Office{Feria V infra Hebdomadam I Adventus}{Thursday of the First Week of Advent}{Feria major}
\begin{paracol}{2}
\LA
\LessonHead{Lectio I}
\switchcolumn
\EN
\LessonHead{Lesson I}
\switchcolumn*

\LA
\LessonTitle{De Isaía Prophéta}
\switchcolumn
\EN
\LessonTitle{Lesson from the book of Isaias}
\switchcolumn*

\LA
\Cite{Isa 4:1-3}
\switchcolumn
\EN
\Cite{Isa 4:1-3}
\switchcolumn*

\LA
\noindent Et apprehéndent septem mulíeres virum unum in die illa, dicéntes: Panem nostrum comedémus, et vestiméntis nostris operiémur: tantúmmodo invocétur nomen tuum super nos, aufer oppróbrium nostrum. In die illa erit germen Dómini in magnificéntia, et glória, et fructus terræ sublímis, et exsultátio his, qui salváti fúerint de Israël. Et erit: Omnis qui relíctus fúerit in Sion, et resíduus in Jerúsalem, sanctus vocábitur, omnis qui scriptus est in vita in Jerúsalem.\par
\switchcolumn
\EN
\noindent And in that day seven women shall take hold of one man, saying: We will eat our own bread, and wear our own apparel: only let us be called by thy name, take away our reproach. In that day the bud of the Lord shall be in magnificence and glory, and the fruit of the earth shall be high, and a great joy to them that shall have escaped of Israel. And it shall come to pass, that every one that shall be left in Sion, and that shall remain in Jerusalem, shall be called holy, every one that is written in life in Jerusalem.\par
\switchcolumn*

\LA
\begin{Resp}\Rsign Súscipe verbum, Virgo María, quod tibi a Dómino per Angelum transmíssum est: concípies et páries Deum páriter et hóminem, \Star{} Ut benedícta dicáris inter omnes mulíeres.\end{Resp}
\begin{Resp}\Vsign Páries quidem fílium, et virginitátis non patiéris detriméntum: efficiéris grávida, et eris mater semper intácta.\end{Resp}
\begin{Resp}\Rsign Ut benedícta dicáris inter omnes mulíeres.\end{Resp}
\switchcolumn
\EN
\begin{Resp}\Rsign Receive, O Virgin Mary, receive the word of the Lord, which is sent thee by His Angel; thou shalt conceive, and shalt bring forth God and Man together. \Star{} And thou shalt be called blessed among all women.\end{Resp}
\begin{Resp}\Vsign Thou shalt bring forth a son, and remain a maiden undefiled thou shalt conceive and be a Mother, still Virgin unspotted.\end{Resp}
\begin{Resp}\Rsign And thou shalt be called blessed among all women.\end{Resp}
\switchcolumn*

\LA
\LessonHead{Lectio II}
\switchcolumn
\EN
\LessonHead{Lesson II}
\switchcolumn*

\LA
\Cite{Isa 5:1-4}
\switchcolumn
\EN
\Cite{Isa 5:1-4}
\switchcolumn*

\LA
\noindent Cantábo dilécto meo cánticum patruélis mei víneæ suæ: Vínea facta est dilécto meo in cornu fílio ólei. Et sepívit eam, et lápides elégit ex illa, et plantávit eam eléctam, et ædificávit turrem in médio ejus, et tórcular exstrúxit in ea: et exspectávit ut fáceret uvas, et fecit labrúscas. Nunc ergo, habitatóres Jerúsalem, et viri Juda, judicáte inter me et víneam meam. Quid est quod débui ultra fácere víneæ meæ, et non feci ei? An quod exspectávi, ut fáceret uvas, et fecit labrúscas?\par
\switchcolumn
\EN
\noindent I will sing to my beloved the canticle of my cousin concerning his vineyard. My beloved had a vineyard on a hill in a fruitful place. And he fenced it in, and picked the stones out of it, and planted it with the choicest vines, and built a tower in the midst thereof, and set up a winepress therein: and he looked that it should bring forth grapes, and it brought forth wild grapes. And now, O ye inhabitants of Jerusalem, and ye men of Juda, judge between me and my vineyard. What is there that I ought to do more to my vineyard, that I have not done to it? was it that I looked that it should bring forth grapes, and it hath brought forth wild grapes?\par
\switchcolumn*

\LA
\begin{Resp}\Rsign Aspiciébam in visu noctis, et ecce in núbibus cæli Fílius hóminis veniébat: et datum est ei regnum, et honor: \Star{} Et omnis pópulus, tribus, et linguæ sérvient ei.\end{Resp}
\begin{Resp}\Vsign Potéstas ejus, potéstas ætérna, quæ non auferétur: et regnum ejus, quod non corrumpétur.\end{Resp}
\begin{Resp}\Rsign Et omnis pópulus, tribus, et linguæ sérvient ei.\end{Resp}
\switchcolumn
\EN
\begin{Resp}\Rsign I saw in the night visions, and, behold, the Son of man came with the clouds of heaven, and there was given Him a Kingdom, and glory; \Star{} And all people, nations, and languages shall serve Him.\end{Resp}
\begin{Resp}\Vsign His dominion is an everlasting dominion which shall not pass away, and His Kingdom that which shall not be destroyed.\end{Resp}
\begin{Resp}\Rsign And all people, nations, and languages shall serve Him.\end{Resp}
\switchcolumn*

\LA
\LessonHead{Lectio III}
\switchcolumn
\EN
\LessonHead{Lesson III}
\switchcolumn*

\LA
\Cite{Isa 5:5-7}
\switchcolumn
\EN
\Cite{Isa 5:5-7}
\switchcolumn*

\LA
\noindent Et nunc osténdam vobis quid ego fáciam víneæ meæ: áuferam sepem ejus, et erit in direptiónem: díruam macériam ejus, et erit in conculcatiónem. Et ponam eam desértam: non putábitur, et non fodiétur. Et ascéndent vepres et spinæ; et núbibus mandábo ne pluant super eam imbrem. Vínea enim Dómini exercítuum domus Israël est: et vir Juda germen ejus delectábile: et exspectávi ut fáceret judícium, et ecce iníquitas: et justítiam, et ecce clamor.\par
\switchcolumn
\EN
\noindent And now I will shew you what I will do to my vineyard. I will take away the hedge thereof, and it shall be wasted: I will break down the wall thereof, and it shall be trodden down. And I will make it desolate: it shall not be pruned, and it shall not be digged: but briers and thorns shall come up: and I will command the clouds to rain no rain upon it. For the vineyard of the Lord of hosts is the house of Israel: and the man of Juda, his pleasant plant: and I looked that he should do judgment, and behold iniquity: and do justice, and behold a cry.\par
\switchcolumn*

\LA
\begin{Resp}\Rsign Missus est Gábriel Angelus ad Maríam Vírginem desponsátam Joseph, núntians ei verbum; et expavéscit Virgo de lúmine: ne tímeas, María, invenísti grátiam apud Dóminum: \Star{} Ecce concípies, et páries, et vocábitur Altíssimi Fílius.\end{Resp}
\begin{Resp}\Vsign Dabit ei Dóminus Deus sedem David, patris ejus, et regnábit in domo Jacob in ætérnum.\end{Resp}
\begin{Resp}\Rsign Ecce concípies, et páries, et vocábitur Altíssimi Fílius.\end{Resp}
\begin{Resp}\Vsign Glória Patri, et Fílio, \Star{} et Spirítui Sancto.\end{Resp}
\begin{Resp}\Rsign Ecce concípies, et páries, et vocábitur Altíssimi Fílius.\end{Resp}
\switchcolumn
\EN
\begin{Resp}\Rsign The Angel Gabriel was sent to Mary, a Virgin espoused to Joseph, to bring unto her the word of the Lord, and when the Virgin saw the light, she was afraid. Fear not, Mary, for thou hast found grace from the Lord. \Star{} Behold, thou shalt conceive and bring forth a son, and He shall be called the Son of the Highest.\end{Resp}
\begin{Resp}\Vsign The Lord God shall give unto Him the throne of His father David, and He shall reign over the house of Jacob for ever.\end{Resp}
\begin{Resp}\Rsign Behold, thou shalt conceive, and bring forth a son, and He shall be called the Son of the Highest.\end{Resp}
\begin{Resp}\Vsign Glory be to the Father, and to the Son, \Star{} and to the Holy Ghost.\end{Resp}
\begin{Resp}\Rsign Behold, thou shalt conceive and bring forth a son, and He shall be called the Son of the Highest.\end{Resp}
\switchcolumn*

\end{paracol}

\Office{Feria VI infra Hebdomadam I Adventus}{Friday of the First Week of Advent}{Feria major}
\begin{paracol}{2}
\LA
\LessonHead{Lectio I}
\switchcolumn
\EN
\LessonHead{Lesson I}
\switchcolumn*

\LA
\LessonTitle{De Isaía Prophéta}
\switchcolumn
\EN
\LessonTitle{Lesson from the book of Isaias}
\switchcolumn*

\LA
\Cite{Isa 6:1-3}
\switchcolumn
\EN
\Cite{Isa 6:1-3}
\switchcolumn*

\LA
\noindent In anno, quo mórtuus est rex Ozías, vidi Dóminum sedéntem super sólium excélsum et elevátum: et ea, quæ sub ipso erant, replébant templum. Séraphim stabant super illud: sex alæ uni, et sex alæ álteri, duábus velábant fáciem ejus, et duábus velábant pedes ejus, et duábus volábant. Et clamábant alter ad álterum, et dicébant: Sanctus, sanctus, sanctus, Dóminus exercítuum, plena est omnis terra glória ejus.\par
\switchcolumn
\EN
\noindent In the year that king Ozias died, I saw the Lord sitting upon a throne high and elevated: and his train filled the temple. Upon it stood the seraphims: the one had six wings, and the other had six wings: with two they covered his face, and with two they covered his feet, and with two they flew. And they cried one to another, and said: Holy, holy, holy, the Lord God of hosts, all the earth is full of his glory.\par
\switchcolumn*

\LA
\begin{Resp}\Rsign Ave, María, grátia plena, Dóminus tecum: \Star{} Spíritus Sanctus supervéniet in te, et virtus Altíssimi obumbrábit tibi: quod enim ex te nascétur Sanctum, vocábitur Fílius Dei.\end{Resp}
\begin{Resp}\Vsign Quómodo fiet istud, quóniam virum non cognósco? Et respóndens Angelus, dixit ei.\end{Resp}
\begin{Resp}\Rsign Spíritus Sanctus supervéniet in te, et virtus Altíssimi obumbrábit tibi: quod enim ex te nascétur Sanctum, vocábitur Fílius Dei.\end{Resp}
\switchcolumn
\EN
\begin{Resp}\Rsign Hail, Mary, full of grace; the Lord is with thee; \Star{} The Holy Ghost shall come upon thee, and the power of the Highest shall overshadow thee; therefore, also that Holy Thing Which shall be born of thee shall be called the Son of God.\end{Resp}
\begin{Resp}\Vsign How shall this be, seeing I know not a man? And the Angel answered and said unto her,\end{Resp}
\begin{Resp}\Rsign The Holy Ghost shall come upon thee, and the power of the Highest shall overshadow thee; therefore, also that Holy Thing Which shall be born of thee shall be called the Son of God.\end{Resp}
\switchcolumn*

\LA
\LessonHead{Lectio II}
\switchcolumn
\EN
\LessonHead{Lesson II}
\switchcolumn*

\LA
\Cite{Isa 6:4-7}
\switchcolumn
\EN
\Cite{Isa 6:4-7}
\switchcolumn*

\LA
\noindent Et commóta sunt superliminária cárdinum a voce clamántis, et domus repléta est fumo. Et dixi: Væ mihi, quia tácui, quia vir pollútus lábiis ego sum, et in médio pópuli pollúta lábia habéntis ego hábito, et Regem Dóminum exercítuum vidi óculis meis. Et volávit ad me unus de Séraphim, et in manu ejus cálculus, quem fórcipe túlerat de altári. Et tétigit os meum, et dixit: Ecce tétigit hoc lábia tua, et auferétur iníquitas tua, et peccátum tuum mundábitur.\par
\switchcolumn
\EN
\noindent And the lintels of the doors were moved at the voice of him that cried, and the house was filled with smoke. And I said: Woe is me, because I have held my peace; because I am a man of unclean lips, and I dwell in the midst of a people that hath unclean lips, and I have seen with my eyes the King the Lord of hosts. And one of the seraphims flew to me, and in his hand was a live coal, which he had taken with the tongs off the altar. And he touched my mouth, and said: Behold this hath touched thy lips, and thy iniquities shall be taken away, and thy sin shall be cleansed.\par
\switchcolumn*

\LA
\begin{Resp}\Rsign Salvatórem exspectámus, Dóminum Jesum Christum, \Star{} Qui reformábit corpus humilitátis nostræ configurátum córpori claritátis suæ.\end{Resp}
\begin{Resp}\Vsign Sóbrie, et juste, et pie vivámus in hoc sǽculo, exspectántes beátam spem, et advéntum glóriæ magni Dei.\end{Resp}
\begin{Resp}\Rsign Qui reformábit corpus humilitátis nostræ configurátum córpori claritátis suæ.\end{Resp}
\switchcolumn
\EN
\begin{Resp}\Rsign We look for the Saviour, the Lord Jesus Christ; \Star{} Who shall change our vile body, that it may be fashioned like unto His glorious Body.\end{Resp}
\begin{Resp}\Vsign We should live soberly, and righteously, and godly in this present world, looking for that blessed hope, and the glorious appearing of the great God.\end{Resp}
\begin{Resp}\Rsign Who shall change our vile body, that it may be fashioned like unto His glorious Body.\end{Resp}
\switchcolumn*

\LA
\LessonHead{Lectio III}
\switchcolumn
\EN
\LessonHead{Lesson III}
\switchcolumn*

\LA
\Cite{Isa 6:8-10}
\switchcolumn
\EN
\Cite{Isa 6:8-10}
\switchcolumn*

\LA
\noindent Et audívi vocem Dómini dicéntis: Quem mittam? et quis ibit nobis? Et dixi: Ecce ego, mitte me. Et dixit: Vade, et dices pópulo huic: Audíte audiéntes, et nolíte intellégere: et vidéte visiónem, et nolíte cognóscere. Excǽca cor pópuli hujus, et aures ejus ággrava, et óculos ejus claude, ne forte vídeat óculis suis, et áuribus suis áudiat, et corde suo intéllegat, et convertátur, et sanem eum.\par
\switchcolumn
\EN
\noindent And I heard the voice of the Lord, saying: Whom shall I send? and who shall go for us? And I said: Lo, here am I, send me. And he said: Go, and thou shalt say to this people: Hearing, hear, and understand not: and see the vision, and know it not. Blind the heart of this people, and make their ears heavy, and shut their eyes: lest they see with their eyes, and hear with their ears, and understand with their heart, and be converted and I heal them.\par
\switchcolumn*

\LA
\begin{Resp}\Rsign Obsecro, Dómine, mitte quem missúrus es: vide afflictiónem pópuli tui: \Star{} Sicut locútus es, veni, * Et líbera nos.\end{Resp}
\begin{Resp}\Vsign Qui regis Israël, inténde, qui dedúcis velut ovem Joseph, qui sedes super Chérubim.\end{Resp}
\begin{Resp}\Rsign Sicut locútus es, veni.\end{Resp}
\begin{Resp}\Vsign Glória Patri, et Fílio, \Star{} et Spirítui Sancto.\end{Resp}
\begin{Resp}\Rsign Et líbera nos.\end{Resp}
\switchcolumn
\EN
\begin{Resp}\Rsign O my Lord, send I pray thee, Him Whom Thou wilt send; see the affliction of thy people. As Thou hast promised, come; \Star{} And deliver us.\end{Resp}
\begin{Resp}\Vsign Give ear, O Shepherd of Israel, Thou That leadest Joseph like a flock, Thou That sittest upon the Cherubim!\end{Resp}
\begin{Resp}\Rsign As Thou hast promised, come.\end{Resp}
\begin{Resp}\Vsign Glory be to the Father, and to the Son, \Star{} and to the Holy Ghost.\end{Resp}
\begin{Resp}\Rsign And deliver us.\end{Resp}
\switchcolumn*

\end{paracol}

\Office{Sabbato infra Hebdomadam I Adventus}{Saturday of the First Week of Advent}{Feria major}
\begin{paracol}{2}
\LA
\LessonHead{Lectio I}
\switchcolumn
\EN
\LessonHead{Lesson I}
\switchcolumn*

\LA
\LessonTitle{De Isaía Prophéta}
\switchcolumn
\EN
\LessonTitle{Lesson from the book of Isaias}
\switchcolumn*

\LA
\Cite{Isa 7:1-3}
\switchcolumn
\EN
\Cite{Isa 7:1-3}
\switchcolumn*

\LA
\noindent Et factum est in diébus Achaz, fílii Jóatham, fílii Ozíæ, regis Juda, ascéndit Rasin, rex Sýriæ, et Phácee, fílius Romelíæ, rex Israël, in Jerúsalem, ad præliándum contra eam: et non potuérunt debelláre eam. Et nuntiavérunt dómui David, dicéntes: Requiévit Sýria super Ephraim, et commótum est cor ejus, et cor pópuli ejus, sicut movéntur ligna silvárum a fácie venti. Et dixit Dóminus ad Isaíam: Egrédere in occúrsum Achaz tu, et qui derelíctus est Jasub fílius tuus, ad extrémum aquædúctus piscínæ superióris in via Agri fullónis.\par
\switchcolumn
\EN
\noindent And it came to pass in the days of Achaz the son of Joathan, the son of Ozias, king of Juda, that Basin king of Syria, and Phacee the son of Romelia king of Israel, came up to Jerusalem, to fight against it: but they could not prevail over it. And they told the house of David, saying: Syria hath rested upon Ephraim, and his heart was moved, and the heart of his people, as the trees of the woods are moved with the wind. And the Lord said to Isaias: Go forth to meet Achaz, thou and Jasub thy son that is left, to the conduit in the upper pool, in the way of the fuller's field.\par
\switchcolumn*

\LA
\begin{Resp}\Rsign Ecce virgo concípiet, et páriet fílium, dicit Dóminus: \Star{} Et vocábitur nomen ejus Admirábilis, Deus, Fortis.\end{Resp}
\begin{Resp}\Vsign Super sólium David, et super regnum ejus sedébit in ætérnum.\end{Resp}
\begin{Resp}\Rsign Et vocábitur nomen ejus Admirábilis, Deus, Fortis.\end{Resp}
\switchcolumn
\EN
\begin{Resp}\Rsign Behold, the Virgin shall conceive, and bear a son, saith the Lord; \Star{} And His name shall be called Wonderful, the Mighty God.\end{Resp}
\begin{Resp}\Vsign He shall sit upon the throne of David, and upon his kingdom for ever.\end{Resp}
\begin{Resp}\Rsign And His name shall be called Wonderful, the Mighty God.\end{Resp}
\switchcolumn*

\LA
\LessonHead{Lectio II}
\switchcolumn
\EN
\LessonHead{Lesson II}
\switchcolumn*

\LA
\Cite{Isa 7:4-6}
\switchcolumn
\EN
\Cite{Isa 7:4-6}
\switchcolumn*

\LA
\noindent Et dices ad eum: Vide ut síleas: noli timére, et cor tuum ne formídet a duábus caudis titiónum fumigántium istórum, in ira furóris Rasin, regis Sýriæ, et fílii Romelíæ: eo quod consílium iníerit contra te Sýria in malum Ephraim, et fílius Romelíæ, dicéntes: Ascendámus ad Judam, et suscitémus eum, et avellámus eum ad nos, et ponámus regem in médio ejus fílium Tábeel.\par
\switchcolumn
\EN
\noindent And thou shalt say to him: See thou be quiet: fear not, and let not thy heart be afraid of the two tails of these fire brands, smoking with the wrath of the fury of Rasin king of Syria, and of the son of Romelia. Because Syria hath taken counsel against thee, unto the evil of Ephraim and the son of Romelia, saying: Let us go up to Juda, and rouse it up, and draw it away to us, and make the son of Tabeel king in the midst thereof.\par
\switchcolumn*

\LA
\begin{Resp}\Rsign Audíte verbum Dómini, gentes, et annuntiáte illud in fínibus terræ: \Star{} Et ínsulis, quæ procul sunt, dícite: Salvátor noster advéniet.\end{Resp}
\begin{Resp}\Vsign Annuntiáte, et audítum fácite: loquímini, et clamáte.\end{Resp}
\begin{Resp}\Rsign Et ínsulis, quæ procul sunt, dícite: Salvátor noster advéniet.\end{Resp}
\switchcolumn
\EN
\begin{Resp}\Rsign Hear the word of the Lord, O ye nations, and declare it in the ends of the earth; \Star{} And in the isles afar off, and say: Our Saviour shall come.\end{Resp}
\begin{Resp}\Vsign Declare it and make it known, lift up your voice and cry aloud.\end{Resp}
\begin{Resp}\Rsign And in the isles afar off, and say: Our Saviour shall come.\end{Resp}
\switchcolumn*

\LA
\LessonHead{Lectio III}
\switchcolumn
\EN
\LessonHead{Lesson III}
\switchcolumn*

\LA
\Cite{Isa 7:10-15}
\switchcolumn
\EN
\Cite{Isa 7:10-15}
\switchcolumn*

\LA
\noindent Et adjécit Dóminus loqui ad Achaz, dicens: Pete tibi signum a Dómino Deo tuo in profúndum inférni, sive in excélsum supra. Et dixit Achaz: Non petam, et non tentábo Dóminum. Et dixit: Audíte ergo, domus David: Numquid parum vobis est, moléstos esse homínibus, quia molésti estis et Deo meo? Propter hoc dabit Dóminus ipse vobis signum. Ecce virgo concípiet, et páriet fílium, et vocábitur nomen ejus Emmánuel. Butýrum et mel cómedet, ut sciat reprobáre malum, et elígere bonum.\par
\switchcolumn
\EN
\noindent And the Lord spoke again to Achaz, saying: Ask thee a sign of the Lord thy God either unto the depth of hell, or unto the height above. And Achaz said: I will not ask, and I will not tempt the Lord. And he said: Hear ye therefore, O house of David: Is it a small thing for you to be grievous to men, that you are grievous to my God also? Therefore the Lord himself shall give you a sign. Behold a virgin shall conceive, and bear a son, and his name shall be called Emmanuel. He shall eat butter and honey, that he may know to refuse the evil, and to choose the good.\par
\switchcolumn*

\LA
\begin{Resp}\Rsign Ecce dies véniunt, dicit Dóminus, et suscitábo David germen justum: et regnábit rex, et sápiens erit, et fáciet judícium et justítiam in terra: \Star{} Et hoc est nomen quod vocábunt eum: * Dóminus justus noster.\end{Resp}
\begin{Resp}\Vsign In diébus illis salvábitur Juda, et Israël habitábit confidénter.\end{Resp}
\begin{Resp}\Rsign Et hoc est nomen quod vocábunt eum:\end{Resp}
\begin{Resp}\Vsign Glória Patri, et Fílio, \Star{} et Spirítui Sancto.\end{Resp}
\begin{Resp}\Rsign Dóminus justus noster.\end{Resp}
\switchcolumn
\EN
\begin{Resp}\Rsign Behold, the days come, saith the Lord, that I will raise unto David a righteous Branch; and a King shall reign in wisdom and shall execute judgment and justice in the earth: \Star{} And this is His name whereby He shall be called: * The Lord our Righteous one.\end{Resp}
\begin{Resp}\Vsign In His days Judah shall be saved, and Israel shall dwell safely.\end{Resp}
\begin{Resp}\Rsign And this is His name whereby He shall be called.\end{Resp}
\begin{Resp}\Vsign Glory be to the Father, and to the Son, \Star{} and to the Holy Ghost.\end{Resp}
\begin{Resp}\Rsign The Lord our Righteous one.\end{Resp}
\switchcolumn*

\end{paracol}

\Office{Dominica II Adventus}{Second Sunday of Advent}{Semiduplex II classis}
\begin{paracol}{2}
\LA
\LessonHead{Lectio I}
\switchcolumn
\EN
\LessonHead{Lesson I}
\switchcolumn*

\LA
\LessonTitle{De Isaía Prophéta}
\switchcolumn
\EN
\LessonTitle{Lesson from the book of Isaias}
\switchcolumn*

\LA
\Cite{Isa 11:1-4}
\switchcolumn
\EN
\Cite{Isa 11:1-4}
\switchcolumn*

\LA
\noindent Et egrediétur virga de radíce Jesse, et flos de radíce ejus ascéndet. Et requiéscet super eum Spíritus Dómini: spíritus sapiéntiæ et intelléctus, spíritus consílii et fortitúdinis, spíritus sciéntiæ et pietátis; Et replébit eum spíritus timóris Dómini: non secúndum visiónem oculórum judicábit, neque secúndum audítum áurium árguet: Sed judicábit in justítia páuperes, et árguet in æquitáte pro mansuétis terræ:\par
\switchcolumn
\EN
\noindent And there shall come forth a rod out of the root of Jesse, and a flower shall rise up out of his root. And the spirit of the Lord shall rest upon him: the spirit of wisdom, and of understanding, the spirit of counsel, and of fortitude, the spirit of knowledge, and of godliness. And he shall be filled with the spirit of the fear of the Lord. He shall not judge according to the sight of the eyes, nor reprove according to the hearing of the ears. But he shall judge the poor with justice, and shall reprove with equity for the meek of the earth.\par
\switchcolumn*

\LA
\begin{Resp}\Rsign Jerúsalem, cito véniet salus tua: Quare mœróre consúmeris? numquid consiliárius non est tibi, quia innovávit te dolor? \Star{} Salvábo te, et liberábo te, noli timére.\end{Resp}
\begin{Resp}\Vsign Ego enim sum Dóminus, Deus tuus, Sanctus Israël, Redémptor tuus.\end{Resp}
\begin{Resp}\Rsign Salvábo te, et liberábo te, noli timére.\end{Resp}
\switchcolumn
\EN
\begin{Resp}\Rsign Thy salvation cometh quickly, O Jerusalem; why art thou wasted with sorrow? Is there no counselor in thee, that pangs have taken thee? \Star{} Fear not, for I will save thee and deliver thee.\end{Resp}
\begin{Resp}\Vsign For I am the Lord, thy God, the Holy One of Israel, thy Saviour.\end{Resp}
\begin{Resp}\Rsign Fear not, for I will save thee, and deliver thee.\end{Resp}
\switchcolumn*

\LA
\LessonHead{Lectio II}
\switchcolumn
\EN
\LessonHead{Lesson II}
\switchcolumn*

\LA
\Cite{Isa 11:4-7}
\switchcolumn
\EN
\Cite{Isa 11:4-7}
\switchcolumn*

\LA
\noindent Et percútiet terram virga oris sui, et spíritu labiórum suórum interfíciet ímpium. Et erit justítia cíngulum lumbórum ejus: et fides cinctórium renis ejus. Habitábit lupus cum agno, et pardus cum hædo accubábit, vitúlus et leo et ovis simul morabúntur, et puer párvulus minábit eos. Vítulus et ursus pascéntur, simul requiéscent cátuli eórum, et leo quasi bos cómedet páleas.\par
\switchcolumn
\EN
\noindent And he shall strike the earth with the rod of his mouth, and with the breath of his lips he shall slay the wicked. And justice shall be the girdle of his loins: and faith the girdle of his reins. The wolf shall dwell with the lamb: and the leopard shall lie down with the kid: the calf and the lion, and the sheep shall abide together, and a little child shall lead them. The calf and the bear shall feed: their young ones shall rest together: and the lion shall eat straw like the ox.\par
\switchcolumn*

\LA
\begin{Resp}\Rsign Ecce Dóminus véniet, et omnes Sancti ejus cum eo, et erit in die illa lux magna: et exíbunt de Jerúsalem sicut aqua munda: et regnábit Dóminus in ætérnum \Star{} Super omnes gentes.\end{Resp}
\begin{Resp}\Vsign Ecce Dóminus cum virtúte véniet: et regnum in manu ejus, et potéstas, et impérium.\end{Resp}
\begin{Resp}\Rsign Super omnes gentes.\end{Resp}
\switchcolumn
\EN
\begin{Resp}\Rsign Behold, the Lord shall come, and all His saints with Him, and it shall come to pass in that day that the light shall be great; and they shall go out from Jerusalem like clean water; and the Lord shall be King for ever, \Star{} Over all the earth.\end{Resp}
\begin{Resp}\Vsign Behold, the Lord cometh with an host, and in His hand are the kingdom, and power, and dominion.\end{Resp}
\begin{Resp}\Rsign Over all the earth.\end{Resp}
\switchcolumn*

\LA
\LessonHead{Lectio III}
\switchcolumn
\EN
\LessonHead{Lesson III}
\switchcolumn*

\LA
\Cite{Isa 11:8-10}
\switchcolumn
\EN
\Cite{Isa 11:8-10}
\switchcolumn*

\LA
\noindent Et delectábitur infans ab úbere super forámine áspidis, et in cavérna réguli, qui ablactátus fúerit, manum suam mittet. Non nocébunt, et non occídent in univérso monte sancto meo: quia repléta est terra sciéntia Dómini, sicut aquæ maris operiéntes. In die illa radix Jesse, qui stat in signum populórum, ipsum gentes deprecabúntur, et erit sepúlcrum ejus gloriósum.\par
\switchcolumn
\EN
\noindent And the sucking child shall play on the hole of the asp: and the weaned child shall thrust his hand into the den of the basilisk. They shall not hurt, nor shall they kill in all my holy mountain, for the earth is filled with the knowledge of the Lord, as the covering waters of the sea. In that day the root of Jesse, who standeth for an ensign of the people, him the Gentiles shall beseech, and his sepulchre shall be glorious.\par
\switchcolumn*

\LA
\begin{Resp}\Rsign Cívitas Jerúsalem, noli flere: quóniam dóluit Dóminus super te: \Star{} Et áuferet a te omnem tribulatiónem.\end{Resp}
\begin{Resp}\Vsign Ecce Dóminus in fortitúdine véniet: et brácchium ejus dominábitur.\end{Resp}
\begin{Resp}\Rsign Et áuferet a te omnem tribulatiónem.\end{Resp}
\begin{Resp}\Vsign Glória Patri, et Fílio, \Star{} et Spirítui Sancto.\end{Resp}
\begin{Resp}\Rsign Et áuferet a te omnem tribulatiónem.\end{Resp}
\switchcolumn
\EN
\begin{Resp}\Rsign O, thou city of Jerusalem, weep not, for the Lord hath repented Him concerning thee. \Star{} And He will take away from thee all distress.\end{Resp}
\begin{Resp}\Vsign Behold, the Lord shall come with might, and His arm shall rule.\end{Resp}
\begin{Resp}\Rsign And He will take away from thee all distress.\end{Resp}
\begin{Resp}\Vsign Glory be to the Father, and to the Son, \Star{} and to the Holy Ghost.\end{Resp}
\begin{Resp}\Rsign And He will take away from thee all distress.\end{Resp}
\switchcolumn*

\LA
\LessonHead{Lectio IV}
\switchcolumn
\EN
\LessonHead{Lesson IV}
\switchcolumn*

\LA
\LessonTitle{De Expositióne sancti Hierónymi Presbýteri in Isaíam Prophétam}
\switchcolumn
\EN
\LessonTitle{From the Commentary on the Prophecies of Isaiah made by St. Jerome, Priest (at Bethlehem)}
\switchcolumn*

\LA
\Source{Liber 4. in cap. 11 Isaíæ.}
\switchcolumn
\EN
\Source{Book iv, c. xi}
\switchcolumn*

\LA
\noindent Et egrediétur virga de radíce Jesse. Usque ad principium visiónis, vel pónderis Babylónis, quod vidit Isaías, fílius Amos, omnis hæc prophetía de Christo est: quam per partes vólumus explanare, ne simul propósita atque disserta lectóris confundat memóriam. Virgam et florem de radíce Jesse ipsum Dóminum Judæi interpretántur: quod scilicet in virga regnántis poténtia, in flore pulchritúdo monstrétur.\par
\switchcolumn
\EN
\noindent And there shall come forth a rod out of the stem of Jesse. From the beginning of the Book of this Prophet till the xiiith chapter, where commenceth the vision, or burden of Babylon, the whole of the vision of Isaiah, the son of Amoz, is one continual prophecy of Christ. We must explain it part by part, for if we were to take it all at once, the memory of the reader would be confused. According to the Jewish commentators, the rod and the flower would both relate to the Lord Himself. They take the rod to mean the sceptre of His Royal dominion, and the flower the loveliness of His beauty.\par
\switchcolumn*

\LA
\begin{Resp}\Rsign Ecce véniet Dóminus, protéctor noster, Sanctus Israël, \Star{} Corónam regni habens in cápite suo.\end{Resp}
\begin{Resp}\Vsign Et dominábitur a mari usque ad mare, et a flúmine usque ad términos orbis terrárum.\end{Resp}
\begin{Resp}\Rsign Corónam regni habens in cápite suo.\end{Resp}
\switchcolumn
\EN
\begin{Resp}\Rsign Behold, there cometh the Lord, our defender, the Holy One of Israel, \Star{} Wearing a royal crown upon His head.\end{Resp}
\begin{Resp}\Vsign And His dominion shall be from sea even to sea, and from the river even to the ends of the earth.\end{Resp}
\begin{Resp}\Rsign Wearing a royal crown upon His head.\end{Resp}
\switchcolumn*

\LA
\LessonHead{Lectio V}
\switchcolumn
\EN
\LessonHead{Lesson V}
\switchcolumn*

\LA
\noindent Nos autem virgam de radíce Jesse sanctam Mariam Vírginem intelligámus, quæ nullum hábuit sibi frúticem cohæréntem, de qua et supra légimus: Ecce virgo concípiet et pariet fílium. Et florem, Dóminum Salvatórem, qui dicit in Cántico canticórum: Ego flos campi, et lílium convallium.\par
\switchcolumn
\EN
\noindent We, however, understand that the rod out of the root of Jesse signifieth the holy Virgin Mary. She was a clean stem that had as yet put forth no shoot; as we have read above Behold, the Virgin shall conceive and bear a son. (Isa. vii. 14.) And the flower we believe to mean the Lord our Redeemer, Who hath elsewhere compared Himself to a flower; I am a flower of the plain, and a lily of the valleys. (Cant. ii. 1.)\par
\switchcolumn*

\LA
\begin{Resp}\Rsign Sicut mater consolátur fílios suos, ita consolábor vos, dicit Dóminus: et de Jerúsalem civitáte quam elégi, véniet vobis auxílium: \Star{} Et vidébitis et gaudébit cor vestrum.\end{Resp}
\begin{Resp}\Vsign Dabo in Sion salútem, et in Jerúsalem glóriam meam.\end{Resp}
\begin{Resp}\Rsign Et vidébitis et gaudébit cor vestrum.\end{Resp}
\switchcolumn
\EN
\begin{Resp}\Rsign As a mother comforteth her children, so will I comfort you, saith the Lord; My help also cometh unto you out of Jerusalem, the city which I have chosen. \Star{} And when ye see this, your heart shall rejoice.\end{Resp}
\begin{Resp}\Vsign I will place salvation in Zion and in Jerusalem My glory.\end{Resp}
\begin{Resp}\Rsign And when ye shall see this, your heart shall rejoice.\end{Resp}
\switchcolumn*

\LA
\LessonHead{Lectio VI}
\switchcolumn
\EN
\LessonHead{Lesson VI}
\switchcolumn*

\LA
\noindent Super hunc ígitur florem, qui de trunco et radíce Jesse per Mariam Vírginem repénte consúrget, requiéscet Spíritus Dómini: quia in ipso complácuit omnem plenitúdinem divinitátis habitáre corporáliter: nequáquam per partes, ut in ceteris Sanctis: sed juxta Evangélium eórum, quod Hebræo sermóne conscríptum legunt Nazaræi: Descéndet super eum omnis fons Spíritus Sancti. Dóminus autem Spíritus est: et ubi Spíritus Dómini, ibi libertas.\par
\switchcolumn
\EN
\noindent The Spirit of the Lord then shall rest upon this flower; this flower which shall come forth from the stem and roots of Jesse by means of the Virgin Mary. And truly the Spirit of the Lord did rest upon our Redeemer. It is written that In Him dwelleth all the fulness of the Godhead bodily. (Col. ii. 9.) The Spirit was not shed on Him by measure, as it is upon the (4 Isa. Ixvi. 13, 14. Isa. xlvi. 13.) Saints. To Him we may apply the words of the Hebrew Gospel used by the Nazarenes; The whole fountain of the Holy Ghost shall be poured forth upon Him The Lord is a spirit, and where the Spirit of the Lord is, there is liberty. (2 Cor. iii. 17.)\par
\switchcolumn*

\LA
\begin{Resp}\Rsign Jerúsalem, plantábis víneam in móntibus tuis: exsultábis, quóniam dies Dómini véniet: surge, Sion, convértere ad Dóminum Deum tuum: gaude et lætáre, Jacob: \Star{} Quia de médio géntium Salvátor tuus véniet.\end{Resp}
\begin{Resp}\Vsign Exsúlta satis, fília Sion: júbila, fília Jerúsalem.\end{Resp}
\begin{Resp}\Rsign Quia de médio géntium Salvátor tuus véniet.\end{Resp}
\begin{Resp}\Vsign Glória Patri, et Fílio, \Star{} et Spirítui Sancto.\end{Resp}
\begin{Resp}\Rsign Quia de médio géntium Salvátor tuus véniet.\end{Resp}
\switchcolumn
\EN
\begin{Resp}\Rsign Thou shalt yet plant vines upon thy mountains, O Jerusalem thou shalt sing for joy, for the day of the Lord cometh; arise, O Zion, and turn unto the Lord thy God; rejoice and be glad, O Jacob. \Star{} For thy Saviour cometh from the midst of the nations.\end{Resp}
\begin{Resp}\Vsign Sing aloud for joy, O daughter of Zion; shout with gladness, O daughter of Jerusalem.\end{Resp}
\begin{Resp}\Rsign For thy Saviour cometh from the midst of the nations.\end{Resp}
\begin{Resp}\Vsign Glory be to the Father, and to the Son, \Star{} and to the Holy Ghost.\end{Resp}
\begin{Resp}\Rsign For thy Saviour cometh from the midst of the nations.\end{Resp}
\switchcolumn*

\LA
\LessonHead{Lectio VII}
\switchcolumn
\EN
\LessonHead{Lesson VII}
\switchcolumn*

\LA
\LessonTitle{Léctio sancti Evangélii secúndum Matthǽum}
\switchcolumn
\EN
\LessonTitle{From the Holy Gospel according to Matthew}
\switchcolumn*

\LA
\Cite{Matt 11:2-10}
\switchcolumn
\EN
\Cite{Matt 11:2-10}
\switchcolumn*

\LA
\noindent In illo témpore: Cum audísset Joánnes in vínculis ópera Christi, mittens duos de discípulis suis, ait illi: Tu es qui ventúrus es, an álium exspectámus? Et réliqua.\par
\switchcolumn
\EN
\noindent In that time, when John had heard in prison the works of Christ: sending two of his disciples he said to him: Art thou he that art to come, or look we for another? And so on.\par
\switchcolumn*

\LA
\LessonTitle{Homilía sancti Gregórii Papæ}
\switchcolumn
\EN
\LessonTitle{Homily by Pope St. Gregory (the Great.)}
\switchcolumn*

\LA
\Source{Homilia 6 in Evangelia, post initium}
\switchcolumn
\EN
\Source{6th Homily on the Gospels}
\switchcolumn*

\LA
\noindent Visis tot signis tantísque virtútibus, non scandalizári quisque pótuit, sed admirári. Sed infidélium mens grave in illo scándalum pértulit, cum eum post tot mirácula moriéntem vidit. Unde et Paulus dicit: Nos autem prædicámus Christum crucifíxum, Judǽis quidem scándalum, géntibus autem stultítiam. Stultum quippe homínibus visum est, ut pro homínibus Auctor vitæ morerétur: et inde contra eum homo scándalum sumpsit, unde ei ámplius débitor fíeri débuit. Nam tanto Deus ab homínibus dígnius honorándus est, quanto pro homínibus et indígna suscépit.\par
\switchcolumn
\EN
\noindent The sight of so many signs and so many mighty works should have been a source of wonder, and not a stumbling-block. And yet the unfaithful (Jer. xxxi. 5.) found these very works a rock of offence, when they afterwards saw Him Who had worked so many miracles dying on the Cross. Hence Paul saith: “We preach Christ crucified, unto the Jews a stumbling-block and unto the Gentiles foolishness.” (1 Cor. i. 23.) It is indeed folly in the eyes of men to say that the Author of life died for men and thus men put as a stumbling-block to hinder them from coming to Jesus, the very thing that doth oblige them the most unto Him. For the more humbling God hath undergone for man's sake, the more worthy is He that man should worship Him.\par
\switchcolumn*

\LA
\begin{Resp}\Rsign Egrediétur Dóminus de Samaría ad portam, quæ réspicit ad Oriéntem: et véniet in Béthlehem ámbulans super aquas redemptiónis Judæ: \Star{} Tunc salvus erit omnis homo: quia ecce véniet.\end{Resp}
\begin{Resp}\Vsign Et præparábitur in misericórdia sólium ejus, et sedébit super illud in veritáte.\end{Resp}
\begin{Resp}\Rsign Tunc salvus erit omnis homo: quia ecce véniet.\end{Resp}
\switchcolumn
\EN
\begin{Resp}\Rsign The Lord shall go forth out of Samaria unto the gate that looketh toward the East; and He shall come into Bethlehem, walking upon the waters of the redemption of Judah. \Star{} Then shall every one be saved for, behold, He cometh.\end{Resp}
\begin{Resp}\Vsign And in mercy shall His throne be established, and He shall sit upon it in truth.\end{Resp}
\begin{Resp}\Rsign Then shall every one be saved for, behold, He cometh.\end{Resp}
\switchcolumn*

\LA
\LessonHead{Lectio VIII}
\switchcolumn
\EN
\LessonHead{Lesson VIII}
\switchcolumn*

\LA
\noindent Quid est ergo dicere: Beátus qui non fúerit scandalizátus in me; nisi apérta voce abjectiónem mortis suæ humilitátemque signare? Ac si patenter dicat: Mira quidem facio, sed abjecta pérpeti non dedignor. Quia ergo moriéndo te súbsequor, cavéndum valde est homínibus, ne in me mortem despíciant, qui signa venerántur.\par
\switchcolumn
\EN
\noindent And blessed is he, whosoever shall not be offended in Me. Now what is this, but a plain mention of that time, when He afterwards humbled Himself, becoming obedient unto death, even the death of the Cross? It is as if He said: I indeed do wonderful works, but the day will come when I shall not refuse to suffer shame and evil treatment. Take heed then, ye who now worship Me for the works' sake, that when I come to die ye despise Me not for My death's sake.\par
\switchcolumn*

\LA
\begin{Resp}\Rsign Festína, ne tardáveris, Dómine: \Star{} Et líbera pópulum tuum.\end{Resp}
\begin{Resp}\Vsign Veni, Dómine, et noli tardáre: reláxa facínora plebi tuæ.\end{Resp}
\begin{Resp}\Rsign Et líbera pópulum tuum.\end{Resp}
\switchcolumn
\EN
\begin{Resp}\Rsign Make haste, O Lord, make no tarrying. \Star{} And deliver thy people.\end{Resp}
\begin{Resp}\Vsign O Lord, come and make no tarrying loose the bonds of thy people.\end{Resp}
\begin{Resp}\Rsign And deliver thy people.\end{Resp}
\switchcolumn*

\LA
\LessonHead{Lectio IX}
\switchcolumn
\EN
\LessonHead{Lesson IX}
\switchcolumn*

\LA
\noindent Sed dimissis Joánnis discípulis, quid de eodem Joánne turbis dicat, audiámus: Quid exístis in desértum vidére? Arúndinem vento agitatam? Quod videlicet non asseréndo , sed negando intulit. Arúndinem quippe mox ut aura contígerit, in partem álteram inflectit. Et quid per arúndinem, nisi carnalis animus designátur? Qui mox ut favore vel detractióne tangitur, statim in partem quámlibet inclinátur?\par
\switchcolumn
\EN
\noindent And, as the disciples of John departed, what did Jesus say unto the multitudes concerning this same John? Let us hear. What went ye out into the wilderness to see? A reed shaken with the wind? Here our Lord teacheth not by assertion, but by negation. Now a reed is a thing so made that as soon as the wind bloweth upon it, it bendeth it over toward the opposite quarter. And the fleshly-minded man is like a human reed. As he is praised or blamed so he bendeth himself in the one direction or the other.\par
\switchcolumn*

\LA
\begin{Resp}\Rsign Ecce Dóminus véniet cum splendóre descéndens, et virtus ejus cum eo, \Star{} Visitáre pópulum suum in pace, et constitúere super eum vitam sempitérnam.\end{Resp}
\begin{Resp}\Vsign Ecce Dóminus noster cum virtúte véniet.\end{Resp}
\begin{Resp}\Rsign Visitáre pópulum suum in pace, et constitúere super eum vitam sempitérnam.\end{Resp}
\begin{Resp}\Vsign Glória Patri, et Fílio, \Star{} et Spirítui Sancto.\end{Resp}
\begin{Resp}\Rsign Visitáre pópulum suum in pace, et constitúere super eum vitam sempitérnam.\end{Resp}
\switchcolumn
\EN
\begin{Resp}\Rsign Behold, the Lord cometh down with glory, and His host is with Him. \Star{} To visit His people in peace, and to establish them in life everlasting.\end{Resp}
\begin{Resp}\Vsign Behold, our Lord cometh with an host.\end{Resp}
\begin{Resp}\Rsign To visit His people in peace, and to establish them in life everlasting.\end{Resp}
\begin{Resp}\Vsign Glory be to the Father, and to the Son, \Star{} and to the Holy Ghost.\end{Resp}
\begin{Resp}\Rsign To visit His people in peace, and to establish them in life everlasting.\end{Resp}
\switchcolumn*

\end{paracol}

\Office{Feria II infra Hebdomadam II Adventus}{Monday of the Second Week of Advent}{Feria major}
\begin{paracol}{2}
\LA
\LessonHead{Lectio I}
\switchcolumn
\EN
\LessonHead{Lesson I}
\switchcolumn*

\LA
\LessonTitle{De Isaía Prophéta}
\switchcolumn
\EN
\LessonTitle{Lesson from the book of Isaias}
\switchcolumn*

\LA
\Cite{Isa 13:1-4}
\switchcolumn
\EN
\Cite{Isa 13:1-4}
\switchcolumn*

\LA
\noindent Onus Babylónis, quod vidit Isaías fílius Amos. Super montem caliginósum leváte signum, exaltáte vocem, leváte manum, et ingrediántur portas duces. Ego mandávi sanctificátis meis, et vocávi fortes meos in ira mea, exsultántes in glória mea. Vox multitúdinis in móntibus, quasi populórum frequéntium: vox sónitus regum, géntium congregatárum.\par
\switchcolumn
\EN
\noindent The burden of Babylon, which Isaias the son of Amos saw. Upon the dark mountain lift ye up a banner, exalt the voice, lift up the hand, and let the rulers go into the gates. I have commanded my sanctified ones, and have called my strong ones in my wrath, them that rejoice in my glory. The noise of a multitude in the mountains, as it were of many people, the noise of the sound of kings, of nations gathered together.\par
\switchcolumn*

\LA
\begin{Resp}\Rsign Súscipe verbum, Virgo María, quod tibi a Dómino per Angelum transmíssum est: concípies et páries Deum páriter et hóminem, \Star{} Ut benedícta dicáris inter omnes mulíeres.\end{Resp}
\begin{Resp}\Vsign Páries quidem fílium, et virginitátis non patiéris detriméntum: efficiéris grávida, et eris mater semper intácta.\end{Resp}
\begin{Resp}\Rsign Ut benedícta dicáris inter omnes mulíeres.\end{Resp}
\switchcolumn
\EN
\begin{Resp}\Rsign Receive, O Virgin Mary, receive the word of the Lord, which is sent thee by His Angel; thou shalt conceive, and shalt bring forth God and Man together. \Star{} And thou shalt be called blessed among all women.\end{Resp}
\begin{Resp}\Vsign Thou shalt bring forth a son, and remain a maiden undefiled thou shalt conceive and be a Mother, still Virgin unspotted.\end{Resp}
\begin{Resp}\Rsign And thou shalt be called blessed among all women.\end{Resp}
\switchcolumn*

\LA
\LessonHead{Lectio II}
\switchcolumn
\EN
\LessonHead{Lesson II}
\switchcolumn*

\LA
\Cite{Isa 13:4-8}
\switchcolumn
\EN
\Cite{Isa 13:4-8}
\switchcolumn*

\LA
\noindent Dóminus exercítuum præcépit milítiæ belli, veniéntibus de terra procul a summitáte cæli: Dóminus, et vasa furóris ejus, ut dispérdat omnem terram. Ululáte, quia prope est dies Dómini: quasi vástitas a Dómino véniet. Propter hoc omnes manus dissolvéntur, et omne cor hóminis contabéscet, et conterétur. Torsiónes et dolóres tenébunt; quasi partúriens dolébunt: unusquísque ad próximum suum stupébit, fácies combústæ vultus eórum.\par
\switchcolumn
\EN
\noindent The Lord of hosts hath given charge to the troops of war. To them that come from a country afar off, from the end of heaven: the Lord and the instruments of his wrath, to destroy the whole land. Howl ye, for the day of the Lord is near: it shall come as a destruction from the Lord. Therefore shall all hands be faint, and every heart of man shall melt, And shall be broken. Gripings and pains shall take hold of them, they shall be in pain as a woman in labour. Every one shall be amazed at his neighbour, their countenances shall be as faces burnt.\par
\switchcolumn*

\LA
\begin{Resp}\Rsign Læténtur cæli, et exsúltet terra, jubiláte, montes, laudem: quia Dóminus noster véniet, \Star{} Et páuperum suórum miserébitur.\end{Resp}
\begin{Resp}\Vsign Oriétur in diébus ejus justítia, et abundántia pacis.\end{Resp}
\begin{Resp}\Rsign Et páuperum suórum miserébitur.\end{Resp}
\switchcolumn
\EN
\begin{Resp}\Rsign Sing, O heavens; and be joyful, O earth; and break forth into singing, O mountains, for our Lord will come; \Star{} And will have mercy on His afflicted.\end{Resp}
\begin{Resp}\Vsign In His days shall righteousness flourish and abundance of peace.\end{Resp}
\begin{Resp}\Rsign And will have mercy upon His afflicted.\end{Resp}
\switchcolumn*

\LA
\LessonHead{Lectio III}
\switchcolumn
\EN
\LessonHead{Lesson III}
\switchcolumn*

\LA
\Cite{Isa 13:9-11}
\switchcolumn
\EN
\Cite{Isa 13:9-11}
\switchcolumn*

\LA
\noindent Ecce dies Dómini véniet, crudélis, et indignatiónis plenus, et iræ furorísque, ad ponéndam terram in solitúdinem, et peccatóres ejus conteréndos de ea. Quóniam stellæ cæli, et splendor eárum, non expándent lumen suum: obtenebrátus est sol in ortu suo, et luna non splendébit in lúmine suo. Et visitábo super orbis mala, et contra ímpios iniquitátem eórum, et quiéscere fáciam supérbiam infidélium, et arrogántiam fórtium humiliábo.\par
\switchcolumn
\EN
\noindent Behold, the day of the Lord shall come, a cruel day, and full of indignation, and of wrath, and fury, to lay the land desolate, and to destroy the sinners thereof out of it. For the stars of heaven, and their brightness shall not display their light: the sun shall be darkened in his rising, and the moon shall not shine with her light. And I will visit the evils of the world, and against the wicked for their iniquity: and I will make the pride of infidels to cease, and will bring down the arrogancy of the mighty.\par
\switchcolumn*

\LA
\begin{Resp}\Rsign Aliéni non transíbunt per Jerúsalem ámplius: \Star{} Nam in illa die stillábunt montes dulcédinem, et colles fluent lac et mel, dicit Dóminus.\end{Resp}
\begin{Resp}\Vsign Deus a Líbano véniet, et Sanctus de monte umbróso et condénso.\end{Resp}
\begin{Resp}\Rsign Nam in illa die stillábunt montes dulcédinem, et colles fluent lac et mel, dicit Dóminus.\end{Resp}
\begin{Resp}\Vsign Glória Patri, et Fílio, \Star{} et Spirítui Sancto.\end{Resp}
\begin{Resp}\Rsign Nam in illa die stillábunt montes dulcédinem, et colles fluent lac et mel, dicit Dóminus.\end{Resp}
\switchcolumn
\EN
\begin{Resp}\Rsign There shall no strangers pass through Jerusalem any more \Star{} For in that day the mountains shall drop down sweet wine, and the hills shall flow with milk and honey, saith the Lord.\end{Resp}
\begin{Resp}\Vsign God shall come from Lebanon, and the Holy One from the thick and shady mountain.\end{Resp}
\begin{Resp}\Rsign For in that day the mountains shall drop down sweet wine, and the hills shall flow with milk and honey, saith the Lord.\end{Resp}
\begin{Resp}\Vsign Glory be to the Father, and to the Son, \Star{} and to the Holy Ghost.\end{Resp}
\begin{Resp}\Rsign For in that day the mountains shall drop down sweet wine, and the hills shall flow with milk and honey, saith the Lord.\end{Resp}
\switchcolumn*

\end{paracol}

\Office{Feria III infra Hebdomadam II Adventus}{Tuesday of the Second Week of Advent}{Feria major}
\begin{paracol}{2}
\LA
\LessonHead{Lectio I}
\switchcolumn
\EN
\LessonHead{Lesson I}
\switchcolumn*

\LA
\LessonTitle{De Isaía Prophéta}
\switchcolumn
\EN
\LessonTitle{Lesson from the book of Isaias}
\switchcolumn*

\LA
\Cite{Isa 14:1-2}
\switchcolumn
\EN
\Cite{Isa 14:1-2}
\switchcolumn*

\LA
\noindent Prope est ut véniat tempus ejus, et dies ejus non elongabúntur. Miserébitur enim Dóminus Jacob, et éliget adhuc de Israël, et requiéscere eos fáciet super humum suam: adjungétur ádvena ad eos, et adhærébit dómui Jacob. Et tenébunt eos pópuli, et addúcent eos in locum suum: et possidébit eos domus Israël super terram Dómini in servos et ancíllas: et erunt capiéntes eos, qui se céperant, et subícient exactóres suos.\par
\switchcolumn
\EN
\noindent Her time is near at hand, and her days shall not be prolonged. For the Lord will have mercy on Jacob, and will yet choose out of Israel, and will make them rest upon their own ground: and the stranger shall be joined with them, and shall adhere to the house of Jacob. And the people shall take them, and bring them into their place: and the house of Israel shall possess them in the land of the Lord for servants and handmaids: and they shall make them captives that had taken them, and shall subdue their oppressors.\par
\switchcolumn*

\LA
\begin{Resp}\Rsign Montes Israël, ramos vestros expándite, et floréte, et fructus fácite: \Star{} Prope est ut véniat dies Dómini.\end{Resp}
\begin{Resp}\Vsign Roráte, cæli, désuper, et nubes pluant justum: aperiátur terra, et gérminet Salvatórem.\end{Resp}
\begin{Resp}\Rsign Prope est ut véniat dies Dómini.\end{Resp}
\switchcolumn
\EN
\begin{Resp}\Rsign O ye mountains of Israel, shoot forth your branches and blossom and bring forth fruit. \Star{} The day of the Lord is at hand to come.\end{Resp}
\begin{Resp}\Vsign Drop down, ye heavens, from above, and let the skies pour down the Righteous One; let the earth open, and let her bring forth the Saviour.\end{Resp}
\begin{Resp}\Rsign The day of the Lord is at hand to come.\end{Resp}
\switchcolumn*

\LA
\LessonHead{Lectio II}
\switchcolumn
\EN
\LessonHead{Lesson II}
\switchcolumn*

\LA
\Cite{Isa 14:3-6}
\switchcolumn
\EN
\Cite{Isa 14:3-6}
\switchcolumn*

\LA
\noindent Et erit in die illa, cum réquiem déderit tibi Deus a labóre tuo, et a concussióne tua, et a servitúte dura, qua ante servísti: sumes parábolam istam contra regem Babylónis, et dices: Quómodo cessávit exáctor, quiévit tribútum? Contrívit Dóminus báculum impiórum, virgam dominántium, cædéntem pópulos in indignatióne, plaga insanábili, subiciéntem in furóre gentes, persequéntem crudéliter.\par
\switchcolumn
\EN
\noindent And it shall come to pass in that day, that when God shall give thee rest from thy labour, and from thy vexation, and from the hard bondage, wherewith thou didst serve before, Thou shalt take up this parable against the king of Babylon, and shalt say: How is the oppressor come to nothing, the tribute hath ceased? The Lord hath broken the staff of the wicked, the rod of the rulers, That struck the people in wrath with an incurable wound, that brought nations under in fury, that persecuted in a cruel manner.\par
\switchcolumn*

\LA
\begin{Resp}\Rsign Erúmpant montes jucunditátem, et colles justítiam: \Star{} Quia lux mundi Dóminus cum poténtia venit.\end{Resp}
\begin{Resp}\Vsign De Sion exíbit lex, et verbum Dómini de Jerúsalem.\end{Resp}
\begin{Resp}\Rsign Quia lux mundi Dóminus cum poténtia venit.\end{Resp}
\switchcolumn
\EN
\begin{Resp}\Rsign Let the mountains break forth into singing, and the hills bring forth righteousness \Star{} For the Lord, the Light of the world, cometh with power.\end{Resp}
\begin{Resp}\Vsign Out of Zion shall go forth the law, and the word of the Lord from Jerusalem.\end{Resp}
\begin{Resp}\Rsign For the Lord, the Light of the world, cometh with power.\end{Resp}
\switchcolumn*

\LA
\LessonHead{Lectio III}
\switchcolumn
\EN
\LessonHead{Lesson III}
\switchcolumn*

\LA
\Cite{Isa 14:12-15}
\switchcolumn
\EN
\Cite{Isa 14:12-15}
\switchcolumn*

\LA
\noindent Quómodo cecidísti de cælo, lúcifer, qui mane oriebáris? corruísti in terram, qui vulnerábas gentes? qui dicébas in corde tuo: In cælum conscéndam, super astra Dei exaltábo sólium meum, sedébo in monte testaménti, in latéribus Aquilónis. Ascéndam super altitúdinem núbium, símilis ero Altíssimo. Verúmtamen ad inférnum detrahéris in profúndum laci.\par
\switchcolumn
\EN
\noindent How art thou fallen from heaven, O Lucifer, who didst rise in the morning? how art thou fallen to the earth, that didst wound the nations? And thou saidst in thy heart: I will ascend into heaven, I will exalt my throne above the stars of God, I will sit in the mountain of the covenant, in the sides of the north. I will ascend above the height of the clouds, I will be like the most High. But yet thou shalt be brought down to hell, into the depth of the pit.\par
\switchcolumn*

\LA
\begin{Resp}\Rsign Ecce ab Austro vénio, ego Dóminus Deus vester, \Star{} Visitáre vos in pace.\end{Resp}
\begin{Resp}\Vsign Aspíciam vos, et créscere fáciam: multiplicabímini, et firmábo pactum meum vobíscum.\end{Resp}
\begin{Resp}\Rsign Visitáre vos in pace.\end{Resp}
\begin{Resp}\Vsign Glória Patri, et Fílio, \Star{} et Spirítui Sancto.\end{Resp}
\begin{Resp}\Rsign Visitáre vos in pace.\end{Resp}
\switchcolumn
\EN
\begin{Resp}\Rsign Behold, I, the Lord your God, come from the South \Star{} To visit you in peace.\end{Resp}
\begin{Resp}\Vsign I will look again upon you and make you to increase ye shall be multiplied, and I will establish My covenant with you.\end{Resp}
\begin{Resp}\Rsign To visit you in peace.\end{Resp}
\begin{Resp}\Vsign Glory be to the Father, and to the Son, \Star{} and to the Holy Ghost.\end{Resp}
\begin{Resp}\Rsign To visit you in peace.\end{Resp}
\switchcolumn*

\end{paracol}

\Office{Feria IV infra Hebdomadam II Adventus}{Wednesday of the Second Week of Advent}{Feria major}
\begin{paracol}{2}
\LA
\LessonHead{Lectio I}
\switchcolumn
\EN
\LessonHead{Lesson I}
\switchcolumn*

\LA
\LessonTitle{De Isaía Prophéta}
\switchcolumn
\EN
\LessonTitle{Lesson from the book of Isaias}
\switchcolumn*

\LA
\Cite{Isa 16:1-4}
\switchcolumn
\EN
\Cite{Isa 16:1-4}
\switchcolumn*

\LA
\noindent Emítte Agnum, Dómine, Dominatórem terræ de Petra desérti ad montem fíliæ Sion. Et erit: Sicut avis fúgiens, et pulli de nido avolántes, sic erunt fíliæ Moab in transcénsu Arnon. Ini consílium, coge concílium: pone quasi noctem umbram tuam in merídie: abscónde fugiéntes, et vagos ne prodas. Habitábunt apud te prófugi mei: Moab esto latíbulum eórum a fácie vastatóris.\par
\switchcolumn
\EN
\noindent Send forth, O Lord, the lamb, the ruler of the earth, from Petra of the desert, to the mount of the daughter of Sion. And it shall come to pass, that as a bird fleeing away, and as young ones flying out of the nest, so shall the daughters of Moab be in the passage of Arnon. Take counsel, gather a council: make thy shadow as the night in the midday: hide them that flee, and betray not them that wander about. My fugitives shall dwell with thee: O Moab, be thou a covert to them from the face of the destroyer.\par
\switchcolumn*

\LA
\begin{Resp}\Rsign Rex noster advéniet Christus, \Star{} Quem Joánnes prædicávit Agnum esse ventúrum.\end{Resp}
\begin{Resp}\Vsign Super ipsum continébunt reges os suum, ipsum gentes deprecabúntur.\end{Resp}
\begin{Resp}\Rsign Quem Joánnes prædicávit Agnum esse ventúrum.\end{Resp}
\switchcolumn
\EN
\begin{Resp}\Rsign Christ our King cometh. \Star{} And John hath testified of Him, that He is the Lamb that should come!\end{Resp}
\begin{Resp}\Vsign The kings shall shut their mouths at Him, all nations shall serve Him.\end{Resp}
\begin{Resp}\Rsign And John hath testified of Him, that He is the Lamb that should come!\end{Resp}
\switchcolumn*

\LA
\LessonHead{Lectio II}
\switchcolumn
\EN
\LessonHead{Lesson II}
\switchcolumn*

\LA
\Cite{Isa 16:4-6}
\switchcolumn
\EN
\Cite{Isa 16:4-6}
\switchcolumn*

\LA
\noindent Fínitus est enim pulvis, consummátus est miser, defécit qui conculcábat terram. Et præparábitur in misericórdia sólium, et sedébit super illud in veritáte in tabernáculo David, júdicans et quærens judícium, et velóciter reddens quod justum est. Audívimus supérbiam Moab, supérbus est valde: supérbia ejus et arrogántia ejus, et indignátio ejus, plus quam fortitúdo ejus.\par
\switchcolumn
\EN
\noindent For the dust is at an end, the wretch is consumed: he hath failed, that trod the earth under foot. And a throne shall be prepared in mercy, and one shall sit upon it in truth in the tabernacle of David, judging and seeking judgment and quickly rendering that which is just. We have heard of the pride of Moab, he is exceeding proud: his pride and his arrogancy, and his indignation is more than his strength.\par
\switchcolumn*

\LA
\begin{Resp}\Rsign Ante multum tempus prophetávit Ezéchiel: Vidi portam clausam; ecce Deus ante sǽcula ex ea procedébat pro salúte mundi: \Star{} Et erat íterum clausa, demónstrans Vírginem, quia post partum permánsit virgo.\end{Resp}
\begin{Resp}\Vsign Porta quam vidísti, Dóminus solus transíbit per illam.\end{Resp}
\begin{Resp}\Rsign Et erat íterum clausa, demónstrans Vírginem, quia post partum permánsit virgo.\end{Resp}
\switchcolumn
\EN
\begin{Resp}\Rsign Of a long time, said Ezekiel the Prophet, I saw the gate shut behold, God went forth from it before the ages for the salvation of the world. \Star{} And it was shut again, for it is a figure of the Virgin, in that after child-birth she remained a Virgin still.\end{Resp}
\begin{Resp}\Vsign The Lord alone shall enter by the gate that thou savest.\end{Resp}
\begin{Resp}\Rsign And it was shut again, for it is a figure of the Virgin, in that after child-birth she remained a Virgin still.\end{Resp}
\switchcolumn*

\LA
\LessonHead{Lectio III}
\switchcolumn
\EN
\LessonHead{Lesson III}
\switchcolumn*

\LA
\Cite{Isa 16:7-8}
\switchcolumn
\EN
\Cite{Isa 16:7-8}
\switchcolumn*

\LA
\noindent Idcírco ululábit Moab ad Moab, univérsus ululábit: his, qui lætántur super muros cocti láteris, loquímini plagas suas. Quóniam suburbána Hésebon desérta sunt, et víneam Sábama dómini géntium excidérunt: flagélla ejus usque ad Jazer pervenérunt: erravérunt in desérto, propágines ejus relictæ sunt, transiérunt mare.\par
\switchcolumn
\EN
\noindent Therefore shall Moab howl to Moab, every one shall howl: to them that rejoice upon the brick walls, tell ye their stripes. For the suburbs of Hesebon are desolate, and the lords of the nations have destroyed the vineyard of Sabama: the branches thereof have reached even to Jazer: they have wandered in the wilderness, the branches thereof are left, they are gone over the sea.\par
\switchcolumn*

\LA
\begin{Resp}\Rsign Ecce Dóminus véniet cum splendóre descéndens, et virtus ejus cum eo, \Star{} Visitáre pópulum suum in pace, et constitúere super eum vitam sempitérnam.\end{Resp}
\begin{Resp}\Vsign Ecce Dóminus noster cum virtúte véniet.\end{Resp}
\begin{Resp}\Rsign Visitáre pópulum suum in pace, et constitúere super eum vitam sempitérnam.\end{Resp}
\begin{Resp}\Vsign Glória Patri, et Fílio, \Star{} et Spirítui Sancto.\end{Resp}
\begin{Resp}\Rsign Visitáre pópulum suum in pace, et constitúere super eum vitam sempitérnam.\end{Resp}
\switchcolumn
\EN
\begin{Resp}\Rsign Behold, the Lord cometh down with glory, and His host is with Him. \Star{} To visit His people in peace, and to establish them in life everlasting.\end{Resp}
\begin{Resp}\Vsign Behold, our Lord cometh with an host.\end{Resp}
\begin{Resp}\Rsign To visit His people in peace, and to establish them in life everlasting.\end{Resp}
\begin{Resp}\Vsign Glory be to the Father, and to the Son, \Star{} and to the Holy Ghost.\end{Resp}
\begin{Resp}\Rsign To visit His people in peace, and to establish them in life everlasting.\end{Resp}
\switchcolumn*

\end{paracol}

\Office{Feria V infra Hebdomadam II Adventus}{Thursday of the Second Week of Advent}{Feria major}
\begin{paracol}{2}
\LA
\LessonHead{Lectio I}
\switchcolumn
\EN
\LessonHead{Lesson I}
\switchcolumn*

\LA
\LessonTitle{De Isaía Prophéta}
\switchcolumn
\EN
\LessonTitle{Lesson from the book of Isaias}
\switchcolumn*

\LA
\Cite{Isa 19:1-2}
\switchcolumn
\EN
\Cite{Isa 19:1-2}
\switchcolumn*

\LA
\noindent Onus Ægýpti. Ecce Dóminus ascéndet super nubem levem, et ingrediétur Ægýptum, et commovebúntur simulácra Ægýpti a fácie ejus, et cor Ægýpti tabéscet in médio ejus. Et concúrrere fáciam Ægýptios advérsus Ægýptios: et pugnábit vir contra fratrem suum, et vir contra amícum suum, cívitas advérsus civitátem, regnum advérsus regnum.\par
\switchcolumn
\EN
\noindent The burden of Egypt. Behold the Lord will ascend upon a swift cloud, and will enter into Egypt, and the idols of Egypt shall be moved at his presence, and the heart of Egypt shall melt in the midst thereof. And I will set the Egyptians to fight against the Egyptians: and they shall fight brother against brother, and friend against friend, city against city, kingdom against kingdom.\par
\switchcolumn*

\LA
\begin{Resp}\Rsign Jerúsalem, cito véniet salus tua: Quare mœróre consúmeris? numquid consiliárius non est tibi, quia innovávit te dolor? \Star{} Salvábo te, et liberábo te, noli timére.\end{Resp}
\begin{Resp}\Vsign Ego enim sum Dóminus, Deus tuus, Sanctus Israël, Redémptor tuus.\end{Resp}
\begin{Resp}\Rsign Salvábo te, et liberábo te, noli timére.\end{Resp}
\switchcolumn
\EN
\begin{Resp}\Rsign Thy salvation cometh quickly, O Jerusalem; why art thou wasted with sorrow? Is there no counselor in thee, that pangs have taken thee? \Star{} Fear not, for I will save thee and deliver thee.\end{Resp}
\begin{Resp}\Vsign For I am the Lord, thy God, the Holy One of Israel, thy Saviour.\end{Resp}
\begin{Resp}\Rsign Fear not, for I will save thee, and deliver thee.\end{Resp}
\switchcolumn*

\LA
\LessonHead{Lectio II}
\switchcolumn
\EN
\LessonHead{Lesson II}
\switchcolumn*

\LA
\Cite{Isa 19:3-6}
\switchcolumn
\EN
\Cite{Isa 19:3-6}
\switchcolumn*

\LA
\noindent Et dirumpétur spíritus Ægýpti in viscéribus ejus, et consílium ejus præcipitábo: et interrogábunt simulácra sua, et divínos suos, et pythónes, et aríolos. Et tradam Ægýptum in manu dominórum crudélium, et rex fortis dominábitur eórum, ait Dóminus Deus exercítuum. Et aréscet aqua de mari, et flúvius desolábitur, atque siccábitur. Et defícient flúmina: attenuabúntur, et siccabúntur rivi ággerum.\par
\switchcolumn
\EN
\noindent And the spirit of Egypt shall be broken in the bowels thereof, and I will cast down their counsel: and they shall consult their idols, and their diviners, and their wizards, and soothsayers. And I will deliver Egypt into the hand of cruel masters, and a strong king shall rule over them, saith the Lord the God of hosts. And the water of the sea shall be dried up, and the river shall be wasted and dry. And the rivers shall fail: the streams of the banks shall be diminished, and be dried up. The reed and the bulrush shall wither away.\par
\switchcolumn*

\LA
\begin{Resp}\Rsign Ecce Dóminus véniet, et omnes Sancti ejus cum eo, et erit in die illa lux magna: et exíbunt de Jerúsalem sicut aqua munda: et regnábit Dóminus in ætérnum \Star{} Super omnes gentes.\end{Resp}
\begin{Resp}\Vsign Ecce Dóminus cum virtúte véniet: et regnum in manu ejus, et potéstas, et impérium.\end{Resp}
\begin{Resp}\Rsign Super omnes gentes.\end{Resp}
\switchcolumn
\EN
\begin{Resp}\Rsign Behold, the Lord shall come, and all His saints with Him, and it shall come to pass in that day that the light shall be great; and they shall go out from Jerusalem like clean water; and the Lord shall be King for ever, \Star{} Over all the earth.\end{Resp}
\begin{Resp}\Vsign Behold, the Lord cometh with an host, and in His hand are the kingdom, and power, and dominion.\end{Resp}
\begin{Resp}\Rsign Over all the earth.\end{Resp}
\switchcolumn*

\LA
\LessonHead{Lectio III}
\switchcolumn
\EN
\LessonHead{Lesson III}
\switchcolumn*

\LA
\Cite{Isa 19:11-13}
\switchcolumn
\EN
\Cite{Isa 19:11-13}
\switchcolumn*

\LA
\noindent Stulti príncipes Táneos, sapiéntes consiliárii Pharaónis dedérunt consílium insípiens. Quómodo dicétis Pharaóni: Fílius sapiéntium ego, fílius regum antiquórum? Ubi nunc sunt sapiéntes tui? annúntient tibi, et índicent quid cogitáverit Dóminus exercítuum super Ægýptum. Stulti facti sunt príncipes Táneos, emarcuérunt príncipes Mémpheos, decepérunt Ægýptum, ángulum populórum ejus.\par
\switchcolumn
\EN
\noindent The princes of Tanis are become fools, the wise counsellors of Pharao have given foolish counsel: how will you say to Pharao: I am the son of the wise, the son of ancient kings? Where are now thy wise men? let them tell thee, and shew what the Lord of hosts hath purposed upon Egypt. The princes of Tanis are become fools, the princes of Memphis are gone astray, they have deceived Egypt, the stay of the people thereof.\par
\switchcolumn*

\LA
\begin{Resp}\Rsign Cívitas Jerúsalem, noli flere: quóniam dóluit Dóminus super te: \Star{} Et áuferet a te omnem tribulatiónem.\end{Resp}
\begin{Resp}\Vsign Ecce Dóminus in fortitúdine véniet: et brácchium ejus dominábitur.\end{Resp}
\begin{Resp}\Rsign Et áuferet a te omnem tribulatiónem.\end{Resp}
\begin{Resp}\Vsign Glória Patri, et Fílio, \Star{} et Spirítui Sancto.\end{Resp}
\begin{Resp}\Rsign Et áuferet a te omnem tribulatiónem.\end{Resp}
\switchcolumn
\EN
\begin{Resp}\Rsign O, thou city of Jerusalem, weep not, for the Lord hath repented Him concerning thee. \Star{} And He will take away from thee all distress.\end{Resp}
\begin{Resp}\Vsign Behold, the Lord shall come with might, and His arm shall rule.\end{Resp}
\begin{Resp}\Rsign And He will take away from thee all distress.\end{Resp}
\begin{Resp}\Vsign Glory be to the Father, and to the Son, \Star{} and to the Holy Ghost.\end{Resp}
\begin{Resp}\Rsign And He will take away from thee all distress.\end{Resp}
\switchcolumn*

\end{paracol}

\Office{Feria VI infra Hebdomadam II Adventus}{Friday of the Second Week of Advent}{Feria major}
\begin{paracol}{2}
\LA
\LessonHead{Lectio I}
\switchcolumn
\EN
\LessonHead{Lesson I}
\switchcolumn*

\LA
\LessonTitle{De Isaía Prophéta}
\switchcolumn
\EN
\LessonTitle{Lesson from the book of Isaias}
\switchcolumn*

\LA
\Cite{Isa 24:1-3}
\switchcolumn
\EN
\Cite{Isa 24:1-3}
\switchcolumn*

\LA
\noindent Ecce Dóminus dissipábit terram, et nudábit eam, et afflíget fáciem ejus, et dispérget habitatóres ejus. Et erit sicut pópulus, sic sacérdos: et sicut servus, sic dóminus ejus: sicut ancílla, sic dómina ejus: sicut emens, sic ille qui vendit: sicut fœnerátor, sic is qui mútuum áccipit: sicut qui répetit, sic qui debet. Dissipatióne dissipábitur terra, et direptióne prædábitur. Dóminus enim locútus est verbum hoc.\par
\switchcolumn
\EN
\noindent Behold the Lord shall lay waste the earth, and shall strip it, and shall afflict the face thereof, and scatter abroad the inhabitants thereof. And it shall be as with the people, so with the priest: and as with the servant, so with his master: as with the handmaid, so with her mistress: as with the buyer, so with the seller: as with the lender, so with the borrower: as with him that calleth for his money, so with him that oweth. With desolation shall the earth be laid waste, and it shall be utterly spoiled: for the Lord hath spoken this word.\par
\switchcolumn*

\LA
\begin{Resp}\Rsign Ecce véniet Dóminus, protéctor noster, Sanctus Israël, \Star{} Corónam regni habens in cápite suo.\end{Resp}
\begin{Resp}\Vsign Et dominábitur a mari usque ad mare, et a flúmine usque ad términos orbis terrárum.\end{Resp}
\begin{Resp}\Rsign Corónam regni habens in cápite suo.\end{Resp}
\switchcolumn
\EN
\begin{Resp}\Rsign Behold, there cometh the Lord, our defender, the Holy One of Israel, \Star{} Wearing a royal crown upon His head.\end{Resp}
\begin{Resp}\Vsign And His dominion shall be from sea even to sea, and from the river even to the ends of the earth.\end{Resp}
\begin{Resp}\Rsign Wearing a royal crown upon His head.\end{Resp}
\switchcolumn*

\LA
\LessonHead{Lectio II}
\switchcolumn
\EN
\LessonHead{Lesson II}
\switchcolumn*

\LA
\Cite{Isa 24:4-6}
\switchcolumn
\EN
\Cite{Isa 24:4-6}
\switchcolumn*

\LA
\noindent Luxit et deflúxit terra, et infirmáta est: deflúxit orbis, infirmáta est altitúdo pópuli terræ. Et terra infécta est ab habitatóribus suis: quia transgréssi sunt leges, mutavérunt jus, dissipavérunt fœdus sempitérnum. Propter hoc maledíctio vorábit terram, et peccábunt habitatóres ejus: ideóque insánient cultóres ejus, et relinquéntur hómines pauci.\par
\switchcolumn
\EN
\noindent The earth mourned, and faded away, and is weakened: the world faded away, the height of the people of the earth is weakened. And the earth is infected by the inhabitants thereof: because they have transgressed the laws, they have changed the ordinance, they have broken the everlasting covenant. Therefore shall a curse devour the earth, and the inhabitants thereof shall sin: and therefore they that dwell therein shall be mad, and few men shall be left.\par
\switchcolumn*

\LA
\begin{Resp}\Rsign Sicut mater consolátur fílios suos, ita consolábor vos, dicit Dóminus: et de Jerúsalem civitáte quam elégi, véniet vobis auxílium: \Star{} Et vidébitis et gaudébit cor vestrum.\end{Resp}
\begin{Resp}\Vsign Dabo in Sion salútem, et in Jerúsalem glóriam meam.\end{Resp}
\begin{Resp}\Rsign Et vidébitis et gaudébit cor vestrum.\end{Resp}
\switchcolumn
\EN
\begin{Resp}\Rsign As a mother comforteth her children, so will I comfort you, saith the Lord; My help also cometh unto you out of Jerusalem, the city which I have chosen. \Star{} And when ye see this, your heart shall rejoice.\end{Resp}
\begin{Resp}\Vsign I will place salvation in Zion and in Jerusalem My glory.\end{Resp}
\begin{Resp}\Rsign And when ye shall see this, your heart shall rejoice.\end{Resp}
\switchcolumn*

\LA
\LessonHead{Lectio III}
\switchcolumn
\EN
\LessonHead{Lesson III}
\switchcolumn*

\LA
\Cite{Isa 24:7-16}
\switchcolumn
\EN
\Cite{Isa 24:7-16}
\switchcolumn*

\LA
\noindent Luxit vindémia, infirmáta est vitis, ingemuérunt omnes qui lætabántur corde. Cessávit gáudium tympanórum, quiévit sónitus lætántium, contícuit dulcédo cítharæ. Cum cántico non bibent vinum, amára erit pótio bibéntibus illam. Attríta est cívitas vanitátis, clausa est omnis domus, nullo introëúnte. Clamor erit super vino in platéis: desérta est omnis lætítia: translátum est gáudium terræ. Relícta est in urbe solitudo, et calámitas ópprimet portas. Quia hæc erunt in médio terræ, in médio populórum: quómodo si paucæ olívæ, quæ remansérunt, excutiántur ex ólea: et racémi, cum fúerit fínita vindémia. Hi levábunt vocem suam, atque laudábunt: cum glorificátus fúerit Dóminus, hínnient de mari. Propter hoc in doctrínis glorificáte Dóminum, in ínsulis maris nomen Dómini Dei Israël. A fínibus terræ laudes audívimus, glóriam justi.\par
\switchcolumn
\EN
\noindent The vintage hath mourned, the vine hath languished away, all the merryhearted have sighed. The mirth of timbrels hath ceased, the noise of them that rejoice is ended, the melody of the harp is silent. They shall not drink wine with a song: the drink shall be bitter to them that drink it. The city of vanity is broken down, every house is shut up, no man cometh in. There shall be a crying for wine in the streets: all mirth is forsaken: the joy of the earth is gone away. Desolation is left in the city, and calamity shall oppress the gates. For it shall be thus in the midst of the earth, in the midst of the people, as if a few olives, that remain, should be shaken out of the olive tree: or grapes, when the vintage is ended. These shall lift up their voice, and shall give praise: when the Lord shall be glorified, they shall make a joyful noise from the sea. Therefore glorify ye the Lord in instruction: the name of the Lord God of Israel in the islands of the sea. From the ends of the earth we have heard praises, the glory of the Just one.\par
\switchcolumn*

\LA
\begin{Resp}\Rsign Jerúsalem, plantábis víneam in móntibus tuis: exsultábis, quóniam dies Dómini véniet: surge, Sion, convértere ad Dóminum Deum tuum: gaude et lætáre, Jacob: \Star{} Quia de médio géntium Salvátor tuus véniet.\end{Resp}
\begin{Resp}\Vsign Exsúlta satis, fília Sion: júbila, fília Jerúsalem.\end{Resp}
\begin{Resp}\Rsign Quia de médio géntium Salvátor tuus véniet.\end{Resp}
\begin{Resp}\Vsign Glória Patri, et Fílio, \Star{} et Spirítui Sancto.\end{Resp}
\begin{Resp}\Rsign Quia de médio géntium Salvátor tuus véniet.\end{Resp}
\switchcolumn
\EN
\begin{Resp}\Rsign Thou shalt yet plant vines upon thy mountains, O Jerusalem thou shalt sing for joy, for the day of the Lord cometh; arise, O Zion, and turn unto the Lord thy God; rejoice and be glad, O Jacob. \Star{} For thy Saviour cometh from the midst of the nations.\end{Resp}
\begin{Resp}\Vsign Sing aloud for joy, O daughter of Zion; shout with gladness, O daughter of Jerusalem.\end{Resp}
\begin{Resp}\Rsign For thy Saviour cometh from the midst of the nations.\end{Resp}
\begin{Resp}\Vsign Glory be to the Father, and to the Son, \Star{} and to the Holy Ghost.\end{Resp}
\begin{Resp}\Rsign For thy Saviour cometh from the midst of the nations.\end{Resp}
\switchcolumn*

\end{paracol}

\Office{Sabbato infra Hebdomadam II Adventus}{Saturday of the Second Week of Advent}{Feria major}
\begin{paracol}{2}
\LA
\LessonHead{Lectio I}
\switchcolumn
\EN
\LessonHead{Lesson I}
\switchcolumn*

\LA
\LessonTitle{De Isaía Prophéta}
\switchcolumn
\EN
\LessonTitle{Lesson from the book of Isaias}
\switchcolumn*

\LA
\Cite{Isa 25:1-4}
\switchcolumn
\EN
\Cite{Isa 25:1-4}
\switchcolumn*

\LA
\noindent Dómine, Deus meus es tu; exaltábo te, et confitébor nómini tuo: quóniam fecísti mirabília, cogitatiónes antíquas fidéles. Amen. Quia posuísti civitátem in túmulum, urbem fortem in ruínam, domum alienórum: ut non sit cívitas, et in sempitérnum non ædificétur. Super hoc laudábit te pópulus fortis; cívitas géntium robustárum timébit te: quia factus es fortitúdo páuperi, fortitúdo egéno in tribulatióne sua, spes a túrbine, umbráculum ab æstu.\par
\switchcolumn
\EN
\noindent O Lord, thou art my God, I will exalt thee, and give glory to thy name: for thou hast done wonderful things, thy designs of old faithful, amen. For thou hast reduced the city to a heap, the strong city to ruin, the house of strangers, to be no city, and to be no more built up for ever. Therefore shall a strong people praise thee, the city of mighty nations shall fear thee. Because thou hast been a strength to the poor, a strength to the needy in his distress: a refuge from the whirlwind, a shadow from the heat.\par
\switchcolumn*

\LA
\begin{Resp}\Rsign Egrediétur Dóminus de Samaría ad portam, quæ réspicit ad Oriéntem: et véniet in Béthlehem ámbulans super aquas redemptiónis Judæ: \Star{} Tunc salvus erit omnis homo: quia ecce véniet.\end{Resp}
\begin{Resp}\Vsign Et præparábitur in misericórdia sólium ejus, et sedébit super illud in veritáte.\end{Resp}
\begin{Resp}\Rsign Tunc salvus erit omnis homo: quia ecce véniet.\end{Resp}
\switchcolumn
\EN
\begin{Resp}\Rsign The Lord shall go forth out of Samaria unto the gate that looketh toward the East; and He shall come into Bethlehem, walking upon the waters of the redemption of Judah. \Star{} Then shall every one be saved for, behold, He cometh.\end{Resp}
\begin{Resp}\Vsign And in mercy shall His throne be established, and He shall sit upon it in truth.\end{Resp}
\begin{Resp}\Rsign Then shall every one be saved for, behold, He cometh.\end{Resp}
\switchcolumn*

\LA
\LessonHead{Lectio II}
\switchcolumn
\EN
\LessonHead{Lesson II}
\switchcolumn*

\LA
\Cite{Isa 25:4-7}
\switchcolumn
\EN
\Cite{Isa 25:4-7}
\switchcolumn*

\LA
\noindent Spíritus enim robustórum quasi turbo impéllens paríetem. Sicut æstus in siti, tumúltum alienórum humiliábis; et quasi calóre sub nube torrénte, propáginem fórtium marcéscere fácies. Et fáciet Dóminus exercítuum ómnibus pópulis in monte hoc convívium píngujum, convívium vindémiæ, píngujum medullatórum, vindémiæ defæcátæ. Et præcipitábit in monte isto fáciem vínculi colligáti super omnes pópulos, et telam quam ordítus est super omnes natiónes.\par
\switchcolumn
\EN
\noindent For the blast of the mighty is like a whirlwind beating against a wall. Thou shalt bring down the tumult of strangers, as heat in thirst: and as with heat under a burning cloud, thou shalt make the branch of the mighty to wither away. And the Lord of hosts shall make unto all people in this mountain, a feast of fat things, a feast of wine, of fat things full of marrow, of wine purified from the lees. And he shall destroy in this mountain the face of the bond with which all people were tied, and the web that he began over all nations.\par
\switchcolumn*

\LA
\begin{Resp}\Rsign Festína, ne tardáveris, Dómine: \Star{} Et líbera pópulum tuum.\end{Resp}
\begin{Resp}\Vsign Veni, Dómine, et noli tardáre: reláxa facínora plebi tuæ.\end{Resp}
\begin{Resp}\Rsign Et líbera pópulum tuum.\end{Resp}
\switchcolumn
\EN
\begin{Resp}\Rsign Make haste, O Lord, make no tarrying. \Star{} And deliver thy people.\end{Resp}
\begin{Resp}\Vsign O Lord, come and make no tarrying loose the bonds of thy people.\end{Resp}
\begin{Resp}\Rsign And deliver thy people.\end{Resp}
\switchcolumn*

\LA
\LessonHead{Lectio III}
\switchcolumn
\EN
\LessonHead{Lesson III}
\switchcolumn*

\LA
\Cite{Isa 25:8-12}
\switchcolumn
\EN
\Cite{Isa 25:8-12}
\switchcolumn*

\LA
\noindent Præcipitábit mortem in sempitérnum; et áuferet Dóminus Deus lácrimam ab omni fácie, et oppróbrium pópuli sui áuferet de univérsa terra: quia Dóminus locútus est. Et dicet in die illa: Ecce Deus noster iste; exspectávimus eum, et salvábit nos; iste Dóminus, sustinúimus eum: exsultábimus, et lætábimur in salutári ejus. Quia requiéscet manus Dómini in monte isto; et triturábitur Moab sub eo, sícuti terúntur páleæ in plaustro. Et exténdet manus suas sub eo sicut exténdit natans ad natándum; et humiliábit glóriam ejus cum allisióne mánuum ejus. Et muniménta sublímium murórum tuórum cóncident, et humiliabúntur, et detrahéntur in terram usque ad púlverem.\par
\switchcolumn
\EN
\noindent He shall cast death down headlong forever: and the Lord God shall wipe away tears from every face, and the reproach of his people he shall take away from off the whole earth: for the Lord hath spoken it. And they shall say in that day: Lo, this is our God, we have waited for him, and he will save us: this is the Lord, we have patiently waited for him, we shall rejoice and be joyful in his salvation. For the hand of the Lord shall rest in this mountain: and Moab shall be trodden down under him, as straw is broken in pieces with the wain. And he shall stretch forth his hands under him, as he that swimmeth stretcheth forth his hands to swim: and he shall bring down his glory with the dashing of his hands. And the bulwarks of thy high walls shall fall, and be brought low, and shall be pulled down to the ground, even to the dust.\par
\switchcolumn*

\LA
\begin{Resp}\Rsign Ecce Dóminus véniet cum splendóre descéndens, et virtus ejus cum eo, \Star{} Visitáre pópulum suum in pace, et constitúere super eum vitam sempitérnam.\end{Resp}
\begin{Resp}\Vsign Ecce Dóminus noster cum virtúte véniet.\end{Resp}
\begin{Resp}\Rsign Visitáre pópulum suum in pace, et constitúere super eum vitam sempitérnam.\end{Resp}
\begin{Resp}\Vsign Glória Patri, et Fílio, \Star{} et Spirítui Sancto.\end{Resp}
\begin{Resp}\Rsign Visitáre pópulum suum in pace, et constitúere super eum vitam sempitérnam.\end{Resp}
\switchcolumn
\EN
\begin{Resp}\Rsign Behold, the Lord cometh down with glory, and His host is with Him. \Star{} To visit His people in peace, and to establish them in life everlasting.\end{Resp}
\begin{Resp}\Vsign Behold, our Lord cometh with an host.\end{Resp}
\begin{Resp}\Rsign To visit His people in peace, and to establish them in life everlasting.\end{Resp}
\begin{Resp}\Vsign Glory be to the Father, and to the Son, \Star{} and to the Holy Ghost.\end{Resp}
\begin{Resp}\Rsign To visit His people in peace, and to establish them in life everlasting.\end{Resp}
\switchcolumn*

\end{paracol}

\Office{Dominica III Adventus}{Third Sunday of Advent}{Semiduplex II classis}
\begin{paracol}{2}
\LA
\LessonHead{Lectio I}
\switchcolumn
\EN
\LessonHead{Lesson I}
\switchcolumn*

\LA
\LessonTitle{De Isaía Prophéta}
\switchcolumn
\EN
\LessonTitle{Lesson from the book of Isaias}
\switchcolumn*

\LA
\Cite{Isa 26:1-6}
\switchcolumn
\EN
\Cite{Isa 26:1-6}
\switchcolumn*

\LA
\noindent In die illa cantábitur cánticum istud in terra Juda: Urbs fortitúdinis nostræ Sion Salvátor, ponétur in ea murus et antemurále. Aperíte portas et ingrediátur gens justa, custódiens veritátem. Vetus error ábiit: servábis pacem: pacem, quia in te sperávimus. Sperástis in Dómino in sǽculis ætérnis, in Dómino Deo forti, in perpétuum. Quia incurvábit habitántes in excélso, civitátem sublímem humiliábit. Humiliábit eam usque ad terram, détrahet eam usque ad púlverem. Conculcábit eam pes, pedes páuperis, gressus egenórum.\par
\switchcolumn
\EN
\noindent In that day shall this canticle be sung in the land of Juda. Sion the city of our strength a saviour, a wall and a bulwark shall be set therein. Open ye the gates, and let the just nation, that keepeth the truth, enter in. The old error is passed away: thou wilt keep peace: peace, because we have hoped in thee. You have hoped in the Lord for evermore, in the Lord God mighty for ever. For he shall bring down them that dwell on high, the high city he shall lay low. He shall bring it down even to the ground, he shall pull it down even to the dust. The foot shall tread it down, the feet of the poor, the steps of the needy.\par
\switchcolumn*

\LA
\begin{Resp}\Rsign Ecce apparébit Dóminus super nubem cándidam, \Star{} Et cum eo Sanctórum míllia: et habébit in vestiménto, et in fémore suo scriptum: Rex regum, et Dóminus dominántium.\end{Resp}
\begin{Resp}\Vsign Apparébit in finem, et non mentiétur; si moram fécerit, exspécta eum, quia véniens véniet.\end{Resp}
\begin{Resp}\Rsign Et cum eo Sanctórum míllia: et habébit in vestiménto, et in fémore suo scriptum: Rex regum, et Dóminus dominántium.\end{Resp}
\switchcolumn
\EN
\begin{Resp}\Rsign Behold, the Lord shall appear upon a white cloud, \Star{} And ten thousand of His saints with Him; and He shall have on His vesture, and on His thigh a name written King of kings, and Lord of lords.\end{Resp}
\begin{Resp}\Vsign He shall appear and not lie; though He tarry, wait for Him, because He will surely come.\end{Resp}
\begin{Resp}\Rsign And ten thousand of His saints with Him; and He shall have on His vesture, and on His thigh a name written King of kings, and Lord of lords.\end{Resp}
\switchcolumn*

\LA
\LessonHead{Lectio II}
\switchcolumn
\EN
\LessonHead{Lesson II}
\switchcolumn*

\LA
\Cite{Isa 26:7-10}
\switchcolumn
\EN
\Cite{Isa 26:7-10}
\switchcolumn*

\LA
\noindent Sémita justi recta est, rectus callis justi ad ambulándum. Et in sémita judiciórum tuórum, Dómine, sustinúimus te: nomen tuum, et memoriále tuum in desidério ánimæ. Anima mea desiderávit te in nocte, sed et spíritu meo in præcórdiis meis de mane vigilábo ad te. Cum féceris judícia tua in terra, justítiam discent habitatóres orbis. Misereámur ímpio, et non discet justítiam: in terra sanctórum iníqua gessit, et non vidébit glóriam Dómini.\par
\switchcolumn
\EN
\noindent The way of the just is right, the path of the just is right to walk in. And in the way of thy judgments, O Lord, we have patiently waited for thee: thy name, and thy remembrance are the desire of the soul. My soul hath desired thee in the night: yea, and with my spirit within me in the morning early I will watch to thee. When thou shalt do thy judgments on the earth, the inhabitants of the world shall learn justice. Let us have pity on the wicked, but he will not learn justice: in the land of the saints he hath done wicked things, and he shall not see the glory of the Lord.\par
\switchcolumn*

\LA
\begin{Resp}\Rsign Béthlehem cívitas Dei summi, ex te éxiet Dominátor Israël, et egréssus ejus sicut a princípio diérum æternitátis, et magnificábitur in médio univérsæ terræ: \Star{} Et pax erit in terra nostra, dum vénerit.\end{Resp}
\begin{Resp}\Vsign Loquétur pacem in géntibus, et potéstas ejus a mari usque ad mare.\end{Resp}
\begin{Resp}\Rsign Et pax erit in terra nostra, dum vénerit.\end{Resp}
\switchcolumn
\EN
\begin{Resp}\Rsign Thou, Bethlehem, art the city of the Most High God, out of thee shall He come forth That is to be Ruler in Israel; Whose goings forth have been from of old, from everlasting, and now shall He be great unto the ends of the earth. \Star{} And this Man shall be the peace in our land, when He shall come.\end{Resp}
\begin{Resp}\Vsign He shall speak peace unto the Gentiles, and shall have dominion from sea to sea.\end{Resp}
\begin{Resp}\Rsign And this Man shall be the peace in our land, when He shall come.\end{Resp}
\switchcolumn*

\LA
\LessonHead{Lectio III}
\switchcolumn
\EN
\LessonHead{Lesson III}
\switchcolumn*

\LA
\Cite{Isa 26:11-14}
\switchcolumn
\EN
\Cite{Isa 26:11-14}
\switchcolumn*

\LA
\noindent Dómine, exaltétur manus tua, et non vídeant: vídeant, et confundántur zelántes pópuli: et ignis hostes tuos dévoret. Dómine, dabis pacem nobis: ómnia enim ópera nostra operátus es nobis. Dómine Deus noster, possedérunt nos dómini absque te, tantum in te recordémur nóminis tui. Moriéntes non vivant, gigántes non resúrgant: proptérea visitásti et contrivísti eos, et perdidísti omnem memóriam eórum.\par
\switchcolumn
\EN
\noindent Lord, let thy hand be exalted, and let them not see: let the envious people see, and be confounded: and let fire devour thy enemies. Lord, thou wilt give us peace: for thou hast wrought all our works for us. O Lord our God, other lords besides thee have had dominion over us, only in thee let us remember thy name. Let not the dead live, let not the giants rise again: therefore hast thou visited and destroyed them, and hast destroyed all their memory.\par
\switchcolumn*

\LA
\begin{Resp}\Rsign Qui ventúrus est, véniet, et non tardábit: et jam non erit timor in fínibus nostris: \Star{} Quóniam ipse est Salvátor noster.\end{Resp}
\begin{Resp}\Vsign Depónet omnes iniquitátes nostras, et proíciet in profúndum maris ómnia peccáta nostra.\end{Resp}
\begin{Resp}\Rsign Quóniam ipse est Salvátor noster.\end{Resp}
\begin{Resp}\Vsign Glória Patri, et Fílio, \Star{} et Spirítui Sancto.\end{Resp}
\begin{Resp}\Rsign Quóniam ipse est Salvátor noster.\end{Resp}
\switchcolumn
\EN
\begin{Resp}\Rsign He That shall come, will come, and will not tarry; and there shall no more be fear in our borders. \Star{} For He is our Saviour.\end{Resp}
\begin{Resp}\Vsign He shall tread down all our iniquities, and cast all our sins into the depths of the sea.\end{Resp}
\begin{Resp}\Rsign For He is our Saviour.\end{Resp}
\begin{Resp}\Vsign Glory be to the Father, and to the Son, \Star{} and to the Holy Ghost.\end{Resp}
\begin{Resp}\Rsign For He is our Saviour.\end{Resp}
\switchcolumn*

\LA
\LessonHead{Lectio IV}
\switchcolumn
\EN
\LessonHead{Lesson IV}
\switchcolumn*

\LA
\LessonTitle{Sermo sancti Leónis Papæ}
\switchcolumn
\EN
\LessonTitle{From the Sermons of Pope St. Leo (the Great)}
\switchcolumn*

\LA
\Source{Sermo 2 de jejúnio décimi mensis et collectis}
\switchcolumn
\EN
\Source{Second on the December Fast, and Almsgiving}
\switchcolumn*

\LA
\noindent Quod témporis rátio, et devotiónis nostræ ádmonet consuetúdo, pastoráli vobis, dilectíssimi, sollicitúdine prædicámus, décimi mensis celebrándum esse jejúnium, quo pro consummáta perceptióne ómnium frúctuum, digníssime largitóri eórum Deo continéntiæ libámen offértur. Quid enim potest efficácius esse jejúnio? cujus observántia appropinquámus Deo, et resisténtes diábolo, vítia blanda superámus.\par
\switchcolumn
\EN
\noindent Dearly beloved brethren, with the care which becometh us as the shepherd of your souls, we urge upon you the rigid observance of this December Fast. The month of December hath come round again, and with it this devout custom of the Church. The fruits of the year, which is drawing to a close, are now all gathered in, and we most meetly offer our abstinence to God as a sacrifice of thanksgiving. And what can be more useful than fasting, that exercise by which we draw nigh to God, make a stand against the devil, and overcome the softer enticements of sin?\par
\switchcolumn*

\LA
\begin{Resp}\Rsign Ægýpte, noli flere, quia Dominátor tuus véniet tibi, ante cujus conspéctum movebúntur abýssi, \Star{} Liberáre pópulum suum de manu poténtiæ.\end{Resp}
\begin{Resp}\Vsign Ecce véniet Dóminus exercítuum, Deus tuus cum potestáte magna.\end{Resp}
\begin{Resp}\Rsign Liberáre pópulum suum de manu poténtiæ.\end{Resp}
\switchcolumn
\EN
\begin{Resp}\Rsign Weep not, O Egypt, for the Ruler cometh unto thee, and the depths shall be moved at His presence. \Star{} To deliver His people out of the hand of the mighty.\end{Resp}
\begin{Resp}\Vsign Behold, the Lord of hosts, thy God, cometh with great power.\end{Resp}
\begin{Resp}\Rsign To deliver His people out of the hand of the mighty.\end{Resp}
\switchcolumn*

\LA
\LessonHead{Lectio V}
\switchcolumn
\EN
\LessonHead{Lesson V}
\switchcolumn*

\LA
\noindent Semper enim virtúti cibus jejúnium fuit. De abstinéntia dénique pródeunt castæ cogitatiónes, rationábiles voluntátes, salubrióra consília: et per voluntárias afflictiónes caro concupiscéntiis móritur, virtútibus spíritus innovátur. Sed quia non solo jejúnio animárum nostrárum salus acquíritur, jejúnium nostrum misericórdiis páuperum suppleámus. Impendámus virtúti, quod subtráhimus voluptáti. Fiat reféctio páuperis abstinéntia jejunántis.\par
\switchcolumn
\EN
\noindent Fasting hath ever been the bread of strength. From abstinence proceed pure thoughts, reasonable desires, and healthy counsels. By voluntary mortifications the flesh dieth to lust, and the soul is renewed in might. But since fasting is not the only mean whereby we get health for our souls, let us add to our fasting works of mercy. Let us spend in good deeds what we take from indulgence. Let our fast become the banquet of the poor.\par
\switchcolumn*

\LA
\begin{Resp}\Rsign Prope est ut véniat tempus ejus, et dies ejus non elongabúntur: \Star{} Miserébitur Dóminus Jacob, et Israël salvábitur.\end{Resp}
\begin{Resp}\Vsign Revértere, virgo Israël, revértere ad civitátes tuas.\end{Resp}
\begin{Resp}\Rsign Miserébitur Dóminus Jacob, et Israël salvábitur.\end{Resp}
\switchcolumn
\EN
\begin{Resp}\Rsign Her time is near to come, and her days shall not be prolonged. \Star{} Lord shall have mercy on Jacob, and Israel shall be saved.\end{Resp}
\begin{Resp}\Vsign Turn again, O Virgin of Israel, turn again to thy cities.\end{Resp}
\begin{Resp}\Rsign For the Lord shall have mercy on Judah, and Israel shall be saved.\end{Resp}
\switchcolumn*

\LA
\LessonHead{Lectio VI}
\switchcolumn
\EN
\LessonHead{Lesson VI}
\switchcolumn*

\LA
\noindent Studeámus viduárum defensióni, pupillórum utilitáti, lugéntium consolatióni, dissidéntium paci. Suscipiátur peregrínus, adjuvétur oppréssus, vestiátur nudus, foveátur ægrótus: ut quicúmque nostrum de justis labóribus auctóri bonórum ómnium Deo sacrifícium hujus pietátis obtúlerit, ab eodem regni cæléstis prǽmium percípere mereátur. Quarta ígitur et sexta Feria jejunémus; Sábbato autem apud beátum Petrum Apóstolum páriter vigilémus: cujus suffragántibus méritis, quæ póscimus, impetráre possímus per Dóminum nostrum Jesum Christum, qui cum Patre et Sancto Spíritu vivit et regnat in sǽcula sæculórum. Amen.\par
\switchcolumn
\EN
\noindent Let us defend the widow and serve the orphan; let us comfort the afflicted and reconcile the estranged; let us take in the wanderer and succour the oppressed; let us clothe the naked and cherish the sick. And may every one of us that shall offer to the God of all goodness this Advent sacrifice of fasting and alms be by Him fitted to receive an eternal reward in His heavenly kingdom! We fast on Wednesday and Friday; and there is likewise a Vigil on Saturday at the Church of St. Peter, that by his good prayers we may the more effectually obtain what we ask for, through our Lord Jesus Christ, Who with the Father and the Holy Ghost, liveth and reigneth, God, world without end. Amen.\par
\switchcolumn*

\LA
\begin{Resp}\Rsign Descéndet Dóminus sicut plúvia in vellus: \Star{} Oriétur in diébus ejus justítia, et abundántia pacis.\end{Resp}
\begin{Resp}\Vsign Et adorábunt eum omnes reges, omnes gentes sérvient ei.\end{Resp}
\begin{Resp}\Rsign Oriétur in diébus ejus justítia, et abundántia pacis.\end{Resp}
\begin{Resp}\Vsign Glória Patri, et Fílio, \Star{} et Spirítui Sancto.\end{Resp}
\begin{Resp}\Rsign Oriétur in diébus ejus justítia, et abundántia pacis.\end{Resp}
\switchcolumn
\EN
\begin{Resp}\Rsign The Lord shall come down like rain upon a fleece. \Star{} In His days shall righteousness flourish, and abundance of peace.\end{Resp}
\begin{Resp}\Vsign All the kings of the earth shall fall down before Him, all nations shall serve Him.\end{Resp}
\begin{Resp}\Rsign In His days shall righteousness flourish, and abundance of peace.\end{Resp}
\begin{Resp}\Vsign Glory be to the Father, and to the Son, \Star{} and to the Holy Ghost.\end{Resp}
\begin{Resp}\Rsign In His days shall righteousness flourish, and abundance of peace.\end{Resp}
\switchcolumn*

\LA
\LessonHead{Lectio VII}
\switchcolumn
\EN
\LessonHead{Lesson VII}
\switchcolumn*

\LA
\LessonTitle{Léctio sancti Evangélii secúndum Joánnem}
\switchcolumn
\EN
\LessonTitle{From the Holy Gospel according to John}
\switchcolumn*

\LA
\Cite{Joannes 1:19-28}
\switchcolumn
\EN
\Cite{John 1:19-28}
\switchcolumn*

\LA
\noindent In illo témpore: Misérunt Judǽi ab Jerosólymis sacerdótes et levítas ad Joánnem, ut interrogárent eum: Tu quis es? Et réliqua.\par
\switchcolumn
\EN
\noindent In that time were sent from Jerusalem priests and Levites to John, to ask him: Who art thou? And so on.\par
\switchcolumn*

\LA
\LessonTitle{Homilía sancti Gregórii Papæ}
\switchcolumn
\EN
\LessonTitle{Homily by Pope St. Gregory (the Great)}
\switchcolumn*

\LA
\Source{Homilia 7 in Evangelia}
\switchcolumn
\EN
\Source{7th on the Gospels}
\switchcolumn*

\LA
\noindent Ex hujus nobis lectiónis verbis, fratres caríssimi, Joánnis humílitas commendátur: qui cum tantæ virtútis esset, ut Christus credi potuísset, elégit sólide subsístere in se, ne humána opinióne raperétur inániter super se. Nam conféssus est, et non negávit: et conféssus est, Quia non sum ego Christus. Sed qui dixit, Non sum; negávit plane quod non erat, sed non negávit quod erat: ut veritátem loquens, ejus membrum fíeret, cujus sibi nomen falláciter non usurpáret. Cum ergo non vult appétere nomen Christi, factus est membrum Christi: quia dum infirmitátem suam stúduit humíliter agnóscere, illíus celsitúdinem méruit veráciter obtinére.\par
\switchcolumn
\EN
\noindent Dearly beloved brethren, the first thing which striketh us in today's Gospel is the lowly-mindedness of John. He was so great that it was thought he might be the Christ; yet he soberly chose rather to seem only what he really was, than to let the belief of men invest him with a dignity which did not belong to him; for he confessed, and denied not, but confessed, I am not the Christ, at the same time he would not deny what he was in reality; and thus his very truth-speaking made him a member of Him Whose title he would not by falsehood take. In that he arrogated not to himself the name of Christ, he became a member of Christ. While he humbly strove to confess his own weakness, he earned by his simplicity a part in the grandeur of his Master.\par
\switchcolumn*

\LA
\begin{Resp}\Rsign Veni, Dómine, et noli tardáre: reláxa facínora plebi tuæ, \Star{} Et révoca dispérsos in terram suam.\end{Resp}
\begin{Resp}\Vsign Excita, Dómine, poténtiam tuam, et veni, ut salvos fácias nos.\end{Resp}
\begin{Resp}\Rsign Et révoca dispérsos in terram suam.\end{Resp}
\switchcolumn
\EN
\begin{Resp}\Rsign O Lord, come and make no tarrying; loosen the bonds of thy people. \Star{} And gather again into their own land them that are scattered abroad.\end{Resp}
\begin{Resp}\Vsign O Lord, stir up thy strength, and come and save us.\end{Resp}
\begin{Resp}\Rsign And gather again into their own land them that are scattered abroad.\end{Resp}
\switchcolumn*

\LA
\LessonHead{Lectio VIII}
\switchcolumn
\EN
\LessonHead{Lesson VIII}
\switchcolumn*

\LA
\noindent Sed cum ex lectióne alia, Redemptóris nostri senténtia ad mentem redúcitur, ex hujus lectiónis verbis nobis quǽstio valde impléxa generátur. Alio quippe in loco inquisítus a discípulis Dóminus de Elíæ advéntu, respóndit: Elías jam venit, et non cognovérunt eum, sed fecérunt in eum quæcúmque voluérunt: et, si vultis scire, Joánnes ipse est Elías. Requisítus autem Joánnes dicit: Non sum Elías. Quid est hoc, fratres caríssimi, quia quod Véritas affírmat, hoc prophéta Veritátis negat? Valde namque inter se divérsa sunt: Ipse est: et, Non sum. Quómodo ergo prophéta veritátis est, si ejusdem Veritátis sermónibus concors non est?\par
\switchcolumn
\EN
\noindent In considering this subject we find an apparent contradiction between one of John's statements, and the saying of our Redeemer recorded in another part of the Gospel. (Matth. xvii. 10-12.) When His disciples asked our Lord regarding the coming of Elias, He answered: Elias is come already, and they knew him not, but have done unto him whatsoever they listed. And if ye will receive it, this (that is, John) is Elias. (Matth. xi. 14.) But when John was asked if he was Elias, he answered: I am not. How comes it then, dearly beloved brethren, that we find the Truth Itself asserting what the prophet of the Truth denied? It must evidently be that our Lord meant one thing and John another, when the Lord said, This is, and John, I am not. For how can he be the prophet of truth, if he speak not according to the word of Him Who is the Eternal Truth?\par
\switchcolumn*

\LA
\begin{Resp}\Rsign Ecce radix Jesse descéndet in salútem populórum, ipsum gentes deprecabúntur: \Star{} Et erit nomen ejus gloriósum.\end{Resp}
\begin{Resp}\Vsign Dabit ei Dóminus Deus sedem David, patris ejus, et regnábit in domo Jacob in ætérnum.\end{Resp}
\begin{Resp}\Rsign Et erit nomen ejus gloriósum.\end{Resp}
\switchcolumn
\EN
\begin{Resp}\Rsign Behold, there shall be a root of Jesse, which shall come for salvation unto the people, to it shall the Gentiles seek; \Star{} And His name shall be glorious.\end{Resp}
\begin{Resp}\Vsign The Lord God shall give unto Him the throne of His father David, and He shall reign over the house of Jacob for ever.\end{Resp}
\begin{Resp}\Rsign And His name shall be glorious.\end{Resp}
\switchcolumn*

\LA
\LessonHead{Lectio IX}
\switchcolumn
\EN
\LessonHead{Lesson IX}
\switchcolumn*

\LA
\noindent Sed si subtíliter véritas ipsa requirátur, hoc quod inter se contrárium sonat, quómodo contrárium non sit, invenítur. Ad Zacharíam namque de Joánne Angelus dicit: Ipse præcédet ante illum in spíritu et virtúte Elíæ. Qui idcírco ventúrus in spíritu et virtúte Elíæ dícitur, quia sicut Elías secúndum Dómini advéntum prævéniet, ita Joánnes prævénit primum. Sicut ille præcúrsor ventúrus est Júdicis, ita iste præcúrsor est factus Redemptóris. Joánnes ígitur in spíritu Elías erat, in persóna Elías non erat. Quod ergo Dóminus fatétur de spíritu, hoc Joánnes dénegat de persóna.\par
\switchcolumn
\EN
\noindent Let us then more minutely examine these words, and we shall find that there is no real contradiction. When the Angel announced to Zacharias the coming birth of John he said He shall go before Him in the spirit and power of Elias, (Luke i. 17.) As the old Elias will come again before the Second Advent of the Lord, so did John, as the new Elias, go before the First Advent, in the spirit and power of Elias. As the old Elias will be the Fore-runner of the Judge, so the new Elias was the Forerunner of the Saviour. John then was Elias in spirit, but not in person; and our Lord asserteth of the spirit what John denieth of the person.\par
\switchcolumn*

\LA
\begin{Resp}\Rsign Docébit nos Dóminus vias suas, et ambulábimus in sémitis ejus: \Star{} Quia de Sion exíbit lex, et verbum Dómini de Jerúsalem.\end{Resp}
\begin{Resp}\Vsign Veníte, ascendámus ad montem Dómini, et ad domum Dei Jacob.\end{Resp}
\begin{Resp}\Rsign Quia de Sion exíbit lex, et verbum Dómini de Jerúsalem.\end{Resp}
\begin{Resp}\Vsign Glória Patri, et Fílio, \Star{} et Spirítui Sancto.\end{Resp}
\begin{Resp}\Rsign Quia de Sion exíbit lex, et verbum Dómini de Jerúsalem.\end{Resp}
\switchcolumn
\EN
\begin{Resp}\Rsign The Lord will teach us of His ways, and we will walk in His paths. \Star{} For out of Zion shall go forth the law, and the word of the Lord from Jerusalem.\end{Resp}
\begin{Resp}\Vsign Come ye, and let us go up to the mountain of the Lord, and to the house of the God of Jacob.\end{Resp}
\begin{Resp}\Rsign For out of Zion shall go forth the law, and the word of the Lord from Jerusalem.\end{Resp}
\begin{Resp}\Vsign Glory be to the Father, and to the Son, \Star{} and to the Holy Ghost.\end{Resp}
\begin{Resp}\Rsign For out of Zion shall go forth the law, and the word of the Lord from Jerusalem.\end{Resp}
\switchcolumn*

\end{paracol}

\Office{Feria II infra Hebdomadam III Adventus}{Monday of the Third Week of Advent}{Feria major}
\begin{paracol}{2}
\LA
\LessonHead{Lectio I}
\switchcolumn
\EN
\LessonHead{Lesson I}
\switchcolumn*

\LA
\LessonTitle{De Isaía Prophéta}
\switchcolumn
\EN
\LessonTitle{Lesson from the book of Isaias}
\switchcolumn*

\LA
\Cite{Isa 28:1-3}
\switchcolumn
\EN
\Cite{Isa 28:1-3}
\switchcolumn*

\LA
\noindent Væ corónæ supérbiæ, ébriis Ephraim, et flori decidénti, glóriæ exsultatiónis ejus, qui erant in vértice vallis pinguíssimæ, errántes a vino. Ecce válidus et fortis Dóminus, sicut ímpetus grándinis: turbo confríngens, sicut ímpetus aquárum multárum inundántium, et emissárum super terram spatiósam. Pédibus conculcábitur coróna supérbiæ ebriórum Ephraim.\par
\switchcolumn
\EN
\noindent Woe to the crown of pride, to the drunkards of Ephraim, and to the fading flower the glory of his joy, who were on the head of the fat valley, staggering with wine. Behold the Lord is mighty and strong, as a storm of hail: a destroying whirlwind, as the violence of many waters overflowing, and sent forth upon a spacious land. The crown of pride of the drunkards of Ephraim shall be trodden under feet.\par
\switchcolumn*

\LA
\begin{Resp}\Rsign Ecce apparébit Dóminus super nubem cándidam, \Star{} Et cum eo Sanctórum míllia: et habébit in vestiménto, et in fémore suo scriptum: Rex regum, et Dóminus dominántium.\end{Resp}
\begin{Resp}\Vsign Apparébit in finem, et non mentiétur; si moram fécerit, exspécta eum, quia véniens véniet.\end{Resp}
\begin{Resp}\Rsign Et cum eo Sanctórum míllia: et habébit in vestiménto, et in fémore suo scriptum: Rex regum, et Dóminus dominántium.\end{Resp}
\switchcolumn
\EN
\begin{Resp}\Rsign Behold, the Lord shall appear upon a white cloud, \Star{} And ten thousand of His saints with Him; and He shall have on His vesture, and on His thigh a name written King of kings, and Lord of lords.\end{Resp}
\begin{Resp}\Vsign He shall appear and not lie; though He tarry, wait for Him, because He will surely come.\end{Resp}
\begin{Resp}\Rsign And ten thousand of His saints with Him; and He shall have on His vesture, and on His thigh a name written King of kings, and Lord of lords.\end{Resp}
\switchcolumn*

\LA
\LessonHead{Lectio II}
\switchcolumn
\EN
\LessonHead{Lesson II}
\switchcolumn*

\LA
\Cite{Isa 28:4-7}
\switchcolumn
\EN
\Cite{Isa 28:4-7}
\switchcolumn*

\LA
\noindent Et erit flos décidens glóriæ exsultatiónis ejus, qui est super vérticem vallis píngujum, quasi temporáneum ante maturitátem autúmni: quod cum aspéxerit videns, statim ut manu tenúerit, devorábit illud. In die illa erit Dóminus exercítuum coróna glóriæ, et sertum exsultatiónis resíduo pópuli sui; et spíritus judícii sedénti super judícium, et fortitúdo reverténtibus de bello ad portam. Verum hi quoque præ vino nesciérunt, et præ ebrietáte erravérunt: sacérdos et prophéta nesciérunt præ ebrietáte, absórpti sunt a vino.\par
\switchcolumn
\EN
\noindent And the fading flower the glory of his joy, who is on the head of the fat valley, shall be as a hasty fruit before the ripeness of autumn: which when he that seeth it shall behold, as soon as he taketh it in his hand, he will eat it up. In that day the Lord of hosts shall be a crown of glory, and a garland of joy to the residue of his people: And a spirit of judgment to him that sitteth in judgment, and strength to them that return out of the battle to the gate. But these also have been ignorant through wine, and through drunkenness have erred: the priest and the prophet have been ignorant through drunkenness, they are swallowed up with wine, they have gone astray in drunkenness, they have not known him that seeth, they have been ignorant of judgment.\par
\switchcolumn*

\LA
\begin{Resp}\Rsign Béthlehem cívitas Dei summi, ex te éxiet Dominátor Israël, et egréssus ejus sicut a princípio diérum æternitátis, et magnificábitur in médio univérsæ terræ: \Star{} Et pax erit in terra nostra, dum vénerit.\end{Resp}
\begin{Resp}\Vsign Loquétur pacem in géntibus, et potéstas ejus a mari usque ad mare.\end{Resp}
\begin{Resp}\Rsign Et pax erit in terra nostra, dum vénerit.\end{Resp}
\switchcolumn
\EN
\begin{Resp}\Rsign Thou, Bethlehem, art the city of the Most High God, out of thee shall He come forth That is to be Ruler in Israel; Whose goings forth have been from of old, from everlasting, and now shall He be great unto the ends of the earth. \Star{} And this Man shall be the peace in our land, when He shall come.\end{Resp}
\begin{Resp}\Vsign He shall speak peace unto the Gentiles, and shall have dominion from sea to sea.\end{Resp}
\begin{Resp}\Rsign And this Man shall be the peace in our land, when He shall come.\end{Resp}
\switchcolumn*

\LA
\LessonHead{Lectio III}
\switchcolumn
\EN
\LessonHead{Lesson III}
\switchcolumn*

\LA
\Cite{Isa 28:16-18}
\switchcolumn
\EN
\Cite{Isa 28:16-18}
\switchcolumn*

\LA
\noindent Idcírco hæc dicit Dóminus Deus: Ecce ego mittam in fundaméntis Sion lápidem, lápidem probátum, angulárem, pretiósum, in fundaménto fundátum. Qui credíderit, non festínet. Et ponam in póndere judícium, et justítiam in mensúra: et subvértet grando spem mendácii, et protectiónem aquæ inundábunt. Et delébitur fœdus vestrum cum morte, et pactum vestrum cum inférno non stabit.\par
\switchcolumn
\EN
\noindent Therefore thus saith the Lord God: Behold I will lay a stone in the foundations of Sion, a tried stone, a corner stone, a precious stone, founded in the foundation. He that believeth, let him not hasten. And I will set judgment in weight, and justice in measure: and hail shall overturn the hope of falsehood: and waters shall overflow its protection. And your league with death shall be abolished, and your covenant with hell shall not stand: when the overflowing scourge shall pass, you shall be trodden down by it.\par
\switchcolumn*

\LA
\begin{Resp}\Rsign Qui ventúrus est, véniet, et non tardábit: et jam non erit timor in fínibus nostris: \Star{} Quóniam ipse est Salvátor noster.\end{Resp}
\begin{Resp}\Vsign Depónet omnes iniquitátes nostras, et proíciet in profúndum maris ómnia peccáta nostra.\end{Resp}
\begin{Resp}\Rsign Quóniam ipse est Salvátor noster.\end{Resp}
\begin{Resp}\Vsign Glória Patri, et Fílio, \Star{} et Spirítui Sancto.\end{Resp}
\begin{Resp}\Rsign Quóniam ipse est Salvátor noster.\end{Resp}
\switchcolumn
\EN
\begin{Resp}\Rsign He That shall come, will come, and will not tarry; and there shall no more be fear in our borders. \Star{} For He is our Saviour.\end{Resp}
\begin{Resp}\Vsign He shall tread down all our iniquities, and cast all our sins into the depths of the sea.\end{Resp}
\begin{Resp}\Rsign For He is our Saviour.\end{Resp}
\begin{Resp}\Vsign Glory be to the Father, and to the Son, \Star{} and to the Holy Ghost.\end{Resp}
\begin{Resp}\Rsign For He is our Saviour.\end{Resp}
\switchcolumn*

\end{paracol}

\Office{Feria III infra Hebdomadam III Adventus}{Tuesday of the Third Week of Advent}{Feria major}
\begin{paracol}{2}
\LA
\LessonHead{Lectio I}
\switchcolumn
\EN
\LessonHead{Lesson I}
\switchcolumn*

\LA
\LessonTitle{De Isaía Prophéta}
\switchcolumn
\EN
\LessonTitle{Lesson from the book of Isaias}
\switchcolumn*

\LA
\Cite{Isa 30:18-20}
\switchcolumn
\EN
\Cite{Isa 30:18-20}
\switchcolumn*

\LA
\noindent Exspéctat Dóminus ut misereátur vestri, et ídeo exaltábitur parcens vobis: quia Deus judícii Dóminus: beáti omnes qui exspéctant eum. Pópulus enim Sion habitábit in Jerúsalem: plorans nequáquam plorábis, míserans miserébitur tui: ad vocem clamóris tui statim ut audíerit, respondébit tibi. Et dabit vobis Dóminus panem arctum, et aquam brevem: et non fáciet avoláre a te ultra doctórem tuum: et erunt óculi tui vidéntes præceptórem tuum.\par
\switchcolumn
\EN
\noindent The Lord waiteth that he may have mercy on you: and therefore shall he be exalted sparing you: because the Lord is the God of judgment: blessed are all they that wait for him. For the people of Sion shall dwell in Jerusalem: weeping thou shalt not weep, he will surely have pity on thee: at the voice of thy cry, as soon as he shall hear, he will answer thee. And the Lord will give you spare bread, and short water: and will not cause thy teacher to flee away from thee any more, and thy eyes shall see thy teacher.\par
\switchcolumn*

\LA
\begin{Resp}\Rsign Ægýpte, noli flere, quia Dominátor tuus véniet tibi, ante cujus conspéctum movebúntur abýssi, \Star{} Liberáre pópulum suum de manu poténtiæ.\end{Resp}
\begin{Resp}\Vsign Ecce véniet Dóminus exercítuum, Deus tuus cum potestáte magna.\end{Resp}
\begin{Resp}\Rsign Liberáre pópulum suum de manu poténtiæ.\end{Resp}
\switchcolumn
\EN
\begin{Resp}\Rsign Weep not, O Egypt, for the Ruler cometh unto thee, and the depths shall be moved at His presence. \Star{} To deliver His people out of the hand of the mighty.\end{Resp}
\begin{Resp}\Vsign Behold, the Lord of hosts, thy God, cometh with great power.\end{Resp}
\begin{Resp}\Rsign To deliver His people out of the hand of the mighty.\end{Resp}
\switchcolumn*

\LA
\LessonHead{Lectio II}
\switchcolumn
\EN
\LessonHead{Lesson II}
\switchcolumn*

\LA
\Cite{Isa 30:22-25}
\switchcolumn
\EN
\Cite{Isa 30:22-25}
\switchcolumn*

\LA
\noindent Egrédere, dices ei: Et dábitur plúvia sémini tuo, ubicúmque semináveris in terra: et panis frugum terræ erit ubérrimus et pinguis. Pascétur in possessióne tua in die illo agnus spatióse: et tauri tui, et pulli asinórum, qui operántur terram, commístum migma cómedent sicut in área ventilátum est. Et erunt super omnem montem excélsum, et super omnem collem elevátum, rivi curréntium aquárum, in die interfectiónis multórum, cum cecíderint turres.\par
\switchcolumn
\EN
\noindent And thou shalt defile the plates of thy graven things of silver, and the garment of thy molten things of gold, and shalt cast them away as the uncleanness of a menstruous woman. Thou shalt say to it: Get thee hence. And rain shall be given to thy seed, wheresoever thou shalt sow in the land: and the bread of the corn of the land shall be most plentiful, and fat. The lamb in that day shall feed at large in thy possession: And thy oxen, and the ass colts that till the ground, shall eat mingled provender as it was winnowed in the floor. And there shall be upon every high mountain, and upon every elevated hill rivers of running waters in the day of the slaughter of many, when the tower shall fall.\par
\switchcolumn*

\LA
\begin{Resp}\Rsign Prope est ut véniat tempus ejus, et dies ejus non elongabúntur: \Star{} Miserébitur Dóminus Jacob, et Israël salvábitur.\end{Resp}
\begin{Resp}\Vsign Revértere, virgo Israël, revértere ad civitátes tuas.\end{Resp}
\begin{Resp}\Rsign Miserébitur Dóminus Jacob, et Israël salvábitur.\end{Resp}
\switchcolumn
\EN
\begin{Resp}\Rsign Her time is near to come, and her days shall not be prolonged. \Star{} Lord shall have mercy on Jacob, and Israel shall be saved.\end{Resp}
\begin{Resp}\Vsign Turn again, O Virgin of Israel, turn again to thy cities.\end{Resp}
\begin{Resp}\Rsign For the Lord shall have mercy on Judah, and Israel shall be saved.\end{Resp}
\switchcolumn*

\LA
\LessonHead{Lectio III}
\switchcolumn
\EN
\LessonHead{Lesson III}
\switchcolumn*

\LA
\Cite{Isa 30:26-28}
\switchcolumn
\EN
\Cite{Isa 30:26-28}
\switchcolumn*

\LA
\noindent Et erit lux lunæ sicut lux solis, et lux solis erit septemplíciter sicut lux septem diérum in die, qua alligáverit Dóminus vulnus pópuli sui, et percussúram plagæ ejus sanáverit. Ecce nomen Dómini venit de longínquo, ardens furor ejus, et gravis ad portándum: lábia ejus repléta sunt indignatióne, et lingua ejus quasi ignis dévorans. Spíritus ejus velut torrens inúndans usque ad médium colli, ad perdéndas gentes in níhilum, et frenum erróris, quod erat in maxíllis populórum.\par
\switchcolumn
\EN
\noindent And the light of the moon shall be as the light of the sun, and the light of the sun shall be sevenfold, as the light of seven days: in the day when the Lord shall bind up the wound of his people, and shall heal the stroke of their wound. Behold the name of the Lord cometh from afar, his wrath burneth, and is heavy to bear: his lips are filled with indignation, and his tongue as a devouring fire. His breath as a torrent overflowing even to the midst of the neck, to destroy the nations unto nothing, and the bridle of error that was in the jaws of the people.\par
\switchcolumn*

\LA
\begin{Resp}\Rsign Descéndet Dóminus sicut plúvia in vellus: \Star{} Oriétur in diébus ejus justítia, et abundántia pacis.\end{Resp}
\begin{Resp}\Vsign Et adorábunt eum omnes reges, omnes gentes sérvient ei.\end{Resp}
\begin{Resp}\Rsign Oriétur in diébus ejus justítia, et abundántia pacis.\end{Resp}
\begin{Resp}\Vsign Glória Patri, et Fílio, \Star{} et Spirítui Sancto.\end{Resp}
\begin{Resp}\Rsign Oriétur in diébus ejus justítia, et abundántia pacis.\end{Resp}
\switchcolumn
\EN
\begin{Resp}\Rsign The Lord shall come down like rain upon a fleece. \Star{} In His days shall righteousness flourish, and abundance of peace.\end{Resp}
\begin{Resp}\Vsign All the kings of the earth shall fall down before Him, all nations shall serve Him.\end{Resp}
\begin{Resp}\Rsign In His days shall righteousness flourish, and abundance of peace.\end{Resp}
\begin{Resp}\Vsign Glory be to the Father, and to the Son, \Star{} and to the Holy Ghost.\end{Resp}
\begin{Resp}\Rsign In His days shall righteousness flourish, and abundance of peace.\end{Resp}
\switchcolumn*

\end{paracol}

\Office{Feria IV Quattuor Temporum in Adventu}{Ember Wednesday of Advent}{Feria major}
\begin{paracol}{2}
\LA
\LessonHead{Lectio I}
\switchcolumn
\EN
\LessonHead{Lesson I}
\switchcolumn*

\LA
\LessonTitle{Léctio sancti Evangélii secúndum Lucam}
\switchcolumn
\EN
\LessonTitle{From the Holy Gospel according to Luke}
\switchcolumn*

\LA
\Cite{Luc 1:26-38}
\switchcolumn
\EN
\Cite{Luke 1:26-38}
\switchcolumn*

\LA
\noindent In illo témpore: Missus est Angelus Gábriel a Deo in civitátem Galilǽæ, cui nomen Názareth, ad Vírginem desponsátam viro, cui nomen erat Joseph, de domo David, et nomen Vírginis María. Et réliqua.\par
\switchcolumn
\EN
\noindent At that time the angel Gabriel was sent from God into a city of Galilee, called Nazareth, To a virgin espoused to a man whose name was Joseph, of the house of David; and the virgin's name was Mary. And so on.\par
\switchcolumn*

\LA
\LessonTitle{Homilía sancti Ambrósii Epíscopi}
\switchcolumn
\EN
\LessonTitle{Homily by St. Ambrose, Bishop (of Milan.)}
\switchcolumn*

\LA
\Source{Liber 2 in Lucam}
\switchcolumn
\EN
\Source{Bk. ii on Luke}
\switchcolumn*

\LA
\noindent Latent quidem divína mystéria, nec fácile, juxta prophéticum dictum, quisquam hóminum potest scire consílium Dei. Sed tamen ex céteris factis, atque præcéptis Dómini Salvatóris póssumus intellégere, et hoc propensióris fuísse consílii, quod ea potíssimum elécta est, ut Dóminum páreret, quæ erat desponsáta viro. Cur autem non ántequam desponsarétur, impléta est? Fortásse ne dicerétur quod concéperat ex adultério.\par
\switchcolumn
\EN
\noindent The mysteries of God are unsearchable, and it is especially declared by a Prophet, that a man can hardly know His counsels. (Wisd. ix. 13.) Nevertheless, some things have been revealed to us, and we may gather from some of the words and works of the Lord our Saviour, that there was a special purpose of God, in the fact that she who was chosen to be the mother of the Lord was espoused to a man. Why did not the power of the Highest overshadow her before she was so espoused? Perhaps it was lest any might blasphemously say that she had conceived in fornication the Holy One.\par
\switchcolumn*

\LA
\begin{Resp}\Rsign Clama in fortitúdine, qui annúntias pacem in Jerúsalem: \Star{} Dic civitátibus Judæ, et habitatóribus Sion: Ecce Deus noster, quem exspectabámus, advéniet.\end{Resp}
\begin{Resp}\Vsign Supra montem excélsum ascénde tu, qui evangelízas Sion, exálta in fortitúdine vocem tuam.\end{Resp}
\begin{Resp}\Rsign Dic civitátibus Judæ, et habitatóribus Sion: Ecce Deus noster, quem exspectabámus, advéniet.\end{Resp}
\switchcolumn
\EN
\begin{Resp}\Rsign O thou that bringest good tidings of peace to Jerusalem, lift up thy voice with strength! \Star{} Say unto the cities of Judah, and to the inhabitants of Jerusalem: Behold, our God will come, for Whom we waited.\end{Resp}
\begin{Resp}\Vsign O thou that tellest good tidings to Zion, get thee up into the high mountain, lift up thy voice with strength.\end{Resp}
\begin{Resp}\Rsign Say unto the cities of Judah, and to the inhabitants of Jerusalem: Behold, our God will come, for Whom we waited.\end{Resp}
\switchcolumn*

\LA
\LessonHead{Lectio II}
\switchcolumn
\EN
\LessonHead{Lesson II}
\switchcolumn*

\LA
\noindent Et ingréssus ad eam Angelus. Disce vírginem móribus, disce vírginem verecúndia, disce oráculo, disce mystério. Trepidáre vírginum est, et ad omnes viri ingréssus pavére, omnes viri affátus veréri. Discant mulíeres propósitum pudóris imitári. Sola in penetrálibus, quam nemo virórum víderit, solus Angelus repérerit: sola sine cómite, sola sine teste, ne quo degénere depravarétur affátu, ab Angelo salutátur.\par
\switchcolumn
\EN
\noindent And the angel came in unto her. Let us learn from this Virgin how to bear ourselves, let us learn her modesty, let us learn by her devout utterance, above all let us learn by the holy mystery enacted. It is the part of a maiden to be timid, to avoid the advances of men, and to shrink from men's addresses. Would that our women would learn from the example of modesty here set before us. She upon whom the stare of men had never been fixed was alone in her chamber, and was found only by an angel. There was neither companion nor witness there, that what passed might not be debased in gossip and the angel saluted her.\par
\switchcolumn*

\LA
\begin{Resp}\Rsign Oriétur stella ex Jacob, et exsúrget homo de Israël, et confrínget omnes duces alienigenárum: \Star{} Et erit omnis terra posséssio ejus.\end{Resp}
\begin{Resp}\Vsign Adorábunt eum omnes reges terræ, omnes gentes sérvient ei.\end{Resp}
\begin{Resp}\Rsign Et erit omnis terra posséssio ejus.\end{Resp}
\switchcolumn
\EN
\begin{Resp}\Rsign There shall come a Star out of Jacob, and a Man shall rise out of Israel, and shall smite through all the princes of the aliens. \Star{} And all the earth shall be His possession.\end{Resp}
\begin{Resp}\Vsign All kings shall fall down before Him, all nations shall serve Him.\end{Resp}
\begin{Resp}\Rsign And all the earth shall be His possession.\end{Resp}
\switchcolumn*

\LA
\LessonHead{Lectio III}
\switchcolumn
\EN
\LessonHead{Lesson III}
\switchcolumn*

\LA
\noindent Tanti namque mandáti mystérium non hóminis fuit, sed Angeli ore proméndum. Hódie primum audítur: Spíritus Sanctus supervéniet in te. Et audítur, et créditur. Dénique, Ecce, inquit, ancílla Dómini: contíngat mihi secúndum verbum tuum. Vide humilitátem, vide devotiónem. Ancíllam se dicit Dómini, quæ mater elígitur: nec repentíno exaltáta promísso est.\par
\switchcolumn
\EN
\noindent The message of God to the Virgin was a mystery, which it was not lawful for the mouth of men, but only of angels, to utter. For the first time on earth the words are spoken The Holy Ghost shall come upon thee. The holy maiden heareth, and believeth. At length she saith Behold the handmaid of the Lord be it unto me according to thy word. Here is an example of lowliness, here is a pattern of true devotion. At the very moment that she is told she is chosen to be the mother of the Lord she at once declareth herself His handmaid. The knowledge that she was Mother of God caused in the heart of Mary only an act of humility.\par
\switchcolumn*

\LA
\begin{Resp}\Rsign Modo véniet Dominátor Dóminus: \Star{} Et nomen ejus Emmánuel vocábitur.\end{Resp}
\begin{Resp}\Vsign Oriétur in diébus ejus justítia, et abundántia pacis.\end{Resp}
\begin{Resp}\Rsign Et nomen ejus Emmánuel vocábitur.\end{Resp}
\begin{Resp}\Vsign Glória Patri, et Fílio, \Star{} et Spirítui Sancto.\end{Resp}
\begin{Resp}\Rsign Et nomen ejus Emmánuel vocábitur.\end{Resp}
\switchcolumn
\EN
\begin{Resp}\Rsign The Lord, the Ruler, cometh quickly. \Star{} And His name shall be called Emmanuel.\end{Resp}
\begin{Resp}\Vsign In His days shall righteousness flourish, and abundance of peace.\end{Resp}
\begin{Resp}\Rsign And His name shall be called Emmanuel.\end{Resp}
\begin{Resp}\Vsign Glory be to the Father, and to the Son, \Star{} and to the Holy Ghost.\end{Resp}
\begin{Resp}\Rsign And His name shall be called Emmanuel.\end{Resp}
\switchcolumn*

\end{paracol}

\Office{Feria V infra Hebdomadam III Adventus}{Thursday of the Third Week of Advent}{Feria major}
\begin{paracol}{2}
\LA
\LessonHead{Lectio I}
\switchcolumn
\EN
\LessonHead{Lesson I}
\switchcolumn*

\LA
\LessonTitle{De Isaía Prophéta}
\switchcolumn
\EN
\LessonTitle{Lesson from the book of Isaias}
\switchcolumn*

\LA
\Cite{Isa 33:1-2}
\switchcolumn
\EN
\Cite{Isa 33:1-2}
\switchcolumn*

\LA
\noindent Væ qui prædáris, nonne et ipse prædáberis? et qui spernis, nonne et ipse spernéris? cum consummáveris deprædatiónem, deprædáberis: cum fatigátus desíeris contémnere, contemnéris. Dómine, miserére nostri, te enim exspectávimus: esto brácchium nostrum in mane, et salus nostra in témpore tribulatiónis.\par
\switchcolumn
\EN
\noindent Woe to thee that spoilest, shalt not thou thyself also be spoiled? and thou that despisest, shalt not thyself also be despised? when thou shalt have made an end of spoiling, thou shalt be spoiled: when being wearied thou shalt cease to despise, thou shalt be despised. O Lord, have mercy on us: for we have waited for thee: be thou our arm in the morning, and our salvation in the time of trouble.\par
\switchcolumn*

\LA
\begin{Resp}\Rsign Egrediétur Dóminus, et præliábitur contra gentes: \Star{} Et stabunt pedes ejus supra montes olivárum ad Oriéntem.\end{Resp}
\begin{Resp}\Vsign Et elevábitur supra omnes colles, et fluent ad eum omnes gentes.\end{Resp}
\begin{Resp}\Rsign Et stabunt pedes ejus supra montes olivárum ad Oriéntem.\end{Resp}
\switchcolumn
\EN
\begin{Resp}\Rsign The Lord shall go forth and fight against the nations. \Star{} And His feet shall stand upon the mount of Olives on the east.\end{Resp}
\begin{Resp}\Vsign And it shall be exalted above the hills, and all nations shall flow unto it.\end{Resp}
\begin{Resp}\Rsign And His feet shall stand upon the mount of Olives on the east.\end{Resp}
\switchcolumn*

\LA
\LessonHead{Lectio II}
\switchcolumn
\EN
\LessonHead{Lesson II}
\switchcolumn*

\LA
\Cite{Isa 33:3-6}
\switchcolumn
\EN
\Cite{Isa 33:3-6}
\switchcolumn*

\LA
\noindent A voce Angeli fugérunt pópuli, et ab exaltatióne tua dispérsæ sunt gentes. Et congregabúntur spólia vestra sicut collígitur bruchus, velut cum fossæ plenæ fúerint de eo. Magnificátus est Dóminus, quóniam habitávit in excélso: implévit Sion judício et justítia. Et erit fides in tempóribus tuis: divítiæ salútis, sapiéntia, et sciéntia: timor Dómini ipse est thesáurus ejus.\par
\switchcolumn
\EN
\noindent At the voice of the angel the people fled, and at the lifting up thyself the nations are scattered. And your spoils shall be gathered together as the locusts are gathered, as when the ditches are full of them. The Lord is magnified, for he hath dwelt on high: he hath filled Sion with judgment and justice. And there shall be faith in thy times: riches of salvation, wisdom and knowledge: the fear of the Lord is his treasure.\par
\switchcolumn*

\LA
\begin{Resp}\Rsign Præcúrsor pro nobis ingréditur Agnus sine mácula, \Star{} Secúndum órdinem Melchísedech Póntifex factus in ætérnum et in sǽculum sǽculi.\end{Resp}
\begin{Resp}\Vsign Ipse est rex justítiæ, cujus generátio non habet finem.\end{Resp}
\begin{Resp}\Rsign Secúndum órdinem Melchísedech Póntifex factus in ætérnum et in sǽculum sǽculi.\end{Resp}
\switchcolumn
\EN
\begin{Resp}\Rsign The Fore-runner is for us entered, even the Lamb without spot \Star{} Made an High Priest for ever after the order of Melchisedek.\end{Resp}
\begin{Resp}\Vsign This is that King of Righteousness without descent, nor end of life.\end{Resp}
\begin{Resp}\Rsign Made an High Priest for ever after the order of Melchisedek.\end{Resp}
\switchcolumn*

\LA
\LessonHead{Lectio III}
\switchcolumn
\EN
\LessonHead{Lesson III}
\switchcolumn*

\LA
\Cite{Isa 33:14-17}
\switchcolumn
\EN
\Cite{Isa 33:14-17}
\switchcolumn*

\LA
\noindent Contérriti sunt in Sion peccatóres, possédit tremor hypócritas. Quis póterit habitáre de vobis cum igne devoránte? quis habitábit ex vobis cum ardóribus sempitérnis? Qui ámbulat in justítiis, et lóquitur veritátem, qui próicit avarítiam ex calúmnia, et éxcutit manus suas ab omni múnere, qui obtúrat aures suas ne áudiat sánguinem, et claudit óculos suos ne vídeat malum. Iste in excélsis habitábit muniménta saxórum sublímitas ejus: panis ei datus est, aquæ ejus fidéles sunt. Regem in decóre suo vidébunt óculi ejus, cernent terram de longe.\par
\switchcolumn
\EN
\noindent The sinners in Sion are afraid, trembling hath seized upon the hypocrites. Which of you can dwell with devouring fire? which of you shall dwell with everlasting burnings? He that walketh in justices, and speaketh truth, that casteth away avarice by oppression, and shaketh his hands from all bribes, that stoppeth his ears lest he hear blood, and shutteth his eyes that he may see no evil. He shall dwell on high, the fortifications of rocks shall be his highness: bread is given him, his waters are sure. His eyes shall see the king in his beauty, they shall see the land far off.\par
\switchcolumn*

\LA
\begin{Resp}\Rsign Vidébunt gentes justum tuum, et cuncti reges ínclitum tuum: \Star{} Et vocábitur tibi nomen novum, quod os Dómini nominávit.\end{Resp}
\begin{Resp}\Vsign Et eris coróna glóriæ in manu Dómini, et diadéma regni in manu Dei tui.\end{Resp}
\begin{Resp}\Rsign Et vocábitur tibi nomen novum, quod os Dómini nominávit.\end{Resp}
\begin{Resp}\Vsign Glória Patri, et Fílio, \Star{} et Spirítui Sancto.\end{Resp}
\begin{Resp}\Rsign Et vocábitur tibi nomen novum, quod os Dómini nominávit.\end{Resp}
\switchcolumn
\EN
\begin{Resp}\Rsign The Gentiles shall see thy Righteous One, and all kings thy Glorious one. \Star{} And thou shalt be called by a new name, which the mouth of the Lord hath named.\end{Resp}
\begin{Resp}\Vsign Thou shalt also be a crown of glory in the hand of the Lord, and a royal diadem in the hand of thy God.\end{Resp}
\begin{Resp}\Rsign And thou shalt be called by a new name, which the mouth of the Lord hath named.\end{Resp}
\begin{Resp}\Vsign Glory be to the Father, and to the Son, \Star{} and to the Holy Ghost.\end{Resp}
\begin{Resp}\Rsign And thou shalt be called by a new name, which the mouth of the Lord hath named.\end{Resp}
\switchcolumn*

\end{paracol}

\Office{Feria VI Quattuor Temporum in Adventu}{Ember Friday of Advent}{Feria major}
\begin{paracol}{2}
\LA
\LessonHead{Lectio I}
\switchcolumn
\EN
\LessonHead{Lesson I}
\switchcolumn*

\LA
\LessonTitle{Léctio sancti Evangélii secúndum Lucam}
\switchcolumn
\EN
\LessonTitle{Continuation of the Holy Gospel according to Luke}
\switchcolumn*

\LA
\Cite{Luc 1:39-47}
\switchcolumn
\EN
\Cite{Luke 1:39-47}
\switchcolumn*

\LA
\noindent In illo témpore: Exsúrgens María ábiit in montána cum festinatióne in civitátem Juda: et intrávit in domum Zacharíæ, et salutávit Elísabeth. Et réliqua.\par
\switchcolumn
\EN
\noindent In that time, Mary, rising up, went into the hill country with haste into a city of Juda. And she entered into the house of Zachary, and saluted Elizabeth. And so on.\par
\switchcolumn*

\LA
\LessonTitle{Homilía sancti Ambrósii Epíscopi}
\switchcolumn
\EN
\LessonTitle{Homily by St. Ambrose, Bishop (of Milan)}
\switchcolumn*

\LA
\Source{Lib. 2 in Luc. cap. 1 post init.}
\switchcolumn
\EN
\Source{Commentary on Luke, Bk. ii. c}
\switchcolumn*

\LA
\noindent Morále est ómnibus, ut qui fidem éxigunt, fidem ástruant. Et ídeo Angelus, cum abscóndita nuntiáret, ut fidem astrúeret exémplo, senióris féminæ sterilísque concéptum Vírgini Maríæ nuntiávit: ut possíbile Deo omne, quod ei placúerit, asséreret. Ubi audívit hoc María, non quasi incrédula de oráculo, nec quasi incérta de núntio, nec quasi dúbitans de exémplo, sed quasi læta pro voto, religiósa pro offício, festína pro gáudio, in montána perréxit. Quo enim jam Deo plena, nisi ad superióra cum festinatióne conténderet? Nescit tarda molímina Sancti Spíritus grátia.\par
\switchcolumn
\EN
\noindent When any one asketh another for credence, he is bound to give some reasonable ground. And so the Angel, when he announced to Mary the counsel of God, gave, as a proof, the conception of Elizabeth, then aged and barren, that Mary might perceive, by this example, that with God nothing is impossible. When the holy Virgin had heard it, she arose and went to visit her cousin. She did not go to see if what she had heard was true, because she did not believe God, or because she knew not who the messenger had been, or yet because she doubted the fact adduced in proof. She went joyfully as one who hath received a mercy in answer to his vow goeth to pay the same. She went with devotion, as a godly person goeth to execute a religious duty. She went into the hill country in joyful haste. And is it not something that she went up into the hills? God was already in her womb, and her feeling bore her continually upward. The grace of the Holy Spirit knoweth no slow working.\par
\switchcolumn*

\LA
\begin{Resp}\Rsign Emítte Agnum, Dómine, Dominatórem terræ, \Star{} De Petra desérti ad montem fíliæ Sion.\end{Resp}
\begin{Resp}\Vsign Osténde nobis, Dómine, misericórdiam tuam, et salutáre tuum da nobis.\end{Resp}
\begin{Resp}\Rsign De Petra desérti ad montem fíliæ Sion.\end{Resp}
\switchcolumn
\EN
\begin{Resp}\Rsign Send forth the Lamb, O Lord, the Ruler of the land; \Star{} From the rock in the wilderness unto the mount of the daughter of Zion.\end{Resp}
\begin{Resp}\Vsign Show us thy mercy, O Lord, and grant us thy salvation.\end{Resp}
\begin{Resp}\Rsign From the rock in the wilderness unto the mount of the daughter of Zion.\end{Resp}
\switchcolumn*

\LA
\LessonHead{Lectio II}
\switchcolumn
\EN
\LessonHead{Lesson II}
\switchcolumn*

\LA
\noindent Díscite et vos, sanctæ mulíeres, sedulitátem, quam prægnántibus debeátis exhibére cognátis. Maríam, quæ ante sola in íntimis penetrálibus versabátur, non a público virginitátis pudor, non ab stúdio aspéritas móntium, non ab offício prolíxitas itíneris retardávit. In montána Virgo cum festinatióne, Virgo offícii memor, injúriæ ímmemor, afféctu urgénte non sexu, relícta perréxit domo. Díscite, vírgines, non circumcursáre per aliénas ædes, non demorári in platéis, non áliquos in público miscére sermónes. María in domo sera, festína in público, mansit apud cognátam suam tribus ménsibus.\par
\switchcolumn
\EN
\noindent Godly women will learn from the example of the Mother of God to take a tender care of their kinswomen who are with child. In pursuance of this charity, Mary, who had hitherto remained alone at home, was not deterred by her maidenly shyness from entering on a public journey; she faced for this end the hardships of mountain travelling; and encountered with a sense of duty the weary length of the way. The Virgin left her home, and went into the hill country with haste, unmindful of the trouble, and remembering only the office to which her cousinly love prompted her, in spite of the delicacy of her sex. Maidens will learn from her not to idle about from house to house, to loiter in the streets, nor to take part in conversations in public. Mary, as she was hasteful to pass through the public roads, so was she slow again to enter on them she abode with her cousin about three months.\par
\switchcolumn*

\LA
\begin{Resp}\Rsign Roráte, cæli, désuper, et nubes pluant justum: \Star{} Aperiátur terra, et gérminet Salvatórem.\end{Resp}
\begin{Resp}\Vsign Emítte Agnum, Dómine, Dominatórem terræ, de Petra desérti ad montem fíliæ Sion.\end{Resp}
\begin{Resp}\Rsign Aperiátur terra, et gérminet Salvatórem.\end{Resp}
\switchcolumn
\EN
\begin{Resp}\Rsign Drop down, ye heavens, from above, and let the skies pour down the Righteous One. \Star{} Let the earth open, and let her bring forth the Saviour.\end{Resp}
\begin{Resp}\Vsign Send forth the Lamb, O Lord, the Ruler of the land, from the rock in the wilderness unto the mount of the daughter of Zion.\end{Resp}
\begin{Resp}\Rsign Let the earth open, and let her bring forth the Saviour.\end{Resp}
\switchcolumn*

\LA
\LessonHead{Lectio III}
\switchcolumn
\EN
\LessonHead{Lesson III}
\switchcolumn*

\LA
\noindent Didicístis, vírgines, pudórem Maríæ: díscite humilitátem. Venit propínqua ad próximam, júnior ad seniórem: nec solum venit, sed étiam prior salutávit. Decet enim, ut quanto cástior virgo, tanto humílior sit. Nóverit deférre senióribus. Sit magístra humilitátis, in qua est proféssio castitátis. Est et causa pietátis, est étiam norma doctrínæ. Contuéndum est enim, quia supérior venit ad inferiórem, ut inférior adjuvétur: María ad Elísabeth, Christus ad Joánnem.\par
\switchcolumn
\EN
\noindent As the modesty of Mary is a pattern for the imitation of all maidens, so also is her humility. She went to see Elizabeth, like one cousin going to visit another, and as the younger to the elder. Not only did she first go, but she first saluted Elizabeth. Now, the purer a virgin is, the humbler ought she to be. She will know how to submit herself to her elders. She that professeth chastity ought to be a very mistress of humility. Lowly-mindedness is at once the very ground in which devotion groweth, and the first and principal rule of its teaching. In this act of the Virgin then we see the greater going to visit and to succour the lesser: Mary to Elizabeth, Christ to John.\par
\switchcolumn*

\LA
\begin{Resp}\Rsign Germinavérunt campi erémi germen odóris Israël: quia ecce Deus noster cum virtúte véniet, \Star{} Et splendor ejus cum eo.\end{Resp}
\begin{Resp}\Vsign Ex Sion spécies decóris ejus: Deus noster maniféste véniet.\end{Resp}
\begin{Resp}\Rsign Et splendor ejus cum eo.\end{Resp}
\begin{Resp}\Vsign Glória Patri, et Fílio, \Star{} et Spirítui Sancto.\end{Resp}
\begin{Resp}\Rsign Et splendor ejus cum eo.\end{Resp}
\switchcolumn
\EN
\begin{Resp}\Rsign The waste places have brought forth sweet-smelling buds for Israel; for, behold, our God will come with power. \Star{} And His brightness is with Him.\end{Resp}
\begin{Resp}\Vsign Out of Zion the perfection of beauty, our God shall come manifestly.\end{Resp}
\begin{Resp}\Rsign And His brightness is with Him.\end{Resp}
\begin{Resp}\Vsign Glory be to the Father, and to the Son, \Star{} and to the Holy Ghost.\end{Resp}
\begin{Resp}\Rsign And His brightness is with Him.\end{Resp}
\switchcolumn*

\end{paracol}

\Office{Sabbato Quattuor Temporum in Adventu}{Ember Saturday of Advent}{Feria major}
\begin{paracol}{2}
\LA
\LessonHead{Lectio I}
\switchcolumn
\EN
\LessonHead{Lesson I}
\switchcolumn*

\LA
\LessonTitle{Léctio sancti Evangélii secúndum Lucam}
\switchcolumn
\EN
\LessonTitle{Continuation of the Holy Gospel according to Luke}
\switchcolumn*

\LA
\Cite{Luc 3:1-6}
\switchcolumn
\EN
\Cite{Luke 3:1-6}
\switchcolumn*

\LA
\noindent Anno quintodécimo impérii Tibérii Cǽsaris, procuránte Póntio Piláto Judǽam. Et réliqua.\par
\switchcolumn
\EN
\noindent In the fifteenth year of the reign of Tiberius Caesar, Pontius Pilate being governor of Judea. And so on.\par
\switchcolumn*

\LA
\LessonTitle{Homilía sancti Gregórii Papæ}
\switchcolumn
\EN
\LessonTitle{Homily by Pope St. Gregory (the Great.)}
\switchcolumn*

\LA
\Source{Homilia 20 in Evangelia}
\switchcolumn
\EN
\Source{10th on the Gospels.}
\switchcolumn*

\LA
\noindent Redemptóris nostri Præcúrsor, quo témpore prædicatiónis offícium accéperit, memoráto Románæ reipúblicæ príncipe, et Judǽæ régibus, designátur. Quia enim illum prædicáre veniébat, qui et ex Judǽa quosdam, et multos ex géntibus redemptúrus erat: per regem géntium et príncipes Judæórum prædicatiónis ejus témpora designántur. Quia autem gentílitas colligénda erat, et Judǽa pro culpa perfídiæ dispergénda, ipsa quoque descríptio terréni principátus osténdit: quóniam et in Romána república unus præfuísse descríbitur, et in Judǽæ regno per quartam partem plúrimi principabántur.\par
\switchcolumn
\EN
\noindent The date, at which the Fore-runner of our Redeemer entered on his public office of preaching, is indicated to us by the name of the ruler of the Roman Commonwealth, and by those of the princes of Palestine. The time of his preaching is indicated by these names, because he came as the Fore-runner of Him Who was to be the Redeemer of some Jews and many Gentiles. Moreover in the enumeration of these worldly monarchs there is a foreshadowing of the fact, that the Gentiles were about to be gathered into one, and the Jews to be scattered abroad in punishment of their unbelief; in the whole heathen Commonwealth we find the title of one Emperor, but in the small kingdom of Judaea are mentioned four masters.\par
\switchcolumn*

\LA
\begin{Resp}\Rsign Egrediétur virga de radíce Jesse, et flos de radíce ejus ascéndet: \Star{} Et erit justítia cíngulum lumbórum ejus, et fides cinctórium renum ejus.\end{Resp}
\begin{Resp}\Vsign Et requiéscet super eum spíritus Dómini: spíritus sapiéntiæ, et intelléctus: spíritus consílii, et fortitúdinis.\end{Resp}
\begin{Resp}\Rsign Et erit justítia cíngulum lumbórum ejus, et fides cinctórium renum ejus.\end{Resp}
\switchcolumn
\EN
\begin{Resp}\Rsign There shall come forth a rod out of the stem of Jesse, and a flower shall grow out of his roots. \Star{} And righteousness shall be the girdle of his loins, and faithfulness the girdle of his reins.\end{Resp}
\begin{Resp}\Vsign And the Spirit of the Lord shall rest upon him the spirit of wisdom and understanding the spirit of counsel and might.\end{Resp}
\begin{Resp}\Rsign And righteousness shall be the girdle of his loins, and faithfulness the girdle of his reins.\end{Resp}
\switchcolumn*

\LA
\LessonHead{Lectio II}
\switchcolumn
\EN
\LessonHead{Lesson II}
\switchcolumn*

\LA
\noindent Voce étenim nostri Redemptóris dícitur: Omne regnum in seípsum divisum desolábitur. Liquet ergo, quod ad finem regni Judǽa pervénerat, quæ tot régibus divísa subjacébat. Apte quoque non solum quibus régibus, sed étiam quibus sacerdótibus actum sit, demonstrátur: et quia illum Joánnes Baptísta prædicáret, qui simul Rex et Sacérdos exsísteret, Lucas Evangelísta prædicatiónis ejus témpora per regnum et sacerdótium designávit.\par
\switchcolumn
\EN
\noindent The blessed voice of the Saviour itself hath said, Every kingdom divided against itself is brought to desolation (Luke xi. 17.) And we may well look for the ruin of the Jewish state when we see it divided among so many rulers. We observe likewise that the names of the reigning priests as well as kings are given. The Evangelist Luke hath left on record the chiefs both of the monarchy and of the priesthood who held office when John the Baptist began to preach, because John preached Him Who is at once our Priest and our King.\par
\switchcolumn*

\LA
\begin{Resp}\Rsign Radix Jesse, qui exsúrget judicáre gentes, in eum gentes sperábunt: \Star{} Et erit nomen ejus benedíctum in sǽcula.\end{Resp}
\begin{Resp}\Vsign Super ipsum continébunt reges os suum, ipsum gentes deprecabúntur.\end{Resp}
\begin{Resp}\Rsign Et erit nomen ejus benedíctum in sǽcula.\end{Resp}
\switchcolumn
\EN
\begin{Resp}\Rsign Behold, the root of Jesse that shall arise to bring forth judgment to the Gentiles, in him shall the Gentiles trust. \Star{} And his name shall be blessed for ever.\end{Resp}
\begin{Resp}\Vsign The Kings shall shut their mouths at him, to him shall the Gentiles seek.\end{Resp}
\begin{Resp}\Rsign And his name shall be blessed for ever.\end{Resp}
\switchcolumn*

\LA
\LessonHead{Lectio III}
\switchcolumn
\EN
\LessonHead{Lesson III}
\switchcolumn*

\LA
\noindent Et venit in omnem regiónem Jordánis, prǽdicans baptísmum pœniténtiæ in remissiónem peccatórum. Cunctis legéntibus liquet, quia Joánnes non solum baptísmum pœniténtiæ prædicávit, verum étiam quibúsdam dedit: sed tamen baptísmum suum in remissiónem peccatórum dare non pótuit. Remíssio étenim peccatórum in solo nobis baptísmo Christi tribúitur. Notándum ítaque, quod dícitur: Prǽdicans baptísmum pœniténtiæ in remissiónem peccatórum: quóniam baptísmum, quod peccáta sólveret, quia dare non póterat, prædicábat: ut sicut incarnátum Verbum Patris præcurrébat verbo prædicatiónis, ita baptísmum pœniténtiæ, quo peccáta solvúntur, præcúrreret suo baptísmate, quo peccáta solvi non possunt.\par
\switchcolumn
\EN
\noindent And he came into all the country about Jordan, preaching the baptism of repentance for the remission of sins. It is evident from these words that John the Baptist not only preached, but also administered the baptism of repentance, and yet that baptism of repentance which he gave, was not really a baptism for the remission of sins. For there is only one baptism for the remission of sins, and that is our Christian baptism. It is worthy of note here that the words used are, preaching the baptism of repentance for the remission of sins, for he himself owned that his baptism was not the true baptism that washes away sin. Even as the Eternal Word of God made Flesh was greater than the preacher that went before Him, so was His holy baptism, by which our sins are washed away, far greater than that baptism of repentance which the Fore-runner preached, and which could never wash away sin.\par
\switchcolumn*

\LA
\begin{Resp}\Rsign Veni, Dómine, et noli tardáre: reláxa facínora plebi tuæ, \Star{} Et révoca dispérsos in terram suam.\end{Resp}
\begin{Resp}\Vsign Excíta, Dómine, poténtiam tuam, et veni, ut salvos fácias nos.\end{Resp}
\begin{Resp}\Rsign Et révoca dispérsos in terram suam.\end{Resp}
\begin{Resp}\Vsign Glória Patri, et Fílio, \Star{} et Spirítui Sancto.\end{Resp}
\begin{Resp}\Rsign Et révoca dispérsos in terram suam.\end{Resp}
\switchcolumn
\EN
\begin{Resp}\Rsign O Lord, come, and make no tarrying loosen the bonds of thy people. \Star{} And gather together into their own land them that are scattered abroad.\end{Resp}
\begin{Resp}\Vsign Stir up, O Lord, thy power, and come among us, to save us.\end{Resp}
\begin{Resp}\Rsign And gather together into their own land them that are scattered abroad.\end{Resp}
\begin{Resp}\Vsign Glory be to the Father, and to the Son, \Star{} and to the Holy Ghost.\end{Resp}
\begin{Resp}\Rsign And gather together into their own land them that are scattered abroad.\end{Resp}
\switchcolumn*

\end{paracol}

\Office{Dominica IV Adventus}{Fourth Sunday of Advent}{Semiduplex II classis}
\begin{paracol}{2}
\LA
\LessonHead{Lectio I}
\switchcolumn
\EN
\LessonHead{Lesson I}
\switchcolumn*

\LA
\LessonTitle{De Isaía Prophéta}
\switchcolumn
\EN
\LessonTitle{Lesson from the book of Isaias}
\switchcolumn*

\LA
\Cite{Isa 35:1-7}
\switchcolumn
\EN
\Cite{Isa 35:1-7}
\switchcolumn*

\LA
\noindent Lætábitur desérta et ínvia, et exsultábit solitúdo, et florébit quasi lílium. Gérminans germinábit, et exsultábit lætabúnda et laudans: glória Líbani data est ei: decor Carméli et Saron, ipsi vidébunt glóriam Dómini, et decórem Dei nostri. Confortáte manus dissolútas, et génua debília roboráte. Dícite pusillánimis: Confortámini, et nolíte timére: ecce Deus vester ultiónem addúcet retributiónis: Deus ipse véniet, et salvábit vos. Tunc aperiéntur óculi cæcórum, et aures surdórum patébunt. Tunc sáliet sicut cervus claudus, et apérta erit lingua mutórum: quia scissæ sunt in desérto aquæ, et torréntes in solitúdine. Et quæ erat árida, erit in stagnum, et sítiens in fontes aquárum.\par
\switchcolumn
\EN
\noindent The land that was desolate and impassable shall be glad, and the wilderness shall rejoice, and shall flourish like the lily. It shall bud forth and blossom, and shall rejoice with joy and praise: the glory of Libanus is given to it: the beauty of Carmel, and Saron, they shall see the glory of the Lord, and the beauty of our God. Strengthen ye the feeble hands, and confirm the weak knees. Say to the fainthearted: Take courage, and fear not: behold your God will bring the revenge of recompense: God himself will come and will save you. Then shall the eyes of the blind be opened, and the ears of the deaf shall be unstopped. Then shall the lame man leap as a hart, and the tongue of the dumb shall be free: for waters are broken out in the desert, and streams in the wilderness. And that which was dry land, shall become a pool, and the thirsty land springs of water.\par
\switchcolumn*

\LA
\begin{Resp}\Rsign Cánite tuba in Sion, vocáte gentes, annuntiáte pópulis, et dícite: \Star{} Ecce Deus Salvátor noster advéniet.\end{Resp}
\begin{Resp}\Vsign Annuntiáte, et audítum fácite: loquímini, et clamáte.\end{Resp}
\begin{Resp}\Rsign Ecce Deus Salvátor noster advéniet.\end{Resp}
\switchcolumn
\EN
\begin{Resp}\Rsign Blow ye the trumpet in Zion, call together the nations, tell it out among the people, and say: \Star{} Behold, God our Saviour cometh.\end{Resp}
\begin{Resp}\Vsign Tell it out and make it to be heard; speak aloud and cry.\end{Resp}
\begin{Resp}\Rsign Behold, God our Saviour cometh.\end{Resp}
\switchcolumn*

\LA
\LessonHead{Lectio II}
\switchcolumn
\EN
\LessonHead{Lesson II}
\switchcolumn*

\LA
\Cite{Isa 35:7-10}
\switchcolumn
\EN
\Cite{Isa 35:7-10}
\switchcolumn*

\LA
\noindent In cubílibus, in quibus prius dracónes habitábant, oriétur viror cálami et junci. Et erit ibi sémita et via, et via sancta vocábitur: non transíbit per eam pollútus, et hæc erit vobis dirécta via, ita ut stulti non errent per eam. Non erit ibi leo, et mala béstia non ascéndet per eam, nec inveniétur ibi: et ambulábunt, qui liberáti fúerint. Et redémpti a Dómino converténtur, et vénient in Sion cum laude: et lætítia sempitérna super caput eórum: gáudium et lætítiam obtinébunt, et fúgiet dolor et gémitus.\par
\switchcolumn
\EN
\noindent In the dens where dragons dwell before, shall rise up the verdure of the reed and the bulrush. And a path and a way shall be there, and it shall be called the holy way: the unclean shall not pass over it, and this shall be unto you a straight way, so that fools shall not err therein. No lion shall be there, nor shall any mischievous beast go up by it, nor be found there: but they shall walk there that shall be delivered. And the redeemed of the Lord shall return, and shall come into Sion with praise, and everlasting joy shall be upon their heads: they shall obtain joy and gladness, and sorrow and mourning shall flee away.\par
\switchcolumn*

\LA
\begin{Resp}\Rsign Non auferétur sceptrum de Juda, et dux de fémore ejus, donec véniat qui mitténdus est: \Star{} Et ipse erit exspectátio géntium.\end{Resp}
\begin{Resp}\Vsign Pulchrióres sunt óculi ejus vino, et dentes ejus lacte candidióres.\end{Resp}
\begin{Resp}\Rsign Et ipse erit exspectátio géntium.\end{Resp}
\switchcolumn
\EN
\begin{Resp}\Rsign The sceptre shall not depart from Judah, nor the law-giver from his loins, until he that shall be sent cometh. \Star{} And unto him shall the longing of the Gentiles be.\end{Resp}
\begin{Resp}\Vsign His eyes shall be bright with wine, and his teeth white with milk.\end{Resp}
\begin{Resp}\Rsign And unto him shall the desire of the Gentiles be.\end{Resp}
\switchcolumn*

\LA
\LessonHead{Lectio III}
\switchcolumn
\EN
\LessonHead{Lesson III}
\switchcolumn*

\LA
\Cite{Isa 41:1-4}
\switchcolumn
\EN
\Cite{Isa 41:1-4}
\switchcolumn*

\LA
\noindent Táceant ad me ínsulæ, et gentes mutent fortitúdinem: accédant, et tunc loquántur, simul ad judícium propinquémus. Quis suscitávit ab Oriénte justum, vocávit eum ut sequerétur se? dabit in conspéctu ejus gentes et reges obtinébit: dabit quasi púlverem gládio ejus, sicut stípulam vento raptam árcui ejus. Persequétur eos, transíbit in pace, sémita in pédibus ejus non apparébit. Quis hæc operátus est, et fecit, vocans generatiónes ab exórdio? Ego Dóminus, primus et novíssimus ego sum.\par
\switchcolumn
\EN
\noindent Let the islands keep silence before me, and the nations take new strength: let them come near, and then speak, let us come near to judgment together. Who hath raised up the just one from the east, hath called him to follow him? he shall give the nations in his sight, and he shall rule over kings: he shall give them as the dust to his sword, as stubble driven by the wind, to his bow. He shall pursue them, he shall pass in peace, no path shall appear after his feet. Who hath wrought and done these things, calling the generations from the beginning? I the Lord, I am the first and the last.\par
\switchcolumn*

\LA
\begin{Resp}\Rsign Me opórtet mínui, illum autem créscere: qui autem post me venit, ante me factus est: \Star{} Cujus non sum dignus corrígiam calceamentórum sólvere.\end{Resp}
\begin{Resp}\Vsign Ego baptizávi vos aqua: ille autem baptizábit vos Spíritu Sancto.\end{Resp}
\begin{Resp}\Rsign Cujus non sum dignus corrígiam calceamentórum sólvere.\end{Resp}
\begin{Resp}\Vsign Glória Patri, et Fílio, \Star{} et Spirítui Sancto.\end{Resp}
\begin{Resp}\Rsign Cujus non sum dignus corrígiam calceamentórum sólvere.\end{Resp}
\switchcolumn
\EN
\begin{Resp}\Rsign I must decrease, but He must increase; He it is Who, coming after me is preferred before me; \Star{} Whose shoe's latchet I am not worthy to unloose.\end{Resp}
\begin{Resp}\Vsign I baptize you with water; but He shall baptize you with the Holy Ghost.\end{Resp}
\begin{Resp}\Rsign Whose shoe's latchet I am not worthy to unloose.\end{Resp}
\begin{Resp}\Vsign Glory be to the Father, and to the Son, \Star{} and to the Holy Ghost.\end{Resp}
\begin{Resp}\Rsign Whose shoe's latchet I am not worthy to unloose.\end{Resp}
\switchcolumn*

\LA
\LessonHead{Lectio IV}
\switchcolumn
\EN
\LessonHead{Lesson IV}
\switchcolumn*

\LA
\LessonTitle{Sermo sancti Leónis Papæ}
\switchcolumn
\EN
\LessonTitle{From the Sermons of Pope St. Leo (the Great)}
\switchcolumn*

\LA
\Source{Sermo 1 de jejúnio décimi mensis, et collectis}
\switchcolumn
\EN
\Source{1st on the December Fast, and Almsgiving}
\switchcolumn*

\LA
\noindent Si fidéliter, dilectíssimi, atque sapiénter creatiónis nostræ intelligámus exórdium, inveniémus hóminem ideo ad imáginem Dei cónditum, ut imitátor sui esset auctóris: et hanc esse naturálem nostri géneris dignitátem, si in nobis, quasi in quodam spéculo, divínæ benignitátis forma respléndeat. Ad quam quotídie nos útique réparat grátia Salvatóris, dum quod cécidit in Adam primo, erígitur in secúndo.\par
\switchcolumn
\EN
\noindent Dearly beloved brethren, if we study attentively the history of the creation of our race, we shall find that man was made in the image of God, that his ways also might be an imitation of the ways of his Maker. This is the natural, real, and highest dignity to which we are capable of attaining, that the goodness of the Divine nature should have a reflection in us, as in a glass. As a mean of reaching this dignity, we are daily offered the grace of our Saviour, for as in the first Adam all men are fallen, so in the Second Adam can all men be raised up again (i Cor. xv. 22).\par
\switchcolumn*

\LA
\begin{Resp}\Rsign Nascétur nobis párvulus, et vocábitur Deus, Fortis: \Star{} Ipse sedébit super thronum David patris sui, et imperábit: cujus potéstas super húmerum ejus.\end{Resp}
\begin{Resp}\Vsign In ipso benedicéntur omnes tribus terræ, omnes gentes sérvient ei.\end{Resp}
\begin{Resp}\Rsign Ipse sedébit super thronum David patris sui, et imperábit: cujus potéstas super húmerum ejus.\end{Resp}
\switchcolumn
\EN
\begin{Resp}\Rsign Unto us shall a Child be born, and His name shall be called the Mighty God. \Star{} He shall sit upon the throne of His father David, and shall reign, and the government shall be upon His shoulder.\end{Resp}
\begin{Resp}\Vsign In Him shall all the kindreds of the earth be blessed; all nations shall serve Him.\end{Resp}
\begin{Resp}\Rsign He shall sit upon the throne of His father David, and shall reign, and the government shall be upon His shoulder.\end{Resp}
\switchcolumn*

\LA
\LessonHead{Lectio V}
\switchcolumn
\EN
\LessonHead{Lesson V}
\switchcolumn*

\LA
\noindent Causa autem reparatiónis nostræ non est nisi misericórdia Dei: quem non diligerémus, nisi prius nos ipse dilígeret, et ténebras ignorántiæ nostræ, suæ veritátis luce discúteret. Quod per sanctum Isaíam Dóminus denúntians, ait: Addúcam cæcos in viam quam ignorábant, et sémitas quas nesciébant, fáciam illos calcáre: fáciam illis ténebras in lucem, et prava in dirécta. Hæc verba fáciam illis, et non relínquam eos. Et íterum: Invéntus sum, inquit, a non quæréntibus me, et palam appárui iis qui me non interrogábant.\par
\switchcolumn
\EN
\noindent Our restoration from the consequences of Adam's fall is sheer mercy of God, and nothing else; we should not have loved Him unless He had first loved us, (1 John iv. 19,) and scattered the darkness of our ignorance by the light of His truth. This the Lord promised by the mouth of Isaiah, where He saith: (Isa. xlii. 16,) I will bring the blind by a way that they knew not, and I will lead them in paths that they have not known I will make darkness light before them, and crooked things straight. These things will I do unto them and not forsake them. And again: (Isa. lxv. 1, 2; Rom. x. 20,) I was found of them that sought Me not; I was made manifest unto them that asked not after Me.\par
\switchcolumn*

\LA
\begin{Resp}\Rsign Ecce jam venit plenitúdo témporis, in quo misit Deus Fílium suum in terras, natum de Vírgine, factum sub lege: \Star{} Ut eos, qui sub lege erant, redímeret.\end{Resp}
\begin{Resp}\Vsign Propter nímiam caritátem suam, qua diléxit nos Deus, Fílium suum misit in similitúdinem carnis peccáti.\end{Resp}
\begin{Resp}\Rsign Ut eos, qui sub lege erant, redímeret.\end{Resp}
\switchcolumn
\EN
\begin{Resp}\Rsign Behold, the fullness of the time is come, wherein God hath sent forth His Son into the world, born of a Virgin, made under the law; \Star{} To redeem them that were under the law.\end{Resp}
\begin{Resp}\Vsign God, for His great love wherewith He loved us, hath sent His own Son in the likeness of sinful flesh.\end{Resp}
\begin{Resp}\Rsign To redeem them that were under the law.\end{Resp}
\switchcolumn*

\LA
\LessonHead{Lectio VI}
\switchcolumn
\EN
\LessonHead{Lesson VI}
\switchcolumn*

\LA
\noindent Quod quómodo implétum sit, Joánnes Apóstolus docet, dicens: Scimus quóniam Fílius Dei venit, et dedit nobis sensum, ut cognoscámus verum, et simus in vero Fílio ejus. Et íterum: Nos ergo diligámus Deum, quóniam ipse prior diléxit nos. Diligéndo ítaque nos Deus, ad imáginem suam nos réparat: et ut in nobis formam suæ bonitátis invéniat, dat unde ipsi quoque quod operátur operémur, accéndens scílicet méntium nostrárum lucérnas, et igne nos suæ caritátis inflámmans, ut non solum ipsum, sed étiam quidquid díligit, diligámus.\par
\switchcolumn
\EN
\noindent And we know from the Apostle John how God fulfilled His promise, (1 John v. 20.) We know that the Son of God is come, and hath given us an understanding, that we may know Him That is True, and be in Him That is True, even in His Son. And again: (iv. 19,) Let us therefore love God, because He first loved us. For His great love then wherewith he hath loved us, (Eph. ii. 4,) God reneweth His likeness in us. And, moreover, in order that He may find in us the reflection of His goodness, He giveth us that whereby to work along with Himself, (Who worketh all in all,) lighting, as it were, candles in our dark minds, and kindling in us the fire of His love, to make us love not Himself only, but likewise, in Him, whatsoever He loveth.\par
\switchcolumn*

\LA
\begin{Resp}\Rsign Virgo Israël, revértere ad civitátes tuas: \Star{} Usquequo dolens avertéris? generábis Dóminum Salvatórem, oblatiónem novam in terra: * Ambulábunt hómines in salvatiónem.\end{Resp}
\begin{Resp}\Vsign In caritáte perpétua diléxi te: ídeo attráxi te míserans tui.\end{Resp}
\begin{Resp}\Rsign Usquequo dolens avertéris? generábis Dóminum Salvatórem, oblatiónem novam in terra.\end{Resp}
\begin{Resp}\Vsign Glória Patri, et Fílio, \Star{} et Spirítui Sancto.\end{Resp}
\begin{Resp}\Rsign Ambulábunt hómines in salvatiónem.\end{Resp}
\switchcolumn
\EN
\begin{Resp}\Rsign O virgin of Israel, turn again to thy cities. \Star{} How long wilt thou go about sorrowing? Thou shalt bring forth the Lord thy Saviour, a new offering in the earth; men shall walk in paths of salvation.\end{Resp}
\begin{Resp}\Vsign I have loved thee with an everlasting love; therefore with loving-kindness have I drawn thee.\end{Resp}
\begin{Resp}\Rsign How long wilt thou go about sorrowing? Thou shalt bring forth the Lord thy Saviour, a new offering in the earth.\end{Resp}
\begin{Resp}\Vsign Glory be to the Father, and to the Son, \Star{} and to the Holy Ghost.\end{Resp}
\begin{Resp}\Rsign Men shall walk in paths of salvation.\end{Resp}
\switchcolumn*

\LA
\LessonHead{Lectio VII}
\switchcolumn
\EN
\LessonHead{Lesson VII}
\switchcolumn*

\LA
\LessonTitle{Léctio sancti Evangélii secúndum Lucam}
\switchcolumn
\EN
\LessonTitle{From the Holy Gospel according to Luke}
\switchcolumn*

\LA
\Cite{Luc 3:1-6}
\switchcolumn
\EN
\Cite{Luke 3:1-6}
\switchcolumn*

\LA
\noindent Anno quintodécimo impérii Tibérii Cǽsaris, procuránte Póntio Piláto Judǽam. Et réliqua.\par
\switchcolumn
\EN
\noindent In the fifteenth year of the reign of Tiberius Caesar, Pontius Pilate being governor of Judea. And so on.\par
\switchcolumn*

\LA
\LessonTitle{De Homilía sancti Gregórii Papæ}
\switchcolumn
\EN
\LessonTitle{Homily by Pope St. Gregory (the Great)}
\switchcolumn*

\LA
\Source{Homilia 20 in Evangelia ante medium}
\switchcolumn
\EN
\Source{20th on the Gospels}
\switchcolumn*

\LA
\noindent Dicébat Joánnes ad turbas, quæ éxibant ut baptizaréntur ab eo: Genímina viperárum, quis osténdit vobis fúgere a ventúra ira? Ventúra enim ira est animadvérsio ultiónis extrémæ: quam tunc fúgere peccátor non valet, qui nunc ad laménta pœniténtiæ non recúrrit. Et notándum, quod malæ sóboles, malórum paréntum actiónem imitántes, genímina viperárum vocántur: quia per hoc quod bonis ínvident, eósque persequúntur, quod quibúsdam mala retríbuunt, quod læsiónes próximis exquírunt: quóniam in his ómnibus patrum suórum carnálium vias sequúntur, quasi venenáti fílii de venenátis paréntibus nati sunt.\par
\switchcolumn
\EN
\noindent John said unto the multitude, that came forth to be baptised of him: O generation of vipers, who hath warned you to flee from the wrath to come? The wrath to come in one sense signifieth the great vengeance of the Latter Day the sinner that repenteth not of his sin now, will have no mean whereby to flee from punishment then. Let us remark that addressing evil children copying the example of evil parents, the Baptist calleth them a generation of vipers in that they were envious at the righteous, and persecuted them; that they repaid evil for evil; that they hunted out ways of harming their neighbours, in all these things following the pattern of carnal parents, the prophet likeneth them to a venomous brood hatched from a venomous stock.\par
\switchcolumn*

\LA
\begin{Resp}\Rsign Jurávi, dicit Dóminus, ut ultra jam non iráscar super terram: montes enim et colles suscípient justítiam meam, \Star{} Et testaméntum pacis erit in Jerúsalem.\end{Resp}
\begin{Resp}\Vsign Juxta est salus mea, ut véniat: et justítia mea, ut revelétur.\end{Resp}
\begin{Resp}\Rsign Et testaméntum pacis erit in Jerúsalem.\end{Resp}
\switchcolumn
\EN
\begin{Resp}\Rsign I have sworn, saith the Lord, that I will not be wroth any more with the earth; for the mountains and the hills shall receive My righteousness. \Star{} And the covenant of My peace shall be in Jerusalem.\end{Resp}
\begin{Resp}\Vsign My salvation is near to come, and My righteousness to be revealed.\end{Resp}
\begin{Resp}\Rsign And the covenant of My peace shall be in Jerusalem.\end{Resp}
\switchcolumn*

\LA
\LessonHead{Lectio VIII}
\switchcolumn
\EN
\LessonHead{Lesson VIII}
\switchcolumn*

\LA
\noindent Sed quia jam peccávimus, quia usu malæ consuetúdinis involúti sumus: dicat quid nobis faciéndum sit, ut fúgere a ventúra ira valeámus. Séquitur: Fácite ergo fructus dignos pœniténtiæ. In quibus verbis notándum est, quod amícus Sponsi non solum fructus pœniténtiæ, sed dignos pœniténtiæ ádmonet esse faciéndos. Aliud namque est fructum fácere pœniténtiæ, aliud, dignum pœniténtiæ fructum fácere. Ut enim secúndum dignos pœniténtiæ fructus loquámur, sciéndum est, quia quisquis illícita nulla commísit, huic jure concéditur, ut lícitis utátur: sicque pietátis ópera fáciat, ut tamen si volúerit, ea quæ mundi sunt, non relínquat.\par
\switchcolumn
\EN
\noindent We also have sinned, we have fallen into wicked habits. What must we do, if we would flee from the wrath to come? Let us hear John. Bring forth fruits worthy of repentance. In which words let us remark that the Friend of the Bridegroom demandeth not only fruits of repentance, but fruits worthy of repentance. The former are one thing, and the latter another. In considering then what are fruits worthy of repentance, we may remark that if we had done nothing unlawful we might have had free use of things which are lawful, and been able to sanctify ourselves without abstaining from indulgence in the things of the world.\par
\switchcolumn*

\LA
\begin{Resp}\Rsign Non discédimus a te, vivificábis nos, Dómine, et nomen tuum invocábimus: \Star{} Osténde nobis fáciem tuam, et salvi érimus.\end{Resp}
\begin{Resp}\Vsign Meménto nostri, Dómine, in beneplácito pópuli tui: vísita nos in salutári tuo.\end{Resp}
\begin{Resp}\Rsign Osténde nobis fáciem tuam, et salvi érimus.\end{Resp}
\switchcolumn
\EN
\begin{Resp}\Rsign We will not go back from thee. Thou, O Lord, shalt quicken us, and we will call upon thy name. \Star{} Cause thy face to shine upon us, and we shall be saved.\end{Resp}
\begin{Resp}\Vsign Remember us, O Lord, with the favour that Thou showest unto thy people; O visit us with thy salvation.\end{Resp}
\begin{Resp}\Rsign Cause thy face to shine upon us, and we shall be saved.\end{Resp}
\switchcolumn*

\LA
\LessonHead{Lectio IX}
\switchcolumn
\EN
\LessonHead{Lesson IX}
\switchcolumn*

\LA
\noindent At si quis in fornicatiónis culpam, vel fortasse, quod est grávius, in adultérium lapsus est: tanto a se lícita debet abscíndere, quanto se méminit et illícita perpetrásse. Neque enim par fructus boni óperis esse debet, ejus qui minus, et ejus qui ámplius delíquit: aut ejus qui in nullis, et ejus qui in quibúsdam facinóribus cécidit, et ejus qui in multis est lapsus. Per hoc ergo quod dícitur: Fácite fructus dignos pœniténtiæ: uniuscujúsque consciéntia convenítur, ut tanto majóra acquírat bonórum óperum lucra per pœniténtiam, quanto gravióra sibi íntulit damna per culpam.\par
\switchcolumn
\EN
\noindent But if any one, for example, hath fallen into fornication, or perhaps, into what is much worse, adultery, he ought to make up for his lawless pleasure by abstaining in some degree from lawful enjoyments. He that hath sinned less is not bound to mortify himself as much as he that hath sinned more, nor he that is innocent like him that is guilty. Let every one hearing these words bring forth fruits worthy of repentance, proceed to judge himself by his own conscience, and the more he perceiveth that he hath sinned, the greater penance let him do.\par
\switchcolumn*

\LA
\begin{Resp}\Rsign Intuémini, quantus sit iste, qui ingréditur ad salvándas gentes: ipse est Rex justítiæ, \Star{} Cujus generátio non habet finem.\end{Resp}
\begin{Resp}\Vsign Præcúrsor pro nobis ingréditur, secúndum órdinem Melchísedech Póntifex factus in ætérnum.\end{Resp}
\begin{Resp}\Rsign Cujus generátio non habet finem.\end{Resp}
\begin{Resp}\Vsign Glória Patri, et Fílio, \Star{} et Spirítui Sancto.\end{Resp}
\begin{Resp}\Rsign Cujus generátio non habet finem.\end{Resp}
\switchcolumn
\EN
\begin{Resp}\Rsign Consider how great this man is, who is entered in for the salvation of the nations; he is King of Righteousness; \Star{} Without descent, nor end of life.\end{Resp}
\begin{Resp}\Vsign The Fore-runner is for us entered, made an High Priest for ever after the order of Melchisedek.\end{Resp}
\begin{Resp}\Rsign Without descent, nor end of life.\end{Resp}
\begin{Resp}\Vsign Glory be to the Father, and to the Son, \Star{} and to the Holy Ghost.\end{Resp}
\begin{Resp}\Rsign Without descent, nor end of life.\end{Resp}
\switchcolumn*

\end{paracol}

\Office{Feria II infra Hebdomadam IV Adventus}{Monday of the Fourth Week of Advent}{Feria major}
\begin{paracol}{2}
\LA
\LessonHead{Lectio I}
\switchcolumn
\EN
\LessonHead{Lesson I}
\switchcolumn*

\LA
\LessonTitle{De Isaía Prophéta}
\switchcolumn
\EN
\LessonTitle{Lesson from the book of Isaias}
\switchcolumn*

\LA
\Cite{Isa 41:8-10}
\switchcolumn
\EN
\Cite{Isa 41:8-10}
\switchcolumn*

\LA
\noindent Et tu, Israël, serve meus, Jacob quem elégi, semen Abraham amíci mei: in quo apprehéndi te ab extrémis terræ, et a longinquis ejus vocávi te, et dixi tibi: Servus meus es tu, elégi te, et non abjéci te. Ne tímeas, quia ego tecum sum: ne declínes, quia ego Deus tuus: confortávi te, et auxiliátus sum tibi, et suscépit te déxtera justi mei.\par
\switchcolumn
\EN
\noindent But thou Israel, art my servant, Jacob whom I have chosen, the seed of Abraham my friend: In whom I have taken thee from the ends of the earth, and from the remote parts thereof have called thee, and said to thee: Thou art my servant, I have chosen thee, and have not cast thee away. Fear not, for I am with thee: turn not aside, for I am thy God: I have strengthened thee, and have helped thee, and the right hand of my just one hath upheld thee.\par
\switchcolumn*

\LA
\begin{Resp}\Rsign Cánite tuba in Sion, vocáte gentes, annuntiáte pópulis, et dícite: \Star{} Ecce Deus Salvátor noster advéniet.\end{Resp}
\begin{Resp}\Vsign Annuntiáte, et audítum fácite: loquímini, et clamáte.\end{Resp}
\begin{Resp}\Rsign Ecce Deus Salvátor noster advéniet.\end{Resp}
\switchcolumn
\EN
\begin{Resp}\Rsign Blow ye the trumpet in Zion, call together the nations, tell it out among the people, and say: \Star{} Behold, God our Saviour cometh.\end{Resp}
\begin{Resp}\Vsign Tell it out and make it to be heard; speak aloud and cry.\end{Resp}
\begin{Resp}\Rsign Behold, God our Saviour cometh.\end{Resp}
\switchcolumn*

\LA
\LessonHead{Lectio II}
\switchcolumn
\EN
\LessonHead{Lesson II}
\switchcolumn*

\LA
\Cite{Isa 41:11-13}
\switchcolumn
\EN
\Cite{Isa 41:11-13}
\switchcolumn*

\LA
\noindent Ecce confundéntur et erubéscent omnes, qui pugnant advérsum te: erunt quasi non sint, et períbunt viri, qui contradícunt tibi. Quæres eos, et non invénies, viros rebélles tuos: erunt quasi non sint, et véluti consúmptio hómines bellántes advérsum te. Quia ego Dóminus Deus tuus apprehéndens manum tuam, dicénsque tibi: Ne tímeas, ego adjúvi te.\par
\switchcolumn
\EN
\noindent Behold all that fight against thee shall be confounded and ashamed, they shall be as nothing, and the men shall perish that strive against thee. Thou shalt seek them, and shalt not find the men that resist thee: they shall be as nothing: and as a thing consumed the men that war against thee. For I am the Lord thy God, who take thee by the hand, and say to thee: Fear not, I have helped thee.\par
\switchcolumn*

\LA
\begin{Resp}\Rsign Non auferétur sceptrum de Juda, et dux de fémore ejus, donec véniat qui mitténdus est: \Star{} Et ipse erit exspectátio géntium.\end{Resp}
\begin{Resp}\Vsign Pulchrióres sunt óculi ejus vino, et dentes ejus lacte candidióres.\end{Resp}
\begin{Resp}\Rsign Et ipse erit exspectátio géntium.\end{Resp}
\switchcolumn
\EN
\begin{Resp}\Rsign The sceptre shall not depart from Judah, nor the law-giver from his loins, until he that shall be sent cometh. \Star{} And unto him shall the longing of the Gentiles be.\end{Resp}
\begin{Resp}\Vsign His eyes shall be bright with wine, and his teeth white with milk.\end{Resp}
\begin{Resp}\Rsign And unto him shall the desire of the Gentiles be.\end{Resp}
\switchcolumn*

\LA
\LessonHead{Lectio III}
\switchcolumn
\EN
\LessonHead{Lesson III}
\switchcolumn*

\LA
\Cite{Isa 41:14-16}
\switchcolumn
\EN
\Cite{Isa 41:14-16}
\switchcolumn*

\LA
\noindent Noli timére, vermis Jacob, qui mórtui estis ex Israël: ego auxiliátus sum tibi, dicit Dóminus: et redémptor tuus Sanctus Israël. Ego pósui te quasi plaustrum tritúrans novum, habens rostra serrántia: triturábis montes, et commínues: et colles quasi púlverem pones. Ventilábis eos, et ventus tollet, et turbo dispérget eos: et tu exsultábis in Dómino, in Sancto Israël lætáberis.\par
\switchcolumn
\EN
\noindent Fear not, thou worm of Jacob, you that are dead of Israel: I have helped thee, saith the Lord: and thy Redeemer the Holy One of Israel. I have made thee as a new thrashing wain, with teeth like a saw: thou shall thrash the mountains, and break them in pieces: and shalt make the hills as chaff. Thou shalt fan them, and the wind shall carry them away, and the whirlwind shall scatter them: and thou shalt rejoice in the Lord, in the Holy One of Israel thou shalt be joyful.\par
\switchcolumn*

\LA
\begin{Resp}\Rsign Me opórtet mínui, illum autem créscere: qui autem post me venit, ante me factus est: \Star{} Cujus non sum dignus corrígiam calceamentórum sólvere.\end{Resp}
\begin{Resp}\Vsign Ego baptizávi vos aqua: ille autem baptizábit vos Spíritu Sancto.\end{Resp}
\begin{Resp}\Rsign Cujus non sum dignus corrígiam calceamentórum sólvere.\end{Resp}
\begin{Resp}\Vsign Glória Patri, et Fílio, \Star{} et Spirítui Sancto.\end{Resp}
\begin{Resp}\Rsign Cujus non sum dignus corrígiam calceamentórum sólvere.\end{Resp}
\switchcolumn
\EN
\begin{Resp}\Rsign I must decrease, but He must increase; He it is Who, coming after me is preferred before me; \Star{} Whose shoe's latchet I am not worthy to unloose.\end{Resp}
\begin{Resp}\Vsign I baptize you with water; but He shall baptize you with the Holy Ghost.\end{Resp}
\begin{Resp}\Rsign Whose shoe's latchet I am not worthy to unloose.\end{Resp}
\begin{Resp}\Vsign Glory be to the Father, and to the Son, \Star{} and to the Holy Ghost.\end{Resp}
\begin{Resp}\Rsign Whose shoe's latchet I am not worthy to unloose.\end{Resp}
\switchcolumn*

\end{paracol}

\Office{Feria III infra Hebdomadam IV Adventus}{Tuesday of the Fourth Week of Advent}{Feria major}
\begin{paracol}{2}
\LA
\LessonHead{Lectio I}
\switchcolumn
\EN
\LessonHead{Lesson I}
\switchcolumn*

\LA
\LessonTitle{De Isaía Prophéta}
\switchcolumn
\EN
\LessonTitle{Lesson from the book of Isaias}
\switchcolumn*

\LA
\Cite{Isa 42:1-4}
\switchcolumn
\EN
\Cite{Isa 42:1-4}
\switchcolumn*

\LA
\noindent Ecce servus meus, suscípiam eum: eléctus meus, complácuit sibi in illo ánima mea: dedi spíritum meum super eum, judícium géntibus próferet. Non clamábit, neque accípiet persónam, nec audiétur vox ejus foris. Cálamum quassátum non cónteret, et linum fúmigans non exstínguet: in veritáte edúcet judícium. Non erit tristis, neque turbuléntus, donec ponat in terra judícium: et legem ejus ínsulæ exspectábunt.\par
\switchcolumn
\EN
\noindent Behold my servant, I will uphold him: my elect, my soul delighteth in him: I have given my spirit upon him, he shall bring forth judgment to the Gentiles. He shall not cry, nor have respect to person, neither shall his voice be heard abroad. The bruised reed he shall not break, and smoking flax he shall not quench: he shall bring forth judgment unto truth. He shall not be sad, nor troublesome, till he set judgment in the earth: and the islands shall wait for his law.\par
\switchcolumn*

\LA
\begin{Resp}\Rsign Nascétur nobis párvulus, et vocábitur Deus, Fortis: \Star{} Ipse sedébit super thronum David patris sui, et imperábit: cujus potéstas super húmerum ejus.\end{Resp}
\begin{Resp}\Vsign In ipso benedicéntur omnes tribus terræ, omnes gentes sérvient ei.\end{Resp}
\begin{Resp}\Rsign Ipse sedébit super thronum David patris sui, et imperábit: cujus potéstas super húmerum ejus.\end{Resp}
\switchcolumn
\EN
\begin{Resp}\Rsign Unto us shall a Child be born, and His name shall be called the Mighty God. \Star{} He shall sit upon the throne of His father David, and shall reign, and the government shall be upon His shoulder.\end{Resp}
\begin{Resp}\Vsign In Him shall all the kindreds of the earth be blessed; all nations shall serve Him.\end{Resp}
\begin{Resp}\Rsign He shall sit upon the throne of His father David, and shall reign, and the government shall be upon His shoulder.\end{Resp}
\switchcolumn*

\LA
\LessonHead{Lectio II}
\switchcolumn
\EN
\LessonHead{Lesson II}
\switchcolumn*

\LA
\Cite{Isa 42:5-7}
\switchcolumn
\EN
\Cite{Isa 42:5-7}
\switchcolumn*

\LA
\noindent Hæc dicit Dóminus Deus creans cælos, et exténdens eos: firmans terram, et quæ gérminant ex ea: dans flatum pópulo, qui est super eam, et spíritum calcántibus eam. Ego Dóminus vocávi te in justítia, et apprehéndi manum tuam, et servávi te. Et dedi te in fœdus pópuli, in lucem géntium: ut aperíres óculos cæcórum, et edúceres de conclusióne vinctum, de domo cárceris sedéntes in ténebris.\par
\switchcolumn
\EN
\noindent Thus saith the Lord God that created the heavens, and stretched them out: that established the earth, and the things that spring out of it: that giveth breath to the people upon it, and spirit to them that tread thereon. I the Lord have called thee in justice, and taken thee by the hand, and preserved thee. And I have given thee for a covenant of the people, for a light of the Gentiles: That thou mightest open the eyes of the blind, and bring forth the prisoner out of prison, and them that sit in darkness out of the prison house.\par
\switchcolumn*

\LA
\begin{Resp}\Rsign Ecce jam venit plenitúdo témporis, in quo misit Deus Fílium suum in terras, natum de Vírgine, factum sub lege: \Star{} Ut eos, qui sub lege erant, redímeret.\end{Resp}
\begin{Resp}\Vsign Propter nímiam caritátem suam, qua diléxit nos Deus, Fílium suum misit in similitúdinem carnis peccáti.\end{Resp}
\begin{Resp}\Rsign Ut eos, qui sub lege erant, redímeret.\end{Resp}
\switchcolumn
\EN
\begin{Resp}\Rsign Behold, the fullness of the time is come, wherein God hath sent forth His Son into the world, born of a Virgin, made under the law; \Star{} To redeem them that were under the law.\end{Resp}
\begin{Resp}\Vsign God, for His great love wherewith He loved us, hath sent His own Son in the likeness of sinful flesh.\end{Resp}
\begin{Resp}\Rsign To redeem them that were under the law.\end{Resp}
\switchcolumn*

\LA
\LessonHead{Lectio III}
\switchcolumn
\EN
\LessonHead{Lesson III}
\switchcolumn*

\LA
\Cite{Isa 42:10-13}
\switchcolumn
\EN
\Cite{Isa 42:10-13}
\switchcolumn*

\LA
\noindent Cantáte Dómino cánticum novum, laus ejus ab extrémis terræ: qui descénditis in mare, et plenitúdo ejus, ínsulæ, et habitatóres eárum. Sublevétur desértum, et civitátes ejus; in dómibus habitábit Cedar: laudáte habitatóres Petræ, de vértice móntium clamábunt. Ponent Dómino glóriam, et laudem ejus in ínsulis nuntiábunt. Dóminus sicut fortis egrediétur, sicut vir præliátor suscitábit zelum: vociferábitur, et clamábit: super inimícos suos confortábitur.\par
\switchcolumn
\EN
\noindent Sing ye to the Lord a new song, his praise is from the ends of the earth: you that go down to the sea, and all that are therein: ye islands, and ye inhabitants of them. Let the desert and the cities thereof be exalted: Cedar shall dwell in houses: ye inhabitants of Petra, give praise, they shall cry from the top of the mountains. They shall give glory to the Lord, and shall declare his praise in the islands. The Lord shall go forth as a mighty man, as a man of war shall he stir up zeal: he shall shout and cry: he shall prevail against his enemies.\par
\switchcolumn*

\LA
\begin{Resp}\Rsign Virgo Israël, revértere ad civitátes tuas: \Star{} Usquequo dolens avertéris? generábis Dóminum Salvatórem, oblatiónem novam in terra: * Ambulábunt hómines in salvatiónem.\end{Resp}
\begin{Resp}\Vsign In caritáte perpétua diléxi te: ídeo attráxi te míserans tui.\end{Resp}
\begin{Resp}\Rsign Usquequo dolens avertéris? generábis Dóminum Salvatórem, oblatiónem novam in terra.\end{Resp}
\begin{Resp}\Vsign Glória Patri, et Fílio, \Star{} et Spirítui Sancto.\end{Resp}
\begin{Resp}\Rsign Ambulábunt hómines in salvatiónem.\end{Resp}
\switchcolumn
\EN
\begin{Resp}\Rsign O virgin of Israel, turn again to thy cities. \Star{} How long wilt thou go about sorrowing? Thou shalt bring forth the Lord thy Saviour, a new offering in the earth; men shall walk in paths of salvation.\end{Resp}
\begin{Resp}\Vsign I have loved thee with an everlasting love; therefore with loving-kindness have I drawn thee.\end{Resp}
\begin{Resp}\Rsign How long wilt thou go about sorrowing? Thou shalt bring forth the Lord thy Saviour, a new offering in the earth.\end{Resp}
\begin{Resp}\Vsign Glory be to the Father, and to the Son, \Star{} and to the Holy Ghost.\end{Resp}
\begin{Resp}\Rsign Men shall walk in paths of salvation.\end{Resp}
\switchcolumn*

\end{paracol}

\Office{Feria IV infra Hebdomadam IV Adventus}{Wednesday of the Fourth Week of Advent}{Feria major}
\begin{paracol}{2}
\LA
\LessonHead{Lectio I}
\switchcolumn
\EN
\LessonHead{Lesson I}
\switchcolumn*

\LA
\LessonTitle{De Isaía Prophéta}
\switchcolumn
\EN
\LessonTitle{Lesson from the book of Isaias}
\switchcolumn*

\LA
\Cite{Isa 51:1-3}
\switchcolumn
\EN
\Cite{Isa 51:1-3}
\switchcolumn*

\LA
\noindent Audíte me, qui sequímini quod justum est, et quǽritis Dóminum: atténdite ad petram unde excísi estis, et ad cavérnam laci, de qua præcísi estis. Atténdite ad Abraham patrem vestrum, et ad Saram, quæ péperit vos, quia unum vocávi eum, et benedíxi ei, et multiplicávi eum. Consolábitur ergo Dóminus Sion, et consolábitur omnes ruínas ejus: et ponet desértum ejus quasi delícias, et solitúdinem ejus quasi hortum Dómini. Gáudium et lætítia inveniétur in ea, gratiárum áctio et vox laudis.\par
\switchcolumn
\EN
\noindent Give ear to me, you that follow that which is just, and you that seek the Lord: look unto the rock whence you are hewn, and to the hole of the pit from which you are dug out. Look unto Abraham your father, and to Sara that bore you: for I called him alone, and blessed him, and multiplied him. The Lord therefore will comfort Sion, and will comfort all the ruins thereof: and he will make her desert as a place of pleasure, and her wilderness as the garden of the Lord. Joy and gladness shall be found therein, thanksgiving, and the voice of praise.\par
\switchcolumn*

\LA
\begin{Resp}\Rsign Jurávi, dicit Dóminus, ut ultra jam non iráscar super terram: montes enim et colles suscípient justítiam meam, \Star{} Et testaméntum pacis erit in Jerúsalem.\end{Resp}
\begin{Resp}\Vsign Juxta est salus mea, ut véniat: et justítia mea, ut revelétur.\end{Resp}
\begin{Resp}\Rsign Et testaméntum pacis erit in Jerúsalem.\end{Resp}
\switchcolumn
\EN
\begin{Resp}\Rsign I have sworn, saith the Lord, that I will not be wroth any more with the earth; for the mountains and the hills shall receive My righteousness. \Star{} And the covenant of My peace shall be in Jerusalem.\end{Resp}
\begin{Resp}\Vsign My salvation is near to come, and My righteousness to be revealed.\end{Resp}
\begin{Resp}\Rsign And the covenant of My peace shall be in Jerusalem.\end{Resp}
\switchcolumn*

\LA
\LessonHead{Lectio II}
\switchcolumn
\EN
\LessonHead{Lesson II}
\switchcolumn*

\LA
\Cite{Isa 51:4-6}
\switchcolumn
\EN
\Cite{Isa 51:4-6}
\switchcolumn*

\LA
\noindent Atténdite ad me, pópule meus, et tribus mea, me audíte: quia lex a me éxiet, et judícium meum in lucem populórum requiéscet. Prope est justus meus, egréssus est salvátor meus, et brácchia mea pópulos judicábunt: me ínsulæ exspectábunt, et brácchium meum sustinébunt. Leváte in cælum óculos vestros, et vidéte sub terra deórsum: quia cæli sicut fumus liquéscent, et terra sicut vestiméntum atterétur, et habitatóres ejus sicut hæc interíbunt: Salus autem mea in sempitérnum erit, et justítia mea non defíciet.\par
\switchcolumn
\EN
\noindent Hearken unto me, O my people, and give ear to me, O my tribes: for a law shall go forth from me, and my judgment shall rest to be a light of the nations. My just one is near at hand, my saviour is gone forth, and my arms shall judge the people: the islands shall look for me, and shall patiently wait for my arm. Lift up your eyes to heaven, and look down to the earth beneath: for the heavens shall vanish like smoke, and the earth shall be worn away like a garment, and the inhabitants thereof shall perish in like manner: but my salvation shall be for ever, and my justice shall not fail.\par
\switchcolumn*

\LA
\begin{Resp}\Rsign Non discédimus a te, vivificábis nos, Dómine, et nomen tuum invocábimus: \Star{} Osténde nobis fáciem tuam, et salvi érimus.\end{Resp}
\begin{Resp}\Vsign Meménto nostri, Dómine, in beneplácito pópuli tui: vísita nos in salutári tuo.\end{Resp}
\begin{Resp}\Rsign Osténde nobis fáciem tuam, et salvi érimus.\end{Resp}
\switchcolumn
\EN
\begin{Resp}\Rsign We will not go back from thee. Thou, O Lord, shalt quicken us, and we will call upon thy name. \Star{} Cause thy face to shine upon us, and we shall be saved.\end{Resp}
\begin{Resp}\Vsign Remember us, O Lord, with the favour that Thou showest unto thy people; O visit us with thy salvation.\end{Resp}
\begin{Resp}\Rsign Cause thy face to shine upon us, and we shall be saved.\end{Resp}
\switchcolumn*

\LA
\LessonHead{Lectio III}
\switchcolumn
\EN
\LessonHead{Lesson III}
\switchcolumn*

\LA
\Cite{Isa 51:7-8}
\switchcolumn
\EN
\Cite{Isa 51:7-8}
\switchcolumn*

\LA
\noindent Audíte me, qui scitis justum, pópulus meus, lex mea in corde eórum: nolíte timére oppróbrium hóminum, et blasphémias eórum ne metuátis. Sicut enim vestiméntum, sic cómedet eos vermis: et sicut lanam, sic devorábit eos tínea: Salus autem mea in sempitérnum erit, et justítia mea in generatiónes generatiónum.\par
\switchcolumn
\EN
\noindent Hearken to me, you that know what is just, my people who have my law in your heart: fear ye not the reproach of men, and be not afraid of their blasphemies. For the worm shall eat them up as a garment: and the moth shall consume them as wool: but my salvation shall be for ever, and my justice from generation to generation,\par
\switchcolumn*

\LA
\begin{Resp}\Rsign Intuémini, quantus sit iste, qui ingréditur ad salvándas gentes: ipse est Rex justítiæ, \Star{} Cujus generátio non habet finem.\end{Resp}
\begin{Resp}\Vsign Præcúrsor pro nobis ingréditur, secúndum órdinem Melchísedech Póntifex factus in ætérnum.\end{Resp}
\begin{Resp}\Rsign Cujus generátio non habet finem.\end{Resp}
\begin{Resp}\Vsign Glória Patri, et Fílio, \Star{} et Spirítui Sancto.\end{Resp}
\begin{Resp}\Rsign Cujus generátio non habet finem.\end{Resp}
\switchcolumn
\EN
\begin{Resp}\Rsign Consider how great this man is, who is entered in for the salvation of the nations; he is King of Righteousness; \Star{} Without descent, nor end of life.\end{Resp}
\begin{Resp}\Vsign The Fore-runner is for us entered, made an High Priest for ever after the order of Melchisedek.\end{Resp}
\begin{Resp}\Rsign Without descent, nor end of life.\end{Resp}
\begin{Resp}\Vsign Glory be to the Father, and to the Son, \Star{} and to the Holy Ghost.\end{Resp}
\begin{Resp}\Rsign Without descent, nor end of life.\end{Resp}
\switchcolumn*

\end{paracol}

\Office{Feria V infra Hebdomadam IV Adventus}{Thursday of the Fourth Week of Advent}{Feria major}
\begin{paracol}{2}
\LA
\LessonHead{Lectio I}
\switchcolumn
\EN
\LessonHead{Lesson I}
\switchcolumn*

\LA
\LessonTitle{De Isaía Prophéta}
\switchcolumn
\EN
\LessonTitle{Lesson from the book of Isaias}
\switchcolumn*

\LA
\Cite{Isa 64:1-4}
\switchcolumn
\EN
\Cite{Isa 64:1-4}
\switchcolumn*

\LA
\noindent Utinam dirúmperes cælos, et descénderes: a fácie tua montes deflúerent. Sicut exústio ignis tabéscerent, aquæ ardérent igni, ut notum fíeret nomen tuum inimícis tuis: a fácie tua gentes turbaréntur. Cum féceris mirabília, non sustinébimus: descendísti, et a fácie tua montes defluxérunt. A sǽculo non audiérunt, neque áuribus percepérunt: óculus non vidit, Deus absque te, quæ præparásti exspectántibus te.\par
\switchcolumn
\EN
\noindent That thou wouldst rend the heavens, and wouldst come down: the mountains would melt away at thy presence. They would melt as at the burning of fire, the waters would burn with fire, that thy name might be made known to thy enemies: that the nations might tremble at thy presence. When thou shalt do wonderful things, we shall not bear them: thou didst come down, and at thy presence the mountains melted away. From the beginning of the world they have not heard, nor perceived with the ears: the eye hath not seen, O God, besides thee, what things thou hast prepared for them that wait for thee.\par
\switchcolumn*

\LA
\begin{Resp}\Rsign Cánite tuba in Sion, vocáte gentes, annuntiáte pópulis, et dícite: \Star{} Ecce Deus Salvátor noster advéniet.\end{Resp}
\begin{Resp}\Vsign Annuntiáte, et audítum fácite: loquímini, et clamáte.\end{Resp}
\begin{Resp}\Rsign Ecce Deus Salvátor noster advéniet.\end{Resp}
\switchcolumn
\EN
\begin{Resp}\Rsign Blow ye the trumpet in Zion, call together the nations, tell it out among the people, and say: \Star{} Behold, God our Saviour cometh.\end{Resp}
\begin{Resp}\Vsign Tell it out and make it to be heard; speak aloud and cry.\end{Resp}
\begin{Resp}\Rsign Behold, God our Saviour cometh.\end{Resp}
\switchcolumn*

\LA
\LessonHead{Lectio II}
\switchcolumn
\EN
\LessonHead{Lesson II}
\switchcolumn*

\LA
\Cite{Isa 64:5-7}
\switchcolumn
\EN
\Cite{Isa 64:5-7}
\switchcolumn*

\LA
\noindent Occurrísti lætánti, et faciénti justítiam: in viis tuis recordabúntur tui: ecce tu irátus es, et peccávimus: in ipsis fúimus semper, et salvábimur. Et facti sumus ut immúndus omnes nos, et quasi pannus menstruátæ univérsæ justítiæ nostræ: et cecídimus quasi fólium univérsi, et iniquitátes nostræ quasi ventus abstulérunt nos. Non est qui ínvocet nomen tuum: qui consúrgat, et téneat te: abscondísti fáciem tuam a nobis, et allisísti nos in manu iniquitátis nostræ.\par
\switchcolumn
\EN
\noindent Thou hast met him that rejoiceth, and doth justice: in thy ways they shall remember thee: behold thou art angry, and we have sinned: in them we have been always, and we shall be saved. And we are all become as one unclean, and all our justices as the rag of a menstruous woman: and we have all fallen as a leaf, and our iniquities, like the wind, have taken us away. There is none that calleth upon thy name: that riseth up, and taketh hold of thee: thou hast hid thy face from us, and hast crushed us in the hand of our iniquity.\par
\switchcolumn*

\LA
\begin{Resp}\Rsign Non auferétur sceptrum de Juda, et dux de fémore ejus, donec véniat qui mitténdus est: \Star{} Et ipse erit exspectátio géntium.\end{Resp}
\begin{Resp}\Vsign Pulchrióres sunt óculi ejus vino, et dentes ejus lacte candidióres.\end{Resp}
\begin{Resp}\Rsign Et ipse erit exspectátio géntium.\end{Resp}
\switchcolumn
\EN
\begin{Resp}\Rsign The sceptre shall not depart from Judah, nor the law-giver from his loins, until he that shall be sent cometh. \Star{} And unto him shall the longing of the Gentiles be.\end{Resp}
\begin{Resp}\Vsign His eyes shall be bright with wine, and his teeth white with milk.\end{Resp}
\begin{Resp}\Rsign And unto him shall the desire of the Gentiles be.\end{Resp}
\switchcolumn*

\LA
\LessonHead{Lectio III}
\switchcolumn
\EN
\LessonHead{Lesson III}
\switchcolumn*

\LA
\Cite{Isa 64:8-11}
\switchcolumn
\EN
\Cite{Isa 64:8-11}
\switchcolumn*

\LA
\noindent Et nunc, Dómine, pater noster es tu, nos vero lutum: et fictor noster tu, et ópera mánuum tuárum omnes nos. Ne irascáris, Dómine, satis, et ne ultra memíneris iniquitátis nostræ: ecce réspice, pópulus tuus omnes nos. Cívitas sancti tui facta est desérta, Sion desérta facta est, Jerúsalem desoláta est. Domus sanctificatiónis nostræ, et glóriæ nostræ, ubi laudavérunt te patres nostri, facta est in exustiónem ignis, et ómnia desiderabília nostra versa sunt in ruínas.\par
\switchcolumn
\EN
\noindent And now, O Lord, thou art our father, and we are clay: and thou art our maker, and we all are the works of thy hands. Be not very angry, O Lord, and remember no longer our iniquity: behold, see we are all thy people. The city of thy sanctuary is become a desert, Sion is made a desert, Jerusalem is desolate. The house of our holiness, and of our glory, where our fathers praised thee, is burnt with fire, and all our lovely things are turned into ruins.\par
\switchcolumn*

\LA
\begin{Resp}\Rsign Me opórtet mínui, illum autem créscere: qui autem post me venit, ante me factus est: \Star{} Cujus non sum dignus corrígiam calceamentórum sólvere.\end{Resp}
\begin{Resp}\Vsign Ego baptizávi vos aqua: ille autem baptizábit vos Spíritu Sancto.\end{Resp}
\begin{Resp}\Rsign Cujus non sum dignus corrígiam calceamentórum sólvere.\end{Resp}
\begin{Resp}\Vsign Glória Patri, et Fílio, \Star{} et Spirítui Sancto.\end{Resp}
\begin{Resp}\Rsign Cujus non sum dignus corrígiam calceamentórum sólvere.\end{Resp}
\switchcolumn
\EN
\begin{Resp}\Rsign I must decrease, but He must increase; He it is Who, coming after me is preferred before me; \Star{} Whose shoe's latchet I am not worthy to unloose.\end{Resp}
\begin{Resp}\Vsign I baptize you with water; but He shall baptize you with the Holy Ghost.\end{Resp}
\begin{Resp}\Rsign Whose shoe's latchet I am not worthy to unloose.\end{Resp}
\begin{Resp}\Vsign Glory be to the Father, and to the Son, \Star{} and to the Holy Ghost.\end{Resp}
\begin{Resp}\Rsign Whose shoe's latchet I am not worthy to unloose.\end{Resp}
\switchcolumn*

\end{paracol}

\Office{Feria VI infra Hebdomadam IV Adventus}{Friday of the Fourth Week of Advent}{Feria major}
\begin{paracol}{2}
\LA
\LessonHead{Lectio I}
\switchcolumn
\EN
\LessonHead{Lesson I}
\switchcolumn*

\LA
\LessonTitle{De Isaía Prophéta}
\switchcolumn
\EN
\LessonTitle{Lesson from the book of Isaias}
\switchcolumn*

\LA
\Cite{Isa 66:5-8}
\switchcolumn
\EN
\Cite{Isa 66:5-8}
\switchcolumn*

\LA
\noindent Audíte verbum Dómini, qui trémitis ad verbum ejus: dixérunt fratres vestri odiéntes vos, et abiciéntes propter nomen meum: glorificétur Dóminus, et vidébimus in lætítia vestra: ipsi autem confundéntur. Vox pópuli de civitáte, vox de templo, vox Dómini reddéntis retributiónem inimícis suis. Antequam parturíret, péperit: ántequam veníret partus ejus, péperit másculum. Quis audívit umquam tale? et quis vidit huic símile? Numquid partúriet terra in die una? aut pariétur gens simul, quia parturívit et péperit Sion fílios suos?\par
\switchcolumn
\EN
\noindent Hear the word of the Lord, you that tremble at his word: Your brethren that hate you, and cast you out for my name's sake, have said: Let the Lord be glorified, and we shall see in your joy: but they shall be confounded. A voice of the people from the city, a voice from the temple, the voice of the Lord that rendereth recompense to his enemies. Before she was in labour, she brought forth; before her time came to be delivered, she brought forth a man child. Who hath ever heard such a thing? and who hath seen the like to this? shall the earth bring forth in one day? or shall a nation be brought forth at once, because Sion hath been in labour, and hath brought forth her children?\par
\switchcolumn*

\LA
\begin{Resp}\Rsign Nascétur nobis párvulus, et vocábitur Deus, Fortis: \Star{} Ipse sedébit super thronum David patris sui, et imperábit: cujus potéstas super húmerum ejus.\end{Resp}
\begin{Resp}\Vsign In ipso benedicéntur omnes tribus terræ, omnes gentes sérvient ei.\end{Resp}
\begin{Resp}\Rsign Ipse sedébit super thronum David patris sui, et imperábit: cujus potéstas super húmerum ejus.\end{Resp}
\switchcolumn
\EN
\begin{Resp}\Rsign Unto us shall a Child be born, and His name shall be called the Mighty God. \Star{} He shall sit upon the throne of His father David, and shall reign, and the government shall be upon His shoulder.\end{Resp}
\begin{Resp}\Vsign In Him shall all the kindreds of the earth be blessed; all nations shall serve Him.\end{Resp}
\begin{Resp}\Rsign He shall sit upon the throne of His father David, and shall reign, and the government shall be upon His shoulder.\end{Resp}
\switchcolumn*

\LA
\LessonHead{Lectio II}
\switchcolumn
\EN
\LessonHead{Lesson II}
\switchcolumn*

\LA
\Cite{Isa 66:9-12}
\switchcolumn
\EN
\Cite{Isa 66:9-12}
\switchcolumn*

\LA
\noindent Numquid ego, qui álios párere fácio, ipse non páriam, dicit Dóminus? si ego, qui generatiónem céteris tríbuo, stérilis ero, ait Dóminus, Deus tuus? Lætámini cum Jerúsalem, et exsultáte in ea, omnes qui dilígitis eam: gaudéte cum ea gáudio, univérsi qui lugétis super eam, ut sugátis, et repleámini ab úbere consolatiónis ejus: ut mulgeátis, et delíciis affluátis ab omnímoda glória ejus. Quia hæc dicit Dóminus: Ecce ego declinábo super eam quasi flúvium pacis, et quasi torréntem inundántem glóriam géntium, quam sugétis: ad úbera portabímini, et super génua blandiéntur vobis.\par
\switchcolumn
\EN
\noindent Shall not I that make others to bring forth children, myself bring forth, saith the Lord? shall I, that give generation to others, be barren, saith the Lord thy God? Rejoice with Jerusalem, and be glad with her, all you that love her: rejoice for joy with her, all you that mourn for her. That you may suck, and be filled with the breasts of her consolations: that you may milk out, and flow with delights, from the abundance of her glory. For thus saith the Lord: Behold I will bring upon her as it were a river of peace, and as an overflowing torrent the glory of the Gentiles, which you shall suck; you shall be carried at the breasts, end upon the knees they shall caress you.\par
\switchcolumn*

\LA
\begin{Resp}\Rsign Ecce jam venit plenitúdo témporis, in quo misit Deus Fílium suum in terras, natum de Vírgine, factum sub lege: \Star{} Ut eos, qui sub lege erant, redímeret.\end{Resp}
\begin{Resp}\Vsign Propter nímiam caritátem suam, qua diléxit nos Deus, Fílium suum misit in similitúdinem carnis peccáti.\end{Resp}
\begin{Resp}\Rsign Ut eos, qui sub lege erant, redímeret.\end{Resp}
\switchcolumn
\EN
\begin{Resp}\Rsign Behold, the fullness of the time is come, wherein God hath sent forth His Son into the world, born of a Virgin, made under the law; \Star{} To redeem them that were under the law.\end{Resp}
\begin{Resp}\Vsign God, for His great love wherewith He loved us, hath sent His own Son in the likeness of sinful flesh.\end{Resp}
\begin{Resp}\Rsign To redeem them that were under the law.\end{Resp}
\switchcolumn*

\LA
\LessonHead{Lectio III}
\switchcolumn
\EN
\LessonHead{Lesson III}
\switchcolumn*

\LA
\Cite{Isa 66:13-16}
\switchcolumn
\EN
\Cite{Isa 66:13-18}
\switchcolumn*

\LA
\noindent Quómodo si cui mater blandiátur, ita ego consolábor vos, et in Jerúsalem consolabímini. Vidébitis, et gaudébit cor vestrum, et ossa vestra quasi herba germinábunt, et cognoscétur manus Dómini servis ejus, et indignábitur inimícis suis. Quia ecce Dóminus in igne véniet, et quasi turbo quadrígæ ejus: réddere in indignatióne furórem suum, et increpatiónem suam in flamma ignis, quia in igne Dóminus dijudicábit, et in gládio suo ad omnem carnem, et multiplicabúntur interfécti a Dómino.\par
\switchcolumn
\EN
\noindent As one whom the mother caresseth, so will I comfort you, and you shall be comforted in Jerusalem. You shall see and your heart shall rejoice, and your bones shall flourish like an herb, and the hand of the Lord shall be known to his servants, and he shall be angry with his enemies. For behold the Lord will come with fire, and his chariots are like a whirlwind, to render his wrath in indignation, and his rebuke with flames of fire. For the Lord shall judge by fire, and by his sword unto all flesh, and the slain of the Lord shall be many. They that were sanctified, and thought themselves clean in the gardens behind the gate within, they that did eat swine's flesh, and the abomination, and the mouse: they shall be consumed together, saith the Lord. But I know their works, and their thoughts: I come that I may gather them together with all nations and tongues: and they shall come and shall see my glory.\par
\switchcolumn*

\LA
\begin{Resp}\Rsign Virgo Israël, revértere ad civitátes tuas: \Star{} Usquequo dolens avertéris? generábis Dóminum Salvatórem, oblatiónem novam in terra: * Ambulábunt hómines in salvatiónem.\end{Resp}
\begin{Resp}\Vsign In caritáte perpétua diléxi te: ídeo attráxi te míserans tui.\end{Resp}
\begin{Resp}\Rsign Usquequo dolens avertéris? generábis Dóminum Salvatórem, oblatiónem novam in terra.\end{Resp}
\begin{Resp}\Vsign Glória Patri, et Fílio, \Star{} et Spirítui Sancto.\end{Resp}
\begin{Resp}\Rsign Ambulábunt hómines in salvatiónem.\end{Resp}
\switchcolumn
\EN
\begin{Resp}\Rsign O virgin of Israel, turn again to thy cities. \Star{} How long wilt thou go about sorrowing? Thou shalt bring forth the Lord thy Saviour, a new offering in the earth; men shall walk in paths of salvation.\end{Resp}
\begin{Resp}\Vsign I have loved thee with an everlasting love; therefore with loving-kindness have I drawn thee.\end{Resp}
\begin{Resp}\Rsign How long wilt thou go about sorrowing? Thou shalt bring forth the Lord thy Saviour, a new offering in the earth.\end{Resp}
\begin{Resp}\Vsign Glory be to the Father, and to the Son, \Star{} and to the Holy Ghost.\end{Resp}
\begin{Resp}\Rsign Men shall walk in paths of salvation.\end{Resp}
\switchcolumn*

\end{paracol}
\end{document}
